\chapter{Hybrid spaces and degeneration}
\index{hybrid space}\index{degeneration!hybrid convergence}
\label{chap:hybrid-spaces}

Hybrid geometry places a punctured complex family and its
non-Archimedean limit in one topological space. It is tempting to read
this slogan as an automatic convergence theorem. That is false:
topological convergence of points, convergence of metrics, convergence
of Monge--Amp\`ere measures, and convergence of entropy are four
different assertions. This chapter records the construction and then
states the published convergence theorems with their actual hypotheses.

Throughout, $\Delta$ is a complex disc with coordinate $t$,
$\Delta^*=\Delta\setminus\{0\}$, and
\[
                         K=\mathbf C((t)).
\tag{66.1}\label{eq:hybrid-base-field}
\]
We fix $0<r<1$ and give $K$ the $t$-adic absolute value
$|t|_K=r$. We work with projective meromorphic degenerations of smooth complex
manifolds: after shrinking the disc, the punctured family has a flat
projective model over the disc and over the hybrid ring below. Normal
models and SNC modifications are used when their divisors or skeleta enter
the argument. The geometric measure theorems require these hypotheses;
the later Banach-family proof states its separate algebraic domain.

The book's additive valuation is $v_x=-\log|\cdot|_x$, so
$v_K(t)=-\log r$.  For integral-affine coordinates on skeleta we also
use the normalized valuation
\[
                    \widehat v_x:=\frac{v_x}{-\log r},
                    \qquad \widehat v_K(t)=1.
\]
The dependencies are explicit. \Cref{chap:skeleta} supplies monomial
retractions and weight-minimum skeleta; \Cref{chap:compare-ma} supplies the
fieldwise Monge--Amp\`ere dictionary; and
\Cref{chap:rays-filtrations} supplies the filtration ray and its energy
identity. The canonical psh extension is proved before its continuity is
discussed, and continuity is established before a measure limit is taken.

Thus formulas involving $b_E^{-1}\ord_E$ or
$\sum b_iw_i=1$ use $\widehat v$; metric values written with the
book's $-\log$-weight retain the factor $-\log r$.

\section{The hybrid base}
\label{sec:hybrid-base}

\begin{definition}[Hybrid norm]
\label{def:hybrid-norm}
On $\mathbf C$, put
\[
             \lVert a\rVert_{\mathrm{hyb}}
             :=\max\{|a|_\infty,|a|_0\},
\tag{66.2}\label{eq:hybrid-norm}
\]
where $|a|_0=1$ for $a\neq0$ and $|0|_0=0$. The Berkovich spectrum of this Banach ring is denoted $\mathcal M(\mathbf C_{\mathrm{hyb}})$.
\end{definition}

For $0<\lambda\leq1$, the seminorm
\[
                            a\longmapsto |a|_\infty^\lambda
\tag{66.3}\label{eq:hybrid-seminorm}
\]
is bounded by $\lVert\cdot\rVert_{\mathrm{hyb}}$. At
$\lambda=0$ it becomes the trivial absolute value. Thus
\[
       [0,1]\longrightarrow\mathcal M(\mathbf C_{\mathrm{hyb}}),
       \qquad
       \lambda\longmapsto|\cdot|_\infty^\lambda
\tag{66.4}\label{eq:hybrid-spectrum-parametrization}
\]
is a homeomorphism. The endpoint $0$ is not an extra complex
point: it is the trivially valued fibre.

To obtain a $t$-adic central fibre, use the Banach ring of convergent
\emph{Laurent} series
\[
 A_r=\left\{\sum_{k\in\mathbf Z}a_kt^k:
       \|f\|_r:=\sum_{k\in\mathbf Z}\|a_k\|_{\mathrm{hyb}}r^k<\infty\right\}.
\tag{66.4a}\label{eq:hybrid-laurent-ring}
\]
Only finitely many negative coefficients can be nonzero, since
$\|a_k\|_{\mathrm{hyb}}\geq1$ for a nonzero coefficient. Thus
$A_r\subset\mathbf C((t))$, and $t$ is invertible in $A_r$.
By \cite[Lemma~1.20]{PS23Hybrid}, its spectrum is homeomorphic to
$\overline\Delta_r=\{\tau\in\mathbf C:|\tau|\leq r\}$, with evaluations
\[
 |f|_\tau=
 \begin{cases}
 |f(\tau)|_\infty^{\log r/\log|\tau|},&\tau\ne0,\\
 r^{\ord_0 f},&\tau=0.
 \end{cases}
\tag{66.4b}\label{eq:hybrid-base-evaluations}
\]
Here $|0|_\tau=0$. These evaluations are bounded: the powered triangle
inequality and $|\tau|^{\lambda(\tau)}=r$ bound a Laurent sum by
$\|f\|_r$. Conversely, the bounds on $t$ and $t^{-1}$ force $|t|_x=r$.
The classification of bounded scalar seminorms then gives the displayed
complex evaluations or the unique $t$-adic evaluation, as in the cited
lemma. Laurent polynomials are dense, so this determines the evaluation
on all of $A_r$. The completed residue field at the centre is $K$.
Every point has $|t|_x=r$; its disc coordinate is the complex parameter
$\tau$, rather than the seminorm value of $t$.

\begin{remark}[Normalization]
\label{rem:hybrid-normalization}
On the fibre over $t\neq0$, the scalar absolute value is a positive power of $|\cdot|_\infty$. One commonly chooses the parameter so that
\[
                 \lambda(t)=\frac{\log r}{\log|t|}.
\tag{66.5}\label{eq:hybrid-lambda}
\]
Changing $r$ reparametrizes the same limiting geometry. Formulas for weights must use the same choice of $|t|_K$; otherwise an unnoticed factor $\log r$ appears.
\end{remark}

\begin{definition}[Hybrid analytification]
\label{def:hybrid-analytification}
Let $\pi:\mathcal X\to\Delta$ be projective and flat, with generic fibre $X_K$. Its hybrid analytification is the compact Hausdorff space
\[
                         \mathcal X^{\mathrm{hyb}}
\tag{66.6}\label{eq:hybrid-analytification}
\]
obtained by analytifying $\mathcal X$ over the hybrid Banach base. It comes with a continuous map
\[
                \pi^{\mathrm{hyb}}:\mathcal X^{\mathrm{hyb}}
                \longrightarrow\overline\Delta_r
\tag{66.7}\label{eq:hybrid-projection}
\]
such that
\[
 \begin{aligned}
   (\pi^{\mathrm{hyb}})^{-1}(t)&\simeq X_t^{\mathrm{an}},
                      &&t\neq0,\\
   (\pi^{\mathrm{hyb}})^{-1}(0)&\simeq X_K^{\mathrm{an}}.
 \end{aligned}
\tag{66.8}\label{eq:hybrid-fibres}
\]
\end{definition}

Here $X_t^{\mathrm{an}}$ means the usual complex analytic space.
The second line means Berkovich analytification over $K$. In
particular, it is not the complex analytic space associated to the
scheme-theoretic special fibre $\mathcal X_0$.  This distinction is
summarized along a radial path in the hybrid base in
\cref{fig:hybrid-fibres}.

\begin{figure}[htbp]
\centering
\begin{tikzpicture}[x=1cm,y=1cm,>=Stealth]
  \draw[thick] (0.65,0)--(10.65,0);
  \node[below,font=\small] at (5.65,-0.15) {$\overline\Delta_r$ (radial path)};
  \fill (1.15,0) circle (1.6pt);
  \fill (4.25,0) circle (1.4pt);
  \fill (7.25,0) circle (1.4pt);
  \fill (9.75,0) circle (1.4pt);
  \node[below,font=\small] at (1.15,-0.02) {$0$};
  \node[below,font=\small] at (4.25,-0.02) {$t_3$};
  \node[below,font=\small] at (7.25,-0.02) {$t_2$};
  \node[below,font=\small] at (9.75,-0.02) {$t_1$};
  \draw[-{Stealth[length=2mm]},thin] (9.25,-0.75)--(2.0,-0.75)
    node[midway,below,font=\scriptsize] {$t_i\longrightarrow0$};

  % Non-Archimedean central fibre.
  \coordinate (na) at (1.15,2.25);
  \draw[very thick] (na)--(0.35,3.05);
  \draw[very thick] (na)--(1.15,3.35);
  \draw[very thick] (na)--(2.0,3.0);
  \draw[very thick] (na)--(0.55,1.58);
  \fill (na) circle (1.5pt);
  \node[font=\small,align=center] at (1.15,3.72) {$X_K^{\mathrm{an}}$\\Berkovich fibre};

  % Archimedean fibres over nonzero points.
  \foreach \xx/\tt in {4.25/t_3,7.25/t_2,9.75/t_1}{
    \draw[very thick] (\xx,2.42) ellipse (0.92 and 0.48);
    \draw[thin] (\xx-0.72,2.42) .. controls (\xx-0.25,2.82) and
      (\xx+0.25,2.02) .. (\xx+0.72,2.42);
    \node[font=\small] at (\xx,3.18) {$X_{\tt}^{\mathrm{an}}$};
  }

  \draw[-{Stealth[length=2mm]},thin] (1.15,1.5)--(1.15,0.18);
  \draw[-{Stealth[length=2mm]},thin] (4.25,1.92)--(4.25,0.18);
  \draw[-{Stealth[length=2mm]},thin] (7.25,1.92)--(7.25,0.18);
  \draw[-{Stealth[length=2mm]},thin] (9.75,1.92)--(9.75,0.18);
  \node[font=\small,fill=white,inner sep=1pt] at (5.65,1.35)
    {$\pi^{\mathrm{hyb}}$};
  \node[font=\small] at (5.65,4.0) {$\mathcal X^{\mathrm{hyb}}$};
\end{tikzpicture}
\caption{Fibres of the hybrid projection along a radial path.  The nonzero
fibres are complex analytic spaces.  The fibre over the centre is the
Berkovich space $X_K^{\mathrm{an}}$ over $\mathbf C((t))$, not the complex
analytic space of the scheme-theoretic fibre $\mathcal X_0$.  The shapes of
the fibres are schematic.}
\label{fig:hybrid-fibres}
\end{figure}

\begin{example}[The affine line]
\label{ex:hybrid-affine-line}
On $\mathbf A^1_{A_r}$, evaluations of $f\in A_r[z]$ define the topology. A net $(z_\tau,\tau)$ approaching the central point $x$ therefore satisfies
\[
 \lambda(\tau)\log|f(\tau,z_\tau)|_\infty
       \longrightarrow\log|f(x)|_K
       =-v_x(f)=\widehat v_x(f)\log r.
\tag{66.9}\label{eq:hybrid-polynomial-test}
\]
The extended value $-\infty$ is allowed. Take $z_\tau=\tau^{-q}$ for an integer $q\geq0$. Substitution gives a convergent meromorphic germ $g(t)=f(t,t^{-q})$. If $g\ne0$, write $g=t^a u$ with $u(0)\ne0$; then $\lambda(\tau)\log|g(\tau)|\to a\log r$. If $g=0$, both sides are $-\infty$. Hence the net converges to the $K$-rational point $z=t^{-q}$, even though $|z_\tau|\to\infty$ when $q>0$. Arbitrary formal coefficients in $K[z]$ cannot be evaluated at complex $\tau$; the completion $K$ belongs to the central fibre, and the topology uses the convergent chart ring.
\end{example}

\section{Chartwise convergence and normalized valuations}
\label{sec:hybrid-chartwise-convergence}

The topology becomes concrete on a normal-crossings chart.  Start with the
nodal model
\[
                         xy=t.
\tag{66.10}\label{eq:hybrid-node-chart}
\]
Choose complex points with
$|x_t|=|t|^\alpha$ and
$|y_t|=|t|^{1-\alpha}$, where $0\leq\alpha\leq1$.  Since
$\lambda(t)=\log r/\log|t|$,
\[
 \lambda(t)\log|x_t|=\alpha\log r,
 \qquad
 \lambda(t)\log|y_t|=(1-\alpha)\log r.
\tag{66.11}\label{eq:hybrid-node-evaluation}
\]
These rates determine a point of the monomial retraction. To conclude
that the analytic limit itself is monomial, one must also check that every
regular-function evaluation has its monomial value, including the absence
of leading cancellation. The monomial point has
\[
 \widehat v_\alpha(x)=\alpha,
 \qquad
 \widehat v_\alpha(y)=1-\alpha,
 \qquad
 \widehat v_\alpha(t)=1.
\tag{66.12}\label{eq:hybrid-node-monomial-valuation}
\]
As $\alpha$ runs from $0$ to $1$, these points trace the edge of the
dual graph.  The interval records logarithmic rates; those rates alone do not specify
all analytic points retracting to it.

For a general SNC chart
\[
                    t=u z_0^{b_0}\cdots z_p^{b_p},
                  \qquad |u|\asymp1,
\tag{66.13}\label{eq:hybrid-general-snc-chart}
\]
put $w_i=\log|z_i|/\log|t|$.  Taking logarithms and dividing by
$\log|t|$ yields
\[
                     \sum_{i=0}^pb_iw_i=1+o(1).
\tag{66.14}\label{eq:hybrid-chart-barycentric-limit}
\]
This is the barycentric equation
\eqref{eq:hybrid-simplex}.  The factor $b_i$ belongs in the affine
constraint, while $b_i^{-1}$ belongs in the normalized divisorial
valuation at the $i$-th vertex.  They cannot both be suppressed by saying
that the model is ``essentially reduced.''

\subsection{A polynomial convergence test on the node}
\label{sec:hybrid-node-polynomial-test}

Let
$f=\sum_{a,b}c_{ab}(t)x^ay^b$ have convergent meromorphic coefficients,
after shrinking the disc, on the generic fibre of \eqref{eq:hybrid-node-chart}.  Along a net with no leading cancellation at the logarithmic scale,
\[
 \lambda(t)\log|f(x_t,y_t)|
 \longrightarrow
 -(-\log r)\min_{a,b}
 \left\{\widehat v_t(c_{ab})+a\alpha+b(1-\alpha)\right\}.
\tag{66.15}\label{eq:hybrid-node-polynomial-limit}
\]
When least terms cancel, one must evaluate the full polynomial rather than
the displayed minimum; this is why hybrid convergence is defined by all
regular functions, not only by coordinatewise rates.  Passing through all
dominating models in \Cref{prop:hybrid-model-criterion} is the
coordinate-free replacement for choosing enough such tests.

\begin{example}[Identical rates with different analytic limits]
\label{ex:hybrid-node-cancellation}
On $xy=t$, the sections $(x_\tau,y_\tau)=(1,\tau)$ and $(1+\tau,\tau/(1+\tau))$ both have limiting rates $(0,1)$. Their $K$-rational limits have $\widehat v(x)=0$ and $\widehat v(y)=1$, but evaluation of $x-1$ gives $+\infty$ for the first and $1$ for the second. At the monomial endpoint the same polynomial has value $0$, since $x$ has transcendental residue there. All three points retract to the same endpoint. This calculation exhibits exactly the cancellation that coordinate rates fail to detect.
\end{example}

\paragraph{Source boundary for chart convergence.}
The hybrid topology and its logarithmic-coordinate criterion are those of
\cite{BJ17Hybrid}.  Equations \eqref{eq:hybrid-node-chart}--
\eqref{eq:hybrid-node-polynomial-limit} are chart calculations inside that
construction.  They do not assert convergence of metrics or measures.

\begin{proposition}\label{prop:hybrid-model-criterion}
Let $\mathcal X$ be a projective model. For each vertical model function on every dominating model, choose a continuous hybrid extension. A net $x_i\in X_{t_i}^{\mathrm{an}}$, with $t_i\to0$, converges to $x\in X_K^{\mathrm{an}}$ if and only if all these test values converge to their values at $x$.
\end{proposition}

\begin{proof}
A vertical divisor gives a difference of model metrics. Choosing smooth
metrics on the model produces its continuous hybrid extension; different
choices differ by $O(\lambda(t))$ near the centre. On the central fibre,
model functions on all dominating models separate points, by the SNC
inverse-limit description in \cite[Section~4.4]{BJ17Hybrid}.
Necessity follows from continuity. For sufficiency, suppose the net fails
to converge to $x$. It has a subnet outside a neighbourhood of $x$.
Compactness provides a convergent further subnet with limit $y$ outside
that neighbourhood. The base projection places $y$ in the central fibre.
Every test has the same value at $y$ and $x$, contrary to separation.
This proves the criterion without replacing nets by sequences.
\end{proof}

\begin{remark}[Central-fibre convention]
\label{rem:hybrid-central-fibre}
The phrase \emph{central fibre} will always mean $X_K^{\mathrm{an}}$ in a hybrid statement. When the scheme-theoretic fibre $\mathcal X_0$ is meant, it is written explicitly.
\end{remark}

\section{Hybrid line bundles and metrics}
\label{sec:hybrid-metrics}

Let $\mathcal L$ be relatively ample on $\mathcal X$.
Its restrictions are denoted $L_t$ and $L_K$.

\begin{definition}[Hybrid metric]
\label{def:hybrid-metric}
A continuous hybrid metric $\Phi$ on $\mathcal L$ is a continuous metric on the analytification of $\mathcal L$ over $\mathcal X^{\mathrm{hyb}}$. Its restrictions are written
\[
                   \phi_t:=\lambda(t)^{-1}\beta_t^*\Phi,
        \qquad
                   \phi^{\NA}:=\Phi|_{X_K^{\mathrm{an}}}.
\tag{66.16}\label{eq:hybrid-metric-restrictions}
\]
Here $\beta_t$ identifies the usual complex fibre with the fibre carrying the powered absolute value. The formula transports additive local norm weights; it unpowers the metric before complex curvature is taken. A hybrid psh metric is a decreasing net limit of tropical hybrid Fubini--Study metrics, with nontrivial generic-fibre restrictions. A continuous metric is semipositive if it belongs to this class; on the compact hybrid space this is equivalent to uniform approximation by tropical hybrid Fubini--Study metrics \cite[Definitions~2.18--2.19, Proposition~2.21]{PS23Hybrid}. Singular psh metrics need not admit uniform approximation.
\end{definition}

The complex normalization hidden in this definition is important.
If a local frame $s$ has norm weight
$\psi_t=-\log|s|_{\phi_t}$, then the continuous hybrid quantity is
\[
                         \lambda(t)\psi_t,
\tag{66.17}\label{eq:hybrid-weight-normalization}
\]
up to the fixed convention $|t|_K=r$. Thus ordinary locally uniform
convergence of the unscaled $\psi_t$ is not the relevant notion.  In the
book's squared-norm convention the complex potential is
$\varphi_t=2\psi_t=-\log|s|_{\phi_t}^{2}$, and its curvature is
$\ddc\varphi_t$.  This factor of two will be kept visible below.

The three scales that occur later can now be written on one line:
\[
 \psi_t=-\log|s|_{\phi_t},
 \qquad
 \varphi_t=-\log|s|_{\phi_t}^{2}=2\psi_t,
 \qquad
 \phi_t^{\mathrm{hyb}}=\lambda(t)\psi_t
       =\frac{\lambda(t)}2\varphi_t.
\tag{66.18}\label{eq:hybrid-three-weight-scales}
\]
Thus a squared-norm potential with asymptotic form
\[
             \varphi_t=2a\log_r|t|+O(1)
\tag{66.19}\label{eq:hybrid-squared-norm-asymptotic}
\]
has central non-Archimedean weight $a$, not $2a$.  Dividing
\eqref{eq:hybrid-squared-norm-asymptotic} by $\log_r|t|$ rather than by
$2\log_r|t|$ is the forbidden factor-two conversion.  With $s=-\log|t|$ and $\lambda_0=-\log r$, division by $2s$ instead gives
$a/\lambda_0$, the weight in the unit valuation normalization. Thus the
radial conversion in \Cref{sec:rf-factor-two-firewall} also changes the
valuation unit when $r\ne\rme^{-1}$.

\begin{definition}[Logarithmic growth]
\label{def:hybrid-log-growth}
A psh metric on $\mathcal L|_{\Delta^*}$ has logarithmic growth if, relative to a smooth metric $\phi_{\mathcal L}$ on a normal projective model, its ordinary norm weight satisfies, locally on the total space,
\[
 \phi-\phi_{\mathcal L}\leq C\log|t|^{-1}+C.
\tag{66.20}\label{eq:hybrid-log-growth}
\]
Equivalently, after adding $a\log|t|$ for some $a\in\mathbf R$, it extends as a psh metric to that model \cite[Definition~3.9]{PS23Hybrid}.
\end{definition}

The definition is independent of the chosen model. It imposes an upper
logarithmic bound, and permits singularities and unbounded negative
weights. A uniform lower logarithmic bound is a stronger condition; it is
not required for the canonical psh extension.

\begin{theorem}\label{thm:hybrid-psh-extension}
Let $\mathcal X\to\Delta$ be a polarized projective degeneration and let $\mathcal L$ be relatively ample. A psh metric on $\mathcal L|_{\Delta^*}$ with logarithmic growth determines a canonical hybrid psh metric. If the resulting metric is continuous, its restrictions in \eqref{eq:hybrid-metric-restrictions} form one continuous metric on $\mathcal X^{\mathrm{hyb}}$.
\end{theorem}

\begin{proof}
This is \cite[Theorem~3.13]{PS23Hybrid}.  We reconstruct the proof far
enough to exhibit the common approximation on the complex and
non-Archimedean fibres.  The multiplier-ideal regularization theorem and
the convergence of the associated ideal functions, proved in the appendix
of that paper and in the references cited there, remain foundational
imported inputs.

\smallskip
\noindent\emph{Finite-section gluing.}
A finite family of sections defines the ordinary complex norm weight
$(2m)^{-1}\log\sum|s_j|^2$ and the central weight
$m^{-1}\max\log|s_j|$. Their hybrid gluing is psh by
\cite[Proposition~3.5 and Lemma~3.6]{PS23Hybrid}; if the sections have no
common zero, it is continuous. The estimate
$0\leq\tfrac12\log\sum\rme^{2u_j}-\max u_j\leq\tfrac12\log N$
becomes an $O(\lambda(t))$ hybrid error. Common zeroes are treated by
decreasing basepoint-free regularizations in that proposition.

\smallskip
\noindent\emph{A synchronized Bergman approximation.}
For a general logarithmic-growth metric, twist by a multiple of
$\log|t|$ and choose an ample model $(\mathscr X,\mathscr L)$ on which
the twisted metric extends as a psh metric.  Put
$\mathscr I_m=\mathscr I(m\phi)$.  The uniform global-generation result
used in \cite[Section~3.3]{PS23Hybrid} supplies one integer $m_0$,
independent of $m$, such that
\[
 \mathcal O_{\mathscr X}
   \bigl((m+m_0)\mathscr L\bigr)\otimes\mathscr I_m
\tag{66.20a}\label{eq:hybrid-uniform-global-generation}
\]
is generated over a fixed slightly larger disk for every $m\geq1$.
Choose a smooth positive model metric $\chi$. Give these sections the
$L^2$ norm with metric $m\phi+m_0\chi$ on the larger disc, and choose a
Hilbert basis $(s_{m,\ell})_{\ell\geq1}$ of that Hilbert space.  The complex Bergman metric and the non-Archimedean ideal metric
are, in compatible additive conventions,
\[
\begin{aligned}
 \phi_m
  &=\frac{1}{2(m+m_0)}
       \log\!\sum_{\ell\geq1}|s_{m,\ell}|^2,\\
 \phi_m^{\NA}
  &=\phi_{\mathscr L}
       +\frac{1}{m+m_0}\phi_{\mathscr I_m}.
\end{aligned}
\tag{66.20b}\label{eq:hybrid-synchronized-bergman}
\]

For a finite partial sum in \eqref{eq:hybrid-synchronized-bergman}, the
complex logarithm of a sum and the non-Archimedean maximum are the two
fibre restrictions of one hybrid Fubini--Study metric.  As the number of
sections tends to infinity, the complex partial sums converge uniformly on
the chosen model.  On the central fibre, the ideals generated by the first
$q$ sections stabilize to $\mathscr I_m$ by noetherianity and global
generation; Dini's lemma then gives uniform convergence there as well.
Consequently, for each fixed $m$, the pair
$(\phi_m,\phi_m^{\NA})$ is the restriction of one hybrid psh metric
$\Phi_m$.  This is the first place where the error must be uniform over
both kinds of fibres: separate fibrewise convergence would not suffice.

\smallskip
\noindent\emph{The common decreasing limit.}
Theorem~A.3 and its proof in \cite{PS23Hybrid} provide a common subsequence
with controlled scalar corrections for which the complex restrictions
decrease to $\phi$. The corrections tend to zero in the hybrid scaling;
the associated ideal-function convergence gives the same decreasing limit
$\phi^{\NA}$ on the central fibre. We use these corrected approximants in
this step. Monotonicity of the uncorrected approximants in every degree is
not required. These are the precise regularization inputs in the proof of
\cite[Theorem~3.13]{PS23Hybrid}.  Hence the pointwise decreasing limit
of the $\Phi_m$ has restrictions
\[
        \beta_t^*\Phi=\lambda(t)\phi_t\quad(t\neq0),
        \qquad
        \Phi|_{X_K^{\mathrm{an}}}=\phi^{\NA}.
\tag{66.20c}\label{eq:hybrid-extension-restrictions-proof}
\]
The class of hybrid psh metrics is closed under such decreasing limits, so
$\Phi$ is the desired canonical hybrid psh extension.  Notice that the
argument exchanges no unrelated limits: first the finite-section limit is
uniform for fixed $m$, and only afterwards is the synchronized decreasing
$m$-limit taken.

If $\Phi$ is continuous, the two restrictions in
\eqref{eq:hybrid-extension-restrictions-proof} are automatically one
continuous metric on the total hybrid space.  This last implication is
formal.  The preceding decreasing approximation does not provide a
quantity $\varepsilon_m\to0$ that is uniform at the scale
$|\log|t||^{-1}$; such an estimate is precisely the missing regularity
in \cite[Question~3.18]{PS23Hybrid}.
\end{proof}

Endpoint continuity does not supply the uniformity missing from the
decreasing approximation.  The question was asked in
\cite[Question~3.18]{PS23Hybrid}, following
\cite[Question~1]{Favre20Hybrid}.  The next result gives a negative answer
even for a smooth trivial degeneration.

\begin{theorem}\label{thm:hybrid-endpoint-continuity-counterexample}
Fix a hybrid radius \(r_0\in(0,1)\).  There exist a smooth trivial
polarized degeneration
\[
 \pi\colon X=\mathbf P^1_{\mathbf C}\times\mathbb D^*
       \longrightarrow\mathbb D^*,
 \qquad
 L=\operatorname{pr}_1^*\mathcal O_{\mathbf P^1}(1),
\]
and a smooth semipositive metric \(\phi\) on \(L\), with logarithmic
growth at the origin, having the following properties.
\begin{enumerate}
\item The metric \(\phi\) extends as an everywhere finite psh metric on
the product model
\[
 \mathcal X=\mathbf P^1_{\mathbf C}\times\mathbb D,
 \qquad
 \mathcal L=\operatorname{pr}_1^*\mathcal O_{\mathbf P^1}(1).
\]
\item Its associated non-Archimedean metric is the continuous product
Fubini--Study model metric:
\[
                   \phi^{\NA}
                   =\phi_{\mathrm{FS}}^{\NA}.
\]
\item Relative to the product Fubini--Study metric, \(\phi\) has zero
fibrewise supremum on every complex fibre and on the non-Archimedean
fibre:
\[
 \sup_{X_t}(\phi_t-\phi_{\mathrm{FS},t})=0
 \quad(t\ne0),
 \qquad
 \sup_{X_{\mathbf C((t))}^{\mathrm{an}}}
 (\phi^{\NA}-\phi_{\mathrm{FS}}^{\NA})=0.
\]
\item The canonical hybrid psh metric \(\phi^{\mathrm{hyb}}\), obtained
by putting \(\phi\) on the complex fibres and
\(\phi^{\NA}\) on the non-Archimedean fibre, is not continuous.
More precisely, for \(p_*=[0:1]\in\mathbf P^1(\mathbf C)\), there are
two sequences \(x_m^+\) and \(x_m^-\) in the complex locus that both
converge to \((p_*)_{\mathbf C((t))}\) and for which
\[
\begin{aligned}
 \lim_{m\to\infty}
 \bigl(\phi^{\mathrm{hyb}}
       -\phi_{\mathrm{FS}}^{\mathrm{hyb}}\bigr)(x_m^+)
   &=\log r_0,\\
 \lim_{m\to\infty}
 \bigl(\phi^{\mathrm{hyb}}
       -\phi_{\mathrm{FS}}^{\mathrm{hyb}}\bigr)(x_m^-)
   &=0.
\end{aligned}
\]
\end{enumerate}
Consequently, continuity of a semipositive metric on the punctured total
space together with continuity of its associated non-Archimedean metric
does not imply continuity of the canonical hybrid psh extension, even
under the displayed zero-supremum normalizations.

Moreover, relative to the product Fubini--Study metric, the resulting
potential satisfies separately all three conditions in Favre's
Question~1: it is continuous near every punctured
\(\overline X_r\), its curvature is the restriction of a positive
closed \((1,1)\)-current on a model neighborhood, and its
generic-Lelong function extends continuously to the non-Archimedean
analytification.  Nevertheless, that potential is not uniform in the
sense of \cite[Definition~4.1]{Favre20Hybrid}.
\end{theorem}

\begin{proof}
For \(j\geq1\), put
\[
 s_j=2^j,\qquad
 \rho_j=\exp(-s_j),\qquad
 c_j=4^{-j},\qquad
 \epsilon_j=\exp(-s_j/c_j)=\exp(-8^j),
\]
and define on the unit disk
\[
 q_j(t)=\frac{c_j}{2}
   \log\frac{|t-\rho_j|^2+\epsilon_j^2}{5}.
\tag{66.20d}\label{eq:p18-qj}
\]
The functions \(q_j\) are smooth and subharmonic.  Indeed,
\(\log(|t-a|^2+\epsilon^2)\) has nonnegative
\(\mathrm d\mathrm d^c\), and multiplication by \(c_j/2>0\)
preserves this property.  Moreover \(q_j\leq0\) on \(\mathbb D\):
for \(|t|<1\), the numerator in \eqref{eq:p18-qj} is strictly smaller
than \(5\).

Set
\[
                         u(t)=\sum_{j=1}^{\infty}q_j(t).
\tag{66.20e}\label{eq:p18-u}
\]
Let \(K\Subset\mathbb D^*\).  For all sufficiently large \(j\), the
points \(\rho_j\) and the disks of radius \(\epsilon_j\) about them
have a fixed positive distance from \(K\).  Every derivative, on \(K\),
of the logarithm occurring in \eqref{eq:p18-qj} is then bounded by a
constant depending only on \(K\) and the order of the derivative.
Since \(\sum_jc_j<\infty\), the series in \eqref{eq:p18-u} converges in
\(C^\infty(K)\).  Hence \(u\) is smooth and subharmonic on
\(\mathbb D^*\).

At the origin one has
\[
\begin{aligned}
 q_j(0)
 &=
 \frac{c_j}{2}
 \left(
  2\log\rho_j-\log5+
  \log\left(1+\frac{\epsilon_j^2}{\rho_j^2}\right)
 \right)\\
 &=-c_js_j-\frac{c_j}{2}\log5+O(c_j)
  =-2^{-j}+O(4^{-j}).
\end{aligned}
\tag{66.20f}\label{eq:p18-at-zero}
\]
Thus \(\sum_jq_j(0)>-\infty\).  The partial sums in
\eqref{eq:p18-u} form a decreasing sequence of subharmonic functions on
\(\mathbb D\), because every \(q_j\leq0\).  Their limit is not
identically \(-\infty\), by \eqref{eq:p18-at-zero}; therefore the
decreasing-limit theorem gives a subharmonic extension of \(u\) to
\(\mathbb D\), with the finite value prescribed by
\eqref{eq:p18-at-zero}.

This extension has zero Lelong number at the origin.  To see this without
using any pointwise radial assertion, let \(M_u(r)\) be its circular
mean.  The submean inequality and \(u\leq0\) give
\[
                         u(0)\leq M_u(r)\leq0.
\]
The circular-mean characterization of the one-variable Lelong number
therefore yields
\[
                       \nu_0(u)
                       =\lim_{r\to0}\frac{M_u(r)}{\log r}=0.
\tag{66.20g}\label{eq:p18-zero-lelong}
\]

We next establish the logarithmic-growth estimate and isolate two
different normalized limits.  Let
\[
                B_j=\{\,t\in\mathbb D:
                         |t-\rho_j|<\rho_j/2\,\}.
\]
The disks \(B_j\) are pairwise disjoint.  Indeed,
\(\rho_{j+1}=\rho_j^2\), \(\rho_1<1/3\), and hence
\[
                    \frac32\rho_{j+1}
                    <\frac12\rho_j .
\]
If \(t\notin B_j\), then
\[
\begin{aligned}
 q_j(t)
 &\geq\frac{c_j}{2}
       \log\frac{\rho_j^2}{20}\\
 &=-c_js_j-\frac{c_j}{2}\log20.
\end{aligned}
\tag{66.20h}\label{eq:p18-outside-well}
\]
If \(t\) belongs to none of the \(B_j\), summing
\eqref{eq:p18-outside-well} gives \(u(t)\geq-C\), since
\[
                \sum_{j\geq1}c_js_j=\sum_{j\geq1}2^{-j}<\infty.
\]
If \(t\in B_m\), pairwise disjointness permits the use of
\eqref{eq:p18-outside-well} for all \(j\ne m\), while
\[
\begin{aligned}
 q_m(t)
 &\geq\frac{c_m}{2}\log\frac{\epsilon_m^2}{5}\\
 &=-s_m-\frac{c_m}{2}\log5.
\end{aligned}
\tag{66.20i}\label{eq:p18-inside-well}
\]
Moreover \(-\log|t|=s_m+O(1)\) uniformly for \(t\in B_m\).
Equations \eqref{eq:p18-outside-well}--\eqref{eq:p18-inside-well},
together with \(u\leq0\), prove that, after increasing \(C\),
\[
             -\log|t|^{-1}-C\leq u(t)\leq0
             \qquad(0<|t|<1).
\tag{66.20j}\label{eq:p18-log-growth}
\]

At \(t=\rho_m\), the \(m\)-th summand is exactly
\[
                  q_m(\rho_m)
                  =-s_m-\frac{c_m}{2}\log5.
\]
Since \(\rho_m\notin B_j\) for \(j\ne m\), the sum of all other
summands lies between \(-C\) and \(0\), by
\eqref{eq:p18-outside-well}.  It follows that
\[
                         u(\rho_m)=-s_m+O(1).
\tag{66.20k}\label{eq:p18-plus-sequence}
\]
On the other hand, \(-\rho_m\) belongs to none of the \(B_j\), because
\(|-\rho_m-\rho_j|=\rho_m+\rho_j>\rho_j/2\).  Hence
\[
                         -C\leq u(-\rho_m)\leq0.
\tag{66.20l}\label{eq:p18-minus-sequence}
\]
Since \(\log|\rho_m|=-s_m\), equations
\eqref{eq:p18-plus-sequence}--\eqref{eq:p18-minus-sequence} give
\[
\begin{aligned}
 \frac{\log r_0}{\log|\rho_m|}\,u(\rho_m)
   &\longrightarrow\log r_0,\\
 \frac{\log r_0}{\log|-\rho_m|}\,u(-\rho_m)
   &\longrightarrow0.
\end{aligned}
\tag{66.20m}\label{eq:p18-two-limits}
\]

We now put this scalar example on a polarized family without adding it as
a scalar pulled back from the base.  Let
\[
 \mathcal X=\mathbf P^1_{\mathbf C}\times\mathbb D,\qquad
 \mathcal L=\operatorname{pr}_1^*\mathcal O_{\mathbf P^1}(1),
\]
and use homogeneous coordinates \([Z_0:Z_1]\) on \(\mathbf P^1\).
The product Fubini--Study model metric and our new metric are represented
by the homogeneous additive weights
\[
\begin{aligned}
 \phi_{\mathrm{FS}}([Z_0:Z_1],t)
   &=\frac12\log\bigl(|Z_0|^2+|Z_1|^2\bigr),\\
 \phi([Z_0:Z_1],t)
   &=\frac12\log\bigl(|Z_0|^2+\exp(2u(t))|Z_1|^2\bigr).
\end{aligned}
\tag{66.20n}\label{eq:p18n-metrics}
\]
The second formula is smooth on
\(X=\mathcal X|_{\mathbb D^*}\).  It also defines an everywhere finite
psh metric on \(\mathcal L\) over \(\mathcal X\).  Indeed, on the chart
\(Z_1\ne0\), with \(z=Z_0/Z_1\), its local weight is
\[
 F_1(z,t)=\frac12\log\bigl(|z|^2+\exp(2u(t))\bigr),
\tag{66.20o}\label{eq:p18n-weight-one}
\]
which is the logarithmic sum of the exponentials of the two psh
functions \(\log|z|\) and \(u(t)\).  On the chart \(Z_0\ne0\), with
\(w=Z_1/Z_0\), its local weight is
\[
 F_0(w,t)=\frac12\log\bigl(1+\exp(2u(t))|w|^2\bigr),
\tag{66.20p}\label{eq:p18n-weight-zero}
\]
the same construction applied to \(0\) and \(u(t)+\log|w|\).
The convex coordinatewise nondecreasing log-sum-exp function preserves
plurisubharmonicity.  Thus both local weights are psh; their transition
law is the one for \(\mathcal O_{\mathbf P^1}(1)\).  On the punctured
total space they are smooth, so \(\phi\) is a smooth semipositive metric.

Its potential relative to the Fubini--Study metric is the globally
defined function
\[
 \psi([Z_0:Z_1],t)
 =\phi-\phi_{\mathrm{FS}}
 =\frac12\log
   \frac{|Z_0|^2+\exp(2u(t))|Z_1|^2}
        {|Z_0|^2+|Z_1|^2}.
\tag{66.20q}\label{eq:p18n-relative-potential}
\]
Since \(u\leq0\), one has
\[
                         u(t)\leq\psi\leq0.
\tag{66.20r}\label{eq:p18n-potential-bounds}
\]
The upper bound is immediate.  For the lower bound, use
\(|Z_0|^2\geq\exp(2u(t))|Z_0|^2\).  In particular, the extension just
constructed proves logarithmic growth by the first equivalent condition
in \cite[Definition~3.9]{PS23Hybrid}, while
\eqref{eq:p18-log-growth} and \eqref{eq:p18n-potential-bounds} verify
the two-sided logarithmic bound directly.  Moreover,
\(\psi([1:0],t)=0\), and hence
\[
                   \sup_{X_t}\psi_t=0\qquad(t\ne0).
\tag{66.20s}\label{eq:p18n-complex-normalization}
\]

We compute the associated non-Archimedean metric.  Let \(E\) be a prime
vertical divisor on an arbitrary model dominating \(\mathcal X\).  At a
general smooth point of \(E\), one may write
\[
                          t=\zeta^{b_E}h,
\]
where \(\zeta=0\) is a local equation for \(E\), \(h\) is a unit, and
\(b_E=\operatorname{ord}_E(t)\).  In one complex variable the Lelong
number is the atomic mass at the origin of the Riesz measure, up to the
fixed normalization.  Pullback by a holomorphic map of order \(b_E\)
multiplies this atomic mass by \(b_E\).  Thus
\[
                 \nu_E(u\circ\pi)
                 =b_E\nu_0(u)=0
\tag{66.20t}\label{eq:p18n-base-divisorial-zero}
\]
by \eqref{eq:p18-zero-lelong}.

Choose one of the two affine charts containing the image of the general
point of \(E\).  On \(Z_0\ne0\), the psh weight
\eqref{eq:p18n-weight-zero} is bounded below by \(0\), so its generic
Lelong number along \(E\) is zero.  On \(Z_1\ne0\), the psh weight
\eqref{eq:p18n-weight-one} is bounded below by \(u\circ\pi\).
Monotonicity of Lelong numbers, together with
\eqref{eq:p18n-base-divisorial-zero}, again gives
\[
                         \nu_E(\phi)=0.
\tag{66.20u}\label{eq:p18n-metric-divisorial-zero}
\]
The smooth Fubini--Study reference has zero generic Lelong numbers as
well.  \cite[Theorem~3.10]{PS23Hybrid} therefore gives
\[
                  \phi^{\NA}
                  =\phi_{\mathrm{FS}}^{\NA},
\tag{66.20v}\label{eq:p18n-na-model}
\]
the continuous product model metric.  In particular,
\[
 \sup_{X_{\mathbf C((t))}^{\mathrm{an}}}
 (\phi^{\NA}-\phi_{\mathrm{FS}}^{\NA})=0.
\tag{66.20w}\label{eq:p18n-na-normalization}
\]

\cite[Theorem~3.13]{PS23Hybrid} now applies to \(\phi\): the metric and
the endpoint \eqref{eq:p18n-na-model} form one canonical hybrid psh
metric \(\phi^{\mathrm{hyb}}\).  This invokes genuine hybrid psh
gluing, rather than an arbitrary assignment of values on two disjoint
sets.

Put \(p_*=[0:1]\), and write \((p_*)_{\mathbf C((t))}\) for the point
on the non-Archimedean fibre induced by the constant algebraic section.
Functoriality of hybrid analytification, applied to that section, gives
\[
\begin{aligned}
 x_m^+&=(p_*,\rho_m)\longrightarrow(p_*)_{\mathbf C((t))},\\
 x_m^-&=(p_*,-\rho_m)\longrightarrow(p_*)_{\mathbf C((t))}.
\end{aligned}
\tag{66.20x}\label{eq:p18n-hybrid-convergence}
\]
The hybrid metric induced by the product Fubini--Study model is
continuous.  On \(p_*\) the relative potential in
\eqref{eq:p18n-relative-potential} is exactly
\[
                         \psi(p_*,t)=u(t).
\tag{66.20y}\label{eq:p18n-special-section}
\]
In additive weights, the difference on a complex fibre over \(t\ne0\)
is the usual hybrid rescaling
\[
 \bigl(\phi^{\mathrm{hyb}}
       -\phi_{\mathrm{FS}}^{\mathrm{hyb}}\bigr)|_{X_t}
 =\frac{\log r_0}{\log|t|}\,\psi_t,
\tag{66.20z}\label{eq:p18n-hybrid-normalization}
\]
whereas its difference on the central non-Archimedean fibre is zero by
\eqref{eq:p18n-na-model}.  Combining
\eqref{eq:p18-two-limits}, \eqref{eq:p18n-hybrid-convergence},
\eqref{eq:p18n-special-section}, and
\eqref{eq:p18n-hybrid-normalization} proves the two limits in the
statement.  Since \(\log r_0\ne0\), the canonical hybrid psh metric is
not continuous despite the normalizations
\eqref{eq:p18n-complex-normalization} and
\eqref{eq:p18n-na-normalization}.

It remains to check Favre's formulation independently.  Take
\(\varphi=\psi\).  The curvature current
\[
                         T=c_1(\mathcal L,\phi)
\tag{66.20aa}\label{eq:p18n-favre-current}
\]
is positive and closed on the product model because \(\phi\) is a psh
metric, and on the punctured total space it restricts to
\(\omega_{\mathrm{FS}}+\mathrm d\mathrm d^c\varphi\).  The potential
is smooth on \(X\), hence continuous near every punctured
\(\overline X_r\).  Equation \eqref{eq:p18n-metric-divisorial-zero}
says that Favre's generic-Lelong function is \(g_T\equiv0\), so it
extends continuously to the non-Archimedean analytification.  All three
bullets of Favre's Question~1 hold.

Suppose, for contradiction, that \(\varphi\) were uniform.  The remark
following \cite[Definition~4.1]{Favre20Hybrid} states that its normalized
function on the punctured hybrid locus would then extend continuously to
the hybrid space.  For the hybrid radius appearing in that definition,
the two sequences in \eqref{eq:p18n-hybrid-convergence} have the distinct
limits computed above.  Such a continuous extension is impossible.  Thus
\(\varphi\) is not uniform, completing the proof.
\end{proof}

The original base-scalar packet for
\Cref{thm:hybrid-endpoint-continuity-counterexample} passed production
Rethlas run \texttt{20260825T030402Z\_}\allowbreak
\texttt{13e01361eeff}.  The stronger zero-supremum theorem and proof above
are controlled by production run
\texttt{20260825T035945Z\_}\allowbreak
\texttt{070e264c3ef1}, item
\texttt{pi\_db89dd432acabcf33a3f2272}, using
\texttt{gpt-5.6-sol} at \texttt{xhigh}.  Both runs returned final verdict
\texttt{correct}, with zero critical errors and zero gaps.  After the
notation and cross-reference adaptation, the exact theorem-and-proof bytes
displayed above passed the independent production run
\texttt{20260825T043840Z\_}\allowbreak
\texttt{d77f7f93aeb0}, item
\texttt{pi\_b7f015a6ad5570f431769a71}, with the same model and scrutiny and
again with final verdict \texttt{correct}, zero critical errors, and zero
gaps.

\begin{remark}\label{rem:hybrid-extension-not-continuity}
The obstruction in \Cref{thm:hybrid-endpoint-continuity-counterexample} is fibre dependent and already has zero supremum relative to a fixed continuous model metric on every complex fibre and on the non-Archimedean fibre.  Thus neither endpoint continuity nor this normalization forces hybrid continuity.  If, more strongly, the difference from a continuous model metric is locally bounded near the central fibre, then compactness gives a uniform bound on a model neighbourhood of the central fibre, and multiplication by $\log r_0/\log|t|$ sends that difference uniformly to zero, so the canonical hybrid extension is continuous at the non-Archimedean fibre.
\end{remark}

\begin{example}[A model metric]
\label{ex:hybrid-model-metric}
Let $\mathcal D$ be a vertical Cartier divisor on a model $\mathcal X'\to\mathcal X$. The model metric on $L_K$ comes from $\mathcal L+\mathcal D$. Its difference $\phi_{\mathcal D}^{\NA}$ from the fixed model metric of $\mathcal L$ is the vertical model function. If $\mathcal X'_0=\sum_E b_EE$, then the normalized divisorial valuation is $\widehat v_E=b_E^{-1}\ord_E$, while the corresponding analytic valuation is $v_E=(-\log r)\widehat v_E$.  At that analytic point,
\[
             \phi_{\mathcal D}^{\NA}(v_E)
             =(-\log r)b_E^{-1}\ord_E(\mathcal D),
\tag{66.21}\label{eq:hybrid-model-value}
\]
The sign is the book's convention: an effective vertical divisor adds its positive valuation. For example $\mathcal D=a\mathcal X'_0$ gives the constant $a(-\log r)$ at every vertex. For a rational divisor the same formula is divided by its Cartier denominator.
\end{example}

\subsection{A finite polynomial functional across the endpoint}
\label{sec:hybrid-polynomial-functional}

At a point $x$ of a Banach analytic base, the number
\[
             \rho(x)=\frac{\log\max\{1,|2|_x\}}{\log 2}
\]
records the power of the ordinary absolute value on an Archimedean fibre;
it is zero precisely on the non-Archimedean locus.  On the hybrid annulus
above, $\rho=\lambda(t)$ at a complex point and $\rho=0$ at the central
fibre.

Later finite-level constructions reduce locally to multihomogeneous
polynomials whose coefficients vary analytically, sometimes as sections of
a line bundle.  What is needed at the endpoint is a continuous local weight
that is the scaled logarithmic sphere functional for $\rho>0$, becomes the
weighted Gauss maximum for $\rho=0$, and has the correct frame-change law.
\Cref{thm:hybrid-weighted-philippon-gauss} supplies exactly this finite
polynomial input.  It does not construct or compare any later relative
resultant tensor.

\begingroup
\edef\savedWeightedPGEquation{\number\value{equation}}
\renewcommand{\theequation}{66.W\arabic{equation}}
\renewcommand{\theHequation}{66.W.\arabic{equation}}
\setcounter{equation}{0}
\begin{theorem}\label{thm:hybrid-weighted-philippon-gauss}
Let \(A\) be a Banach ring and let \(U\subset\mathcal M(A)\) be open.
Fix finitely many blocks
\[
 \mathbf U_i=(U_{i0},\ldots,U_{i,N_i-1}),\qquad N_i\geq 1,
\]
nonnegative block degrees \(d_i\), a finite set \(S\) of multiexponents
of block degree \((d_i)_i\), and real numbers \(w_\alpha\),
\(\alpha\in S\).  Let \(c_\alpha\in\mathcal O(U)\), and suppose that
the tuple \((c_\alpha)_{\alpha\in S}\) is nowhere zero.  Locally on \(U\),
the finitely many \(c_\alpha\) belong to one Banach algebra of analytic
functions on a Banach analytic domain.

For \(x\in U\), set
\begin{equation}\label{eq:hybrid-philippon-rho}
 \rho(x):=\frac{\log\max\{1,|2|_x\}}{\log 2}.
\end{equation}
If \(\rho(x)=0\), put
\begin{equation}\label{eq:hybrid-philippon-na-value}
 \Phi(x):=\max_{\alpha\in S}
 \bigl(\log|c_\alpha(x)|_x+w_\alpha\bigr),
 \qquad \log 0:=-\infty.
\end{equation}
If \(\rho(x)>0\), choose a Banach analytic neighbourhood \(V\) of
\(x\), a Banach algebra
\(A_V\) of analytic functions on \(V\) containing all
\(c_\alpha|_V\), and a character
\(\chi_x\colon A_V\to\mathbf C\) representing the seminorm at \(x\),
so that
\[
 |a|_x=|\chi_x(a)|_\infty^{\rho(x)}.
\]
Define
\begin{equation}\label{eq:hybrid-philippon-arch-polynomial}
 F_x:=\sum_{\alpha\in S}
 \chi_x(c_\alpha)\mathrm e^{w_\alpha/\rho(x)}\mathbf U^\alpha
\end{equation}
and
\begin{equation}\label{eq:hybrid-philippon-arch-value}
 \Phi(x):=\rho(x)\operatorname{ph}(F_x),
\end{equation}
where \(\operatorname{ph}\) is Philippon's logarithmic sphere measure,
including its harmonic correction.  Then \(\Phi\) is independent of all
choices and is a continuous real-valued function on \(U\).

The construction is homogeneous.  More precisely, if
\(\lambda\) is a nowhere-vanishing analytic function, then
\begin{equation}\label{eq:hybrid-philippon-homogeneity}
 \Phi(\lambda F)=\log|\lambda|+\Phi(F).
\end{equation}
Consequently, let \(\mathscr L\) be an analytic line bundle on \(U\) and
let
\[
 R=\sum_{\alpha\in S}r_\alpha\mathbf U^\alpha
\]
be an \(\mathscr L\)-valued polynomial whose fibrewise coefficient tuple
is nowhere zero.  Write \(R=eF_e\) in a local frame \(e\) of
\(\mathscr L\).  If \(e'=ue\), then
\begin{equation}\label{eq:hybrid-philippon-frame-change}
 F_{e'}=u^{-1}F_e,
 \qquad
 \Phi_{e'}=\Phi_e-\log|u|.
\end{equation}
Thus the local weights \(\Phi_e\) define a continuous additive metric on
\(\mathscr L\).
\end{theorem}

\begin{proof}
We first establish the two facts about Archimedean seminorms that will be
needed in the continuity argument.  Let \(R\) be a commutative ring with a
multiplicative seminorm \(x\), normalized by \(|1|_x=1\), and set
\(\mathfrak p=\ker(x)\) and
\(K=\operatorname{Frac}(R/\mathfrak p)\).  If \(|2|_x\leq 1\), write
\(n^m\) in base \(2\).  The number of nonzero binary digits is \(O(m)\),
whereas every power of \(2\) has absolute value at most \(1\).  The ordinary
triangle inequality gives \(|n|_x^m=O(m)\), and taking \(m\)-th roots shows
that \(|n|_x\leq 1\) for every positive integer \(n\).  Expanding
\((a+b)^m\), using
\(\bigl|\binom mj\bigr|_x\leq 1\), and then taking \(m\)-th roots gives
\[
 |a+b|_x\leq\max\{|a|_x,|b|_x\}.
\]
Thus \(x\) induces a non-Archimedean absolute value on \(K\).  The converse
follows from \(|2|_x=|1+1|_x\leq 1\).  In particular,
\(\rho(x)=0\) exactly on the non-Archimedean locus.

Suppose that \(|2|_x>1\).  The base-expansion proof of Ostrowski's
normalization gives, for integers \(a,b\geq 2\),
\[
 |a|_x^m\leq C_b(1+m\log_ba)
 \max\{1,|b|_x^{m\log_ba+O(1)}\},
 \qquad
 C_b:=\max_{0\leq j<b}|j|_x.
\]
Taking \(m\)-th roots and letting \(m\to\infty\), and then interchanging
\(a\) and \(b\), yields
\[
 |n|_x=n^{\rho(x)}\qquad(n\geq 1).
\]
The triangle inequality gives \(|2|_x\leq 2\), so
\(0<\rho(x)\leq 1\).  Put \(v(a):=|a|_x^{1/\rho(x)}\) on \(K\).  This is
multiplicative and restricts to the usual absolute value on \(\mathbf Q\).
It also satisfies the ordinary triangle inequality.  Indeed, with
\(A_0=v(a)\) and \(B_0=v(b)\), the binomial expansion gives
\[
 |a+b|_x^m
 \leq\sum_{j=0}^m
 \bigl(\tbinom mj A_0^jB_0^{m-j}\bigr)^{\rho(x)}
 \leq(m+1)(A_0+B_0)^{m\rho(x)}.
\]
Taking the power \(1/(m\rho(x))\) and letting \(m\to\infty\) proves
\(v(a+b)\leq A_0+B_0\).

Let \(\widehat K\) be the completion for \(v\).  The closure of
\(\mathbf Q\) in \(\widehat K\) is \(\mathbf R\), and \(\widehat K\) is a
commutative real Banach division algebra.  We recall the short
Gelfand--Mazur argument in the form needed here.  In a complex Banach
division algebra the spectrum of every element is nonempty: otherwise,
after applying continuous linear functionals, its resolvent would be a
bounded entire function tending to \(0\) at infinity, contrary to
Liouville's theorem.  A spectral value then forces the element to be a
complex scalar.  Hence every complex Banach division algebra is
\(\mathbf C\).  If a commutative real Banach division algebra \(E\)
contains a square root of \(-1\), this complex statement gives
\(E=\mathbf C\).  Otherwise \(E[T]/(T^2+1)\), equipped with any equivalent
complete finite-module algebra norm, is a complex Banach division algebra
and hence is \(\mathbf C\).  The closed subfield \(E\subset\mathbf C\)
contains \(\mathbf R\) but not \(\mathrm i\), and therefore equals
\(\mathbf R\).  It follows that \(\widehat K\) is isometrically
\(\mathbf R\) or \(\mathbf C\).  Composing
\(R\to K\to\widehat K\hookrightarrow\mathbf C\) produces a character
\(\chi_x\) satisfying
\begin{equation}\label{eq:hybrid-philippon-character}
 |a|_x=|\chi_x(a)|_\infty^{\rho(x)}.
\end{equation}

This character is unique up to conjugation.  In fact, suppose that
\(\sigma,\tau\colon R\to\mathbf C\) satisfy
\(|\sigma(a)|=|\tau(a)|\) for every \(a\).  They have the same kernel and
induce a modulus-preserving field isomorphism
\[
 \operatorname{Frac}(\sigma(R))\longrightarrow
 \operatorname{Frac}(\tau(R)),
 \qquad
 \frac{\sigma(a)}{\sigma(b)}\longmapsto
 \frac{\tau(a)}{\tau(b)}.
\]
It extends isometrically to the closures \(L,M\subset\mathbf C\).  A closed
subfield of \(\mathbf C\) containing \(\mathbf Q\) contains \(\mathbf R\);
it equals \(\mathbf R\) if it has no nonreal element and equals
\(\mathbf C\) otherwise.  In the real case the extension is the identity.
In the complex case it is a continuous automorphism of \(\mathbf C\)
fixing \(\mathbf R\), so it sends \(\mathrm i\) to \(\mathrm i\) or
\(-\mathrm i\) and is the identity or complex conjugation.  This proves the
asserted uniqueness in \eqref{eq:hybrid-philippon-character}.

We next record the corresponding topology of characters.  Let \(B\) be a
finitely generated \(\mathbf Z\)-algebra, equipped with the pointwise
topology on its multiplicative seminorms, and write
\[
 B=\mathbf Z[t_1,\ldots,t_r]/I.
\]
Suppose that a net of Archimedean seminorms \(x_\lambda\to x\) satisfies
\(\rho(x_\lambda)\geq\varepsilon>0\), and choose arbitrary representing
characters \(\chi_\lambda\).  Since
\(|2|_{x_\lambda}\to|2|_x\), one has
\(\rho(x_\lambda)\to\rho(x)\).  The values
\(|t_j|_{x_\lambda}\) are bounded, say by \(K_j\), and hence
\[
 |\chi_\lambda(t_j)|
 =|t_j|_{x_\lambda}^{1/\rho(x_\lambda)}
 \leq\max\{1,K_j^{1/\varepsilon}\}.
\]
The points
\((\chi_\lambda(t_1),\ldots,\chi_\lambda(t_r))\) therefore lie in one
compact polydisc.  From every subnet we may extract a further subnet
converging to \(z\).  Every polynomial in \(I\) vanishes at \(z\), so
evaluation at \(z\) defines a character
\(\eta\colon B\to\mathbf C\).  If \(b\in B\) is represented by a polynomial
\(P\), then
\[
 \chi_\lambda(b)=P(\chi_\lambda(t_1),\ldots,
 \chi_\lambda(t_r))\longrightarrow P(z)=\eta(b).
\]
Consequently
\[
 |b|_x=\lim_\lambda|b|_{x_\lambda}
 =\lim_\lambda|\chi_\lambda(b)|^{\rho(x_\lambda)}
 =|\eta(b)|^{\rho(x)}.
\]
The uniqueness just proved shows that \(\eta\) is \(\chi_x\) or
\(\overline{\chi_x}\).  Thus every subnet of the chosen characters has a
further subnet converging pointwise on \(B\) to one of these two
characters.

We now turn to the Archimedean polynomial functional.  For a nonzero
multihomogeneous complex polynomial \(F\) in fixed blocks and fixed degrees,
D'Andrea--Krick--Sombra prove
\[
 \bigl|m(F)-\log\|F\|_{\max}\bigr|
 \leq\sum_i\log(N_i+1)d_i
\]
and
\[
 0\leq\operatorname{ph}(F)-m(F)
 \leq\sum_i\log(N_i+1)d_i;
\]
see \cite[Lemma~2.32, Definition~2.33, and (2.34)--(2.35)]{DKS13Heights}.
Here \(m(F)\) is the grouped-variable Mahler measure and
\(\|F\|_{\max}\) is the maximum coefficient modulus.  Thus, with the
harmless fixed-type constant
\begin{equation}\label{eq:hybrid-philippon-fixed-bound}
 C:=2\sum_i\log(N_i+1)d_i,
\end{equation}
one has
\begin{equation}\label{eq:hybrid-philippon-coefficient-bound}
 \bigl|\operatorname{ph}(F)-\log\|F\|_{\max}\bigr|\leq C.
\end{equation}

We also need continuity in the coefficients.  Zhang proves continuity and
positivity of the Chow metric, and hence of its logarithmic spherical
integral, on the nonzero projective coefficient space of a fixed equal
Chow type \cite[Section~1.5 and Theorem~3.6(1)]{Zhang96Semistable}.
The reduction from arbitrary fixed block sizes and multidegrees is
elementary.  If every \(d_i=0\), the coefficient space is one-dimensional
and continuity follows directly from
\(\operatorname{ph}(a)=\log|a|_\infty\).  Otherwise put
\(Q=\max_iN_i\) and \(D=\max_i d_i>0\).  Pad the \(i\)-th
block with \(Q-N_i\) unused variables and send
\begin{equation}\label{eq:hybrid-philippon-padding}
 F\longmapsto\iota(F):=F\prod_iU_{i0}^{D-d_i}.
\end{equation}
This is an injective linear map into the equal-type coefficient space.  On
the unit sphere in \(\mathbf C^Q\), write the first \(N_i\) coordinates as
\(r_i\xi_i\), where \(\xi_i\) is uniformly distributed on the unit sphere
in \(\mathbf C^{N_i}\) and is independent of \(r_i\).  Multihomogeneity
gives, off a null set,
\[
 \log|\iota(F)|=
 \log|F(\xi_1,\ldots,\xi_m)|
 +\sum_i\bigl(D\log r_i+(D-d_i)\log|\xi_{i0}|\bigr).
\]
The added logarithms are integrable and their integrals depend only on the
fixed type.  After including the fixed harmonic corrections, we obtain
\[
 \operatorname{ph}_{\mathrm{equal}}(\iota(F))
 =\operatorname{ph}(F)+\kappa
\]
with \(\kappa\) independent of \(F\).  Zhang's result therefore proves
projective coefficient continuity for the original type.  Homogeneity of
the normalized sphere integral lifts this to the punctured coefficient
vector space.  The same formula also gives
\begin{equation}\label{eq:hybrid-philippon-complex-homogeneity}
 \operatorname{ph}(\lambda F)=\log|\lambda|_\infty+
 \operatorname{ph}(F),
 \qquad
 \operatorname{ph}(\overline F)=\operatorname{ph}(F).
\end{equation}
In particular, \(\operatorname{ph}(F)\) is finite, continuous, homogeneous,
and invariant under complex conjugation on every punctured fixed-type
coefficient space.

Continuity is local on \(U\).  We may therefore work on a Banach analytic
domain \(V\subset U\) on which all coefficients belong to one Banach
algebra \(A_V\).  If a rational localization is used, include its finitely
many denominators and their inverses in what follows; they are units on
\(V\), so evaluation of the inverse is continuous.  Let \(B\subset A_V\)
be the finitely generated \(\mathbf Z\)-algebra generated by the
coefficients and these localization elements.  Define
\begin{equation}\label{eq:hybrid-philippon-max-function}
 M(x):=\max_{\alpha\in S}
 \bigl(\log|c_\alpha(x)|_x+w_\alpha\bigr)
 =\log\max_{\alpha\in S}
 \bigl(\mathrm e^{w_\alpha}|c_\alpha(x)|_x\bigr).
\end{equation}
The quantity inside the last logarithm is a finite maximum of continuous
nonnegative functions and is everywhere positive by hypothesis.  Hence
\(M\) is a finite continuous function, even where some coefficients vanish.

If \(\rho(x)>0\), the character construction applied to \(A_V\) gives
\(\chi_x\).  Only its restriction to \(B\) enters
\eqref{eq:hybrid-philippon-arch-polynomial}.  The polynomial \(F_x\) is
nonzero.  Changing \(\chi_x\) conjugates all coefficients simultaneously,
so \eqref{eq:hybrid-philippon-complex-homogeneity} shows that
\eqref{eq:hybrid-philippon-arch-value} is independent of this choice.

Let \(x_\lambda\to x\) with \(\rho(x)>0\).  Eventually
\(\rho(x_\lambda)\geq\rho(x)/2\), and all these points are Archimedean.
Choose arbitrary representing characters.  The finite-type character
argument shows that every subnet has a further subnet whose characters on
\(B\) converge to \(\chi_x\) or \(\overline{\chi_x}\).  Along it, for each
\(\alpha\),
\[
 \chi_{x_\lambda}(c_\alpha)
 \mathrm e^{w_\alpha/\rho(x_\lambda)}
 \longrightarrow
 \chi_x(c_\alpha)\mathrm e^{w_\alpha/\rho(x)}
\]
or all limits are conjugated.  Coefficient continuity and conjugation
invariance of \(\operatorname{ph}\) give
\[
 \rho(x_\lambda)\operatorname{ph}(F_{x_\lambda})
 \longrightarrow\rho(x)\operatorname{ph}(F_x).
\]
A net converges to a point when every subnet has a further subnet converging
to that point: otherwise one can choose a subnet outside a fixed
neighbourhood.  Thus \(\Phi\) is continuous on the locus \(\rho>0\).

It remains to cross its boundary.  For \(\rho(x)>0\), one has the exact
identity
\begin{align}
 \rho(x)\log\|F_x\|_{\max}
 &=\rho(x)\max_\alpha
 \left(\log|\chi_x(c_\alpha)|_\infty+
 \frac{w_\alpha}{\rho(x)}\right)\notag\\
 &=\max_\alpha\bigl(\log|c_\alpha(x)|_x+w_\alpha\bigr)
 =M(x).
 \label{eq:hybrid-philippon-exact-max}
\end{align}
This remains valid when some coefficients are zero, with
\(\log0=-\infty\).  Combining
\eqref{eq:hybrid-philippon-coefficient-bound} and
\eqref{eq:hybrid-philippon-exact-max} yields
\begin{equation}\label{eq:hybrid-philippon-boundary-estimate}
 |\Phi(x)-M(x)|\leq C\rho(x).
\end{equation}
If \(x_\lambda\to x\) and \(\rho(x)=0\), the left side of
\eqref{eq:hybrid-philippon-boundary-estimate} is zero at every
non-Archimedean \(x_\lambda\), while the displayed estimate applies to
every Archimedean \(x_\lambda\).  Since \(M(x_\lambda)\to M(x)\) and
\(\rho(x_\lambda)\to0\), we obtain
\(\Phi(x_\lambda)\to M(x)=\Phi(x)\).  Conversely, continuity of \(\rho\)
prevents a net with \(\rho=0\) from converging to a point with \(\rho>0\).
This proves continuity everywhere.

If \(\lambda\ne0\), then on the non-Archimedean locus
\[
 \Phi(\lambda F)=\max_\alpha
 \bigl(\log|\lambda c_\alpha|+w_\alpha\bigr)
 =\log|\lambda|+\Phi(F).
\]
On the Archimedean locus,
\eqref{eq:hybrid-philippon-complex-homogeneity} and
\eqref{eq:hybrid-philippon-character} give the same formula.  This proves
\eqref{eq:hybrid-philippon-homogeneity}.  Finally,
\(eF_e=e'F_{e'}=ueF_{e'}\), so \(F_{e'}=u^{-1}F_e\), and homogeneity gives
\eqref{eq:hybrid-philippon-frame-change} with the stated minus sign.

We finish with the basic boundary and failure tests.  On the hybrid spectrum
of \(\mathbf Z\), take the degree-one form \(U_0+U_1+U_2\) and zero weights.
Then \(M=0\), and \eqref{eq:hybrid-philippon-boundary-estimate} gives
\(|\Phi|\leq C\rho\), so the Archimedean values converge to the Gauss value
as \(\rho\downarrow0\).  A tuple such as \((t,1)\) may specialize to
\((0,1)\) without affecting continuity, whereas allowing all coefficients
to vanish would already make the value of \(tU_0\) tend to \(-\infty\).
For a real Archimedean character, the sphere integral is that of its
complexification and is unchanged by the two conjugate representatives.
No sign condition on \(w_\alpha\) entered either
\eqref{eq:hybrid-philippon-exact-max} or
\eqref{eq:hybrid-philippon-boundary-estimate}.  In degree zero with one
coefficient, the frame calculation is simply
\(\log|u^{-1}r|=\log|r|-\log|u|\), confirming the sign once more.
\end{proof}

\setcounter{equation}{\savedWeightedPGEquation}
\endgroup

\section{Monge--Amp\`ere measures}
\label{sec:hybrid-ma}

Put $n=\dim X_t$ and $V=(L_t^n)>0$, independent of $t$.
For a continuous semipositive metric $\phi_t$, let $\varphi_t$ denote
its complex squared-norm potential, as fixed after
\eqref{eq:hybrid-weight-normalization}, and write
\[
          \mu_t:=V^{-1}(\ddc\varphi_t)^n.
\tag{66.22}\label{eq:hybrid-complex-ma}
\]
For a continuous semipositive non-Archimedean metric, write
\[
          \mu_0:=V^{-1}\operatorname{MA}^{\NA}(\phi^{\NA}).
\tag{66.23}\label{eq:hybrid-na-ma}
\]
Both are probability measures.

\begin{theorem}\label{thm:hybrid-ma-continuity}
Let $\mathcal X\to\Delta$ be a projective degeneration of $n$-dimensional complex projective manifolds, polarized by a relatively ample line bundle $\mathcal L$. Let $\Phi$ be a continuous semipositive metric on the hybrid line bundle. Then the fibrewise measures
\[
        \mu_t=V^{-1}(\ddc\varphi_t)^n
        \quad(t\neq0),
        \qquad
        \mu_0=V^{-1}\operatorname{MA}^{\NA}(\phi^{\NA})
\tag{66.24}\label{eq:hybrid-ma-family}
\]
are weakly continuous on $\mathcal X^{\mathrm{hyb}}$. Equivalently, for every continuous $F:\mathcal X^{\mathrm{hyb}}\to\mathbf R$,
\[
              \lim_{t\to0}\int_{X_t}F\,\rd \mu_t
                    =\int_{X_K^{\mathrm{an}}}F\,\rd \mu_0.
\tag{66.25}\label{eq:hybrid-ma-test}
\]
\end{theorem}

\begin{proof}
This is \cite[Theorem~4.11]{PS23Hybrid}, where the result is obtained by
translating Favre's projective continuity theorem into the hybrid language.
We give the complete limiting argument and isolate the one foundational
model calculation imported through that theorem.

\smallskip
\noindent\emph{Step 1: model metrics and model tests.}
For a pure hybrid Fubini--Study metric and a model test, the precise
imported interface is \cite[Theorem~3.4]{Favre20Hybrid}: the integrals on
nonzero fibres converge to the divisorial intersection integral on the
central fibre. Its proof extends the defining sections to an SNC model,
uses continuous-potential Bedford--Taylor products there, and bounds the
mass near the exceptional and intersection loci. Outside those loci the
scaled test converges to its constant divisorial value. The exceptional
mass tends to zero, and the restriction of the curvature product to each
component has its intersection degree. Complex measures on nonzero
fibres are not themselves finite intersection cycles. The dictionary
between Favre's admissible data and pure hybrid Fubini--Study metrics is
\cite[Lemma~4.10 and Theorem~4.11]{PS23Hybrid}. We import this model limit
with its stated projective hypotheses.

\smallskip
\noindent\emph{Step 2: a fibre-independent error estimate.}
Choose pure hybrid Fubini--Study metrics $\Phi^{(k)}$ with
\[
            \delta_k:=
            \|\Phi^{(k)}-\Phi\|_{\mathcal X^{\mathrm{hyb}}}
            \longrightarrow0,
\tag{66.25a}\label{eq:hybrid-ma-uniform-metric-error}
\]
and differences of hybrid Fubini--Study functions $F^{(j)}$ converging
uniformly to $F$, using \cite[Theorem~2.16]{PS23Hybrid}.
Near the central fibre these are the model tests of Step~1, up to an
$O(\lambda(t))$ error; their metric entries can be taken in a fixed
auxiliary ample class. This fixes the degrees used by the estimate.  The
Chern--Levine--Nirenberg estimate used in
\cite[Lemma~4.10 and Theorem~4.11]{PS23Hybrid}, applied after polarization,
gives for each fixed $j$ a constant $C_j$, independent of the fibre,
such that
\[
 \left|
  \int F^{(j)}\,\rd\mu_t-
  \int F^{(j)}\,\rd\mu_t^{(k)}
 \right|
 \leq C_j\delta_k,
 \qquad t\in\overline\Delta_r.
\tag{66.25b}\label{eq:hybrid-ma-uniform-cln-error}
\]
Here $\mu_t^{(k)}$ is the normalized Monge--Amp\`ere measure of
$\Phi^{(k)}|_{X_t}$, with the same definition at $t=0$.  The point of
\eqref{eq:hybrid-ma-uniform-cln-error} is that $k$ can be chosen before
$t$ tends to zero.  A merely fibrewise continuity theorem would not give
this order of quantifiers.

All the measures have mass one, so the test-function error satisfies
\[
 \left|
  \int(F-F^{(j)})\,\rd\mu_t
 \right|
 \leq\|F-F^{(j)}\|_\infty
\tag{66.25c}\label{eq:hybrid-ma-uniform-test-error}
\]
uniformly for $t$, and the same estimate holds for $\mu_t^{(k)}$.

\smallskip
\noindent\emph{Step 3: the order of the limits.}
Fix $\varepsilon>0$.  First choose $j$ so that the right-hand side of
\eqref{eq:hybrid-ma-uniform-test-error} is less than
$\varepsilon$.  With this $j$ fixed, choose $k$ so that
$C_j\delta_k<\varepsilon$.  For this fixed model pair
$(F^{(j)},\Phi^{(k)})$, Step~1 supplies $t_0$ such that the difference
between its nonzero-fibre and central integrals is less than
$\varepsilon$ whenever $0<|t|<t_0$.  Insert the two approximating
integrals between the two sides of \eqref{eq:hybrid-ma-test}.  The triangle
inequality and
\eqref{eq:hybrid-ma-uniform-cln-error}--
\eqref{eq:hybrid-ma-uniform-test-error} bound the resulting five terms by
a fixed multiple of $\varepsilon$.  Since $\varepsilon$ was arbitrary,
\eqref{eq:hybrid-ma-test} follows.

Thus the model limit is taken only after the two approximants have been
fixed; the uniform estimates then permit removal of those approximants.
This is the required limit exchange.
\end{proof}

\begin{remark}
\label{fire:hybrid-general-ma}
\Cref{thm:hybrid-ma-continuity} assumes a projective degeneration and a continuous hybrid psh metric. The corresponding continuity statement for an arbitrary flat projective family over an arbitrary integral Banach ring is posed in \cite[Question~4.9]{PS23Hybrid}.  Its full answer, with the original quantifiers and without extra reducedness or pure-dimensionality assumptions, is \cref{thm:p171-full-assembly} below.  The earlier theorem \cref{thm:hybrid-ma-continuity} remains useful because it identifies the geometric specialization mechanism in the classical hybrid-disc setting.
\end{remark}

\begin{corollary}\label{cor:hybrid-model-ma}
If $\Phi$ is induced by a relatively nef model $(\mathcal X',\mathcal L')$, then $\operatorname{MA}^{\NA}(\phi^{\NA})$ is supported on the divisorial points associated to irreducible components of $\mathcal X'_0$, and
\[
 \operatorname{MA}^{\NA}(\phi^{\NA})
   =\sum_E b_E\bigl(c_1(\mathcal L'|_E)^n\bigr)\delta_{v_E}.
\tag{66.26}\label{eq:hybrid-model-ma}
\]
The sum has total mass $V$.
\end{corollary}

\begin{proof}
This is the definition of the Chambert--Loir measure of a semipositive
model metric. The total mass is $(L_K^n)=V$ by flatness and the
projection formula. This proves the stated support and coefficient
formula directly; \Cref{thm:hybrid-ma-continuity} additionally
identifies it as the weak limit of the complex fibre measures.
\end{proof}

\subsection{Continuity over integral Banach bases}
\label{sec:hybrid-banach-family-ma}

We now answer \cite[Question~4.9]{PS23Hybrid} over an integral Banach
base; \Cref{thm:p171-full-assembly} and every fragment supporting it assume
integrality, so the answer is not asserted for an arbitrary Banach ring.
The proof is assembled in dependency order because three normalizations
must remain visible: the positive orientation of the intersection line,
the power of an Archimedean absolute value, and the generic lengths in the
top-dimensional fundamental cycle.  The algebraic construction uses
Ducrot's intersection line and relative resultant
\cite{Ducrot97CubeIntersection}; the fieldwise comparison uses the local
intersection theory of Chen--Moriwaki
\cite{ChenMoriwaki21AdelicIntersection} and the metrized intersection
formalism of \cite{BE21}.  The argument below is a book-level synthesis of
these inputs and does not attribute the resulting Banach-family theorem to
any one of them.

The following proof blocks use ordinary norm weights and the source
curvature operator. In this block, explicitly,
\[
 \ddc_\infty=\frac{i}{\pi}\partial\bar\partial=2\ddc,
 \qquad \ddc_\infty\psi=\ddc(2\psi).
\tag{66.26a}\label{eq:hybrid-banach-curvature-conversion}
\]
A powered Archimedean weight is divided by its power $\rho$ before this
curvature is taken. The block's mixed Monge--Amp\`ere measures have raw
intersection mass, with generic lengths retained; they are not the
probability measures $\mu_t$ in \eqref{eq:hybrid-ma-family}. These
conventions preserve both the fieldwise intersection formula and the
uniform estimate as $\rho\downarrow0$.

\begingroup
\providecommand{\ph}{\operatorname{ph}}
\providecommand{\cR}{\mathcal R}
\providecommand{\R}{\mathbf R}
\providecommand{\C}{\mathbf C}
\providecommand{\eps}{\varepsilon}
\providecommand{\norm}[1]{\lVert#1\rVert}
\providecommand{\abs}[1]{\lvert#1\rvert}
\edef\savedPseventeenEquation{\number\value{equation}}
\renewcommand{\theequation}{66.P17.1.\arabic{equation}}
\renewcommand{\theHequation}{66.P17.1.\arabic{equation}}
\setcounter{equation}{0}

\subsubsection{Relative resultants and intersection lines}
\label{sec:hybrid-relative-resultants}
\begin{theorem}
\label{thm:p171-revised2-complete-resultant}
Let \(S\) be a locally Noetherian integral scheme and let
\(f:X\to S\) be flat and projective with nonempty fibres of dimension
\(n\).  Let \(C_0,\ldots,C_n\) be sufficiently positive in the precise
sense of Ducrot \cite[Section~2.3]{Ducrot97CubeIntersection}: each \(C_i\)
is relatively very ample and
\(R^qf_*C_i=0\) for \(q>0\).  Put \(V_i=f_*C_i\).  Let
\(P_i=\operatorname{Proj}_S\operatorname{Sym}(V_i^\vee)\) be the projective
bundle of invertible submodules of \(V_i\), so that
\(\mathcal O_{P_i}(-1)\hookrightarrow V_i\otimes\mathcal O_{P_i}\) and
\((\pi_i)_*\mathcal O_{P_i}(d)=\operatorname{Sym}^d(V_i^\vee)\).
Put \(P=\prod_SP_i\).  This is the same coefficient-space convention as
\(\mathbf P(V_i^\vee)\) in Chen--Moriwaki
\cite[Section~1.6]{ChenMoriwaki21AdelicIntersection}.  Let

\[
 \mathscr R=I_{X/S}(C_0,\ldots,C_n)\otimes
 \bigotimes_{i=0}^n\mathcal O_{P_i}(\delta_i),
 \qquad
 \delta_i=(C_0\cdots\widehat C_i\cdots C_n)_{X/S}.
\]

Under Ducrot's canonical isomorphism
\cite[Lemma~6.3.1]{Ducrot97CubeIntersection}, let
\(r\in\Gamma(P,\mathscr R)\) be the section of
\cite[Definition~6.3.4]{Ducrot97CubeIntersection}.  For
every morphism \(\operatorname{Spec}k\to S\) from a field, the pulled-back
section \(r_k\) is the positive-intersection-line lift of the
Chen--Moriwaki resultant of the top-dimensional fundamental cycle
\cite[Definitions~1.6.3 and~1.6.7]{ChenMoriwaki21AdelicIntersection}

\[
 [X_k]_n=\sum_Z
 \ell(\mathcal O_{X_k,\eta_Z})[Z].
\]

For an integral \(n\)-dimensional closed subscheme \(Z\subset X_k\), let
\(\Gamma_k([Z])\) be the fundamental cycle of the product, over \(Z\), of
the \(n+1\) hyperplane-incidence projective bundles.  Extend \(\Gamma_k\)
linearly to \(n\)-cycles, and let \(q_k\) denote projection to \(P_k\).
Equivalently,

\begin{equation}
 \operatorname{div}(r_k)
 =q_{k*}\Gamma_k([X_k]_n)
 =\sum_Z\ell(\mathcal O_{X_k,\eta_Z})
   q_{k*}\Gamma_k([Z]).
 \label{eq:p171-revised2-divisor-cycle}
\end{equation}

The equality includes all residue-field degrees and generic lengths.
After any field extension on which there is a tuple
\(t=(t_0,\ldots,t_n)\) such that the first \(n\) members form a regular
sequence on the entire scheme \(X_k\) and the total intersection is empty,
successive restriction and the last finite norm give
\(e_t\in I_{X_k/k}(C_{0,k},\ldots,C_{n,k})\setminus\{0\}\), and

\begin{equation}
 r_k=e_tR_{X_k}^{\,t},\qquad R_{X_k}^{\,t}(t)=1.
 \label{eq:p171-revised2-positive-normalization}
\end{equation}

Thus the line is the positive intersection line.  In relative dimension
zero it is the finite-flat norm line, not its dual.  If \(n=0\) and
\(\operatorname{rank}_kV_{0,k}=1\), the complete closed immersion is
\(X_k\hookrightarrow\mathbf P(V_{0,k}^\vee)=\mathbf P^0_k\), hence
\(X_k=\operatorname{Spec}k\).  Writing an affine coefficient as
\(a_0v\) for a nonzero basis vector \(v\in V_{0,k}\), its coefficient
\(a_0\) is nonzero at the unique projective tuple, the affine resultant is
the literal linear coordinate \(a_0\), its divisor on \(\mathbf P^0_k\)
is empty, and the normalized positive norm is the unit \(1\).
\end{theorem}

\begin{theorem}
\label{thm:p171-revised2-one-stage-bridge}
Retain the preceding notation.  Let \(E_i\) be finite locally free and let
\(a_i:E_i\to V_i\) be maps such that
\(f^*E_i\twoheadrightarrow C_i\) and the induced maps
\(X\to\mathbf P_S(E_i)\) are closed immersions.  Apply

\[
 \operatorname{Sym}^{\delta_i}(a_i^\vee):
 \operatorname{Sym}^{\delta_i}(V_i^\vee)\longrightarrow
 \operatorname{Sym}^{\delta_i}(E_i^\vee)
\]

to the tensor corresponding to \(r\).  On every field fibre, the resulting
line-valued polynomial is the Chen--Moriwaki resultant
\cite[Definition~1.6.3]{ChenMoriwaki21AdelicIntersection} for the induced
closed immersions and for \([X_k]_n\), with the normalization
\eqref{eq:p171-revised2-positive-normalization}.  The maps \(a_i\) may have
kernels; this is affine polynomial substitution and asserts no morphism of
coefficient projective bundles.

More generally, let \(A\) be an integral ring and let the finitely presented
data \(X/A,C_i,E_i,a_i\) be obtained from one chosen Noetherian integral
model \(A_0\subset A\) satisfying the preceding hypotheses.  Pull the
substituted \(A_0\)-tensor to \(A\).  For every factorization
\(A_0\to A\to k\) through a field, first apply
\cref{thm:p171-revised2-complete-resultant} to the complete-system
section over the \(A_0\)-fibre \(k\), and only then perform the affine
substitution over \(k\).  The resulting fibre is the positive line-valued
fundamental-cycle resultant just described.  No inherited norm open and no
projective map attached to a possibly noninjective \(a_i\) is asserted or
used.  The intersection line itself commutes with arbitrary base change
through its perfect determinant-of-cohomology definition.

This construction deliberately makes no assertion that Ducrot's
construction \cite[Section~6.3]{Ducrot97CubeIntersection} is
independent of the chosen Noetherian model, or that
it commutes with arbitrary non-Noetherian base change.  Fibre correctness,
which is the only assertion used in the analytic comparison below, holds
for the tensor
pulled from the one fixed stage.

If \(n=0\) and \(\operatorname{rank}_kV_{0,k}=1\), the surjective
evaluation forces \((a_0)_k:E_{0,k}\to V_{0,k}\) to be nonzero.  The
substituted affine resultant is the nonzero linear form obtained by pulling
back the coordinate \(a_0\).  Its projective zero divisor is
\(\mathbf P(\ker((a_0)_k))\), which is empty when
\(\operatorname{rank}_kE_{0,k}=1\); on the complete rank-one coefficient
space the projective divisor is always empty and the normalized norm is a
unit.
\end{theorem}

\begin{corollary}
\label{cor:p171-revised2-veronese-resultants}
Let \(X/A\) be flat and projective of finite presentation over an integral
ring, with nonempty \(n\)-dimensional fibres.  Suppose that
\(s_{ij}\in H^0(X,B_i)\), \(0\leq j\leq N_i\), define closed immersions and
that the \(B_i\) are ample.  After choosing one Noetherian integral model,
there is \(p_0\) such that for every \(p\geq p_0\), Ducrot's complete
resultant \cite[Section~6.3]{Ducrot97CubeIntersection} for \(B_i^p\),
followed by the maps

\[
 E_{i,p}:=\operatorname{Sym}^p(A^{N_i+1})\longrightarrow H^0(X,B_i^p),
 \qquad e_{i,\beta}\longmapsto s_i^\beta,
\]

gives a tensor

\[
 \mathcal R_p\in I_{X/A}(B_0^p,\ldots,B_n^p)\otimes
 \bigotimes_i\operatorname{Sym}^{\delta_{i,p}}(E_{i,p}^\vee),
 \qquad
 \delta_{i,p}=p^n(B_0\cdots\widehat B_i\cdots B_n).
\]

Every field fibre of \(\mathcal R_p\) is the positive line-valued
fundamental-cycle resultant for the synchronized Veronese closed
immersions.  All generic multiplicities are retained.
\end{corollary}

% No Cartier-divisor induction occurs in this proof.

\begin{proof}
We prove \cref{thm:p171-revised2-complete-resultant}.  First isolate
exactly what is imported from Ducrot.  Ducrot
\cite[Section~2.3.3]{Ducrot97CubeIntersection} shows that each
\(V_i=f_*C_i\) is finite locally free and has the asserted field base
change.  Ducrot \cite[Section~6.3]{Ducrot97CubeIntersection} constructs the
universal sections on \(X\times_SP\).
\cite[Lemma~6.3.1]{Ducrot97CubeIntersection} identifies their intersection
line with \(\mathscr R\), with the positive exponent \(+\delta_i\), and
\cite[Definition~6.3.4]{Ducrot97CubeIntersection} constructs \(r\).  On
each regular-sequence open \(U_\varphi\),
\cite[formula~(17), Corollary~6.1.2, and Section~6.2]{Ducrot97CubeIntersection}
identifies \(r\) with the finite-flat norm of the last universal section.
\cite[Lemma~6.3.2]{Ducrot97CubeIntersection} says that the complement of
\(U=\bigcup_\varphi U_\varphi\) has fibre-local depth at least two.  The
finite-flat norm there is the determinant norm, whose formation commutes
with arbitrary base change
\cite[Lemma~31.18.6, Tag~0BD2]{stacks-project}.  We use no assertion that
Ducrot's construction \cite[Section~6.3]{Ducrot97CubeIntersection}
commutes with arbitrary base change.

Fix a field-valued point \(\operatorname{Spec}k\to S\).  Pull the stage
section and the opens to \(P_k\), and denote them by \(r_k\) and
\(U_k^0\).  Every codimension-one point \(\eta\in P_k\) belongs to
\(U_k^0\).  Indeed, the complement is the field extension of the
corresponding fibre of \(P\setminus U\).  Ducrot's depth-two assertion
\cite[Lemma~6.3.2]{Ducrot97CubeIntersection},
together with regularity of the product of projective spaces, says that
every component of that closed set has codimension at least two; this
property is preserved by a field extension.  Here \(P_k\) is smooth over
\(k\), hence regular by
\cite[Lemma~33.25.3, Tag~056S]{stacks-project}.

Fix such an \(\eta\), choose \(\varphi\) with
\(\eta\in U_{\varphi,k}^0\), and choose an integral affine open

\[
 \eta\in T=\operatorname{Spec}A\subset U_{\varphi,k}^0.
\]

Let \(\mathfrak p\subset A\) be the height-one prime representing
\(\eta\), let \(R=A_{\mathfrak p}=\mathcal O_{P_k,\eta}\), and put
\(d=\dim P_k\).  The first \(n\) universal sections on
\(X_k\times_kT\), in the order \(\varphi\), form a regular sequence on
the entire scheme and cut out a finite flat scheme \(q_T:W_T\to T\).
Write \(L_T\) for the restriction to \(W_T\) of the last universal
invertible sheaf.  Its section induces a map between finite locally free
\(\mathcal O_T\)-modules of the same rank,

\[
 h_T:q_{T*}\mathcal O_{W_T}\longrightarrow q_{T*}L_T.
\]

It is an isomorphism at the generic point of the integral scheme \(T\),
because the generic universal tuple has empty common zero.  Hence
\(h_T\) is injective.  Since \(q_T\) is finite, the original map
\(\mathcal O_{W_T}\to L_T\) is injective as well.  If
\(\mathcal I_T\subset W_T\) is its zero scheme, exactness of finite
pushforward gives

\begin{equation}
 \mathcal C_T:=\operatorname{coker}(h_T)
 =q_{T*}(L_T|_{\mathcal I_T}).
 \label{eq:p171-revised2-cokernel-sheaf}
\end{equation}

Base change to \(R\).  The finite \(R\)-algebra
\(B=\Gamma(W_T\times_T\operatorname{Spec}R,\mathcal O)\) is finite free.
It is semilocal, so the restriction of \(L_T\) is free.  After choosing a
generator, the last section is an element \(a\in B\), and the norm section
is represented by \(\det(m_a:B\to B)\).  The determinant is nonzero, and
the same torsion-free argument makes \(m_a\) injective.  Smith normal form
over the DVR gives (compare
\cite[\emph{Lengths and determinants}, Tag~02QA]{stacks-project})

\begin{equation}
 \operatorname{ord}_\eta(r_k)
 =\operatorname{ord}_R\det(m_a)
 =\operatorname{length}_R(B/aB)
 =\operatorname{length}_R((\mathcal C_T)_\mathfrak p).
 \label{eq:p171-revised2-dvr-length}
\end{equation}

We now calculate the incidence cycle on the affine neighbourhood \(T\),
before localizing its coefficient.  This order is essential: the morphism
\(\operatorname{Spec}R\to P_k\) is not used as a flat pullback in
\cite[Tag~02RA]{stacks-project}.  The projection
\(p_T:X_k\times_kT\to X_k\) is a flat morphism between schemes locally of
finite type over \(k\), so
\cite[Definition~42.14.1 and Lemma~42.14.4, Tag~02RA]{stacks-project}
apply to \(p_T\).  They identify the starting
top-dimensional coherent-sheaf cycle with
\(p_T^*[X_k]_n\), retaining the generic lengths.

Apply the Stacks Project's effective-Cartier identity for coherent-sheaf
cycles \cite[Lemma~42.29.5, Tag~02TB]{stacks-project}, successively on
\(X_k\times_kT\) to the first \(n\) universal sections.  At every stage
the section is regular on the entire preceding quotient.  The
coefficient-by-coefficient computation in the proof of
\cite[Tag~02TB]{stacks-project} therefore yields the associated-cycle
identity

\begin{equation}
 [\mathcal O_{W_T}]_d
 =c_1(\mathcal L_{\varphi(n-1)})\cdots
  c_1(\mathcal L_{\varphi(0)})\cap p_T^*[X_k]_n.
 \label{eq:p171-revised2-first-n-cuts}
\end{equation}

The injectivity proved above permits one last use of the same identity:

\begin{equation}
 [L_T|_{\mathcal I_T}]_{d-1}
 =c_1(L_T)\cap[\mathcal O_{W_T}]_d.
 \label{eq:p171-revised2-last-cut}
\end{equation}

For an integral top-dimensional component \(Z\subset X_k\), let
\(V_i\otimes\mathcal O_Z\twoheadrightarrow C_i|_Z\) have kernel
\(Q_i|_Z\).  The full incidence over \(Z\) is the product over \(Z\) of
the projective bundles of lines in the \(Q_i|_Z\).  Its restriction over
\(T\) is exactly the common zero scheme used in
\eqref{eq:p171-revised2-first-n-cuts}--\eqref{eq:p171-revised2-last-cut}.
Thus, by linearity in
\([X_k]_n=\sum_Z\ell(\mathcal O_{X_k,\eta_Z})[Z]\), those two identities
say

\begin{equation}
 [L_T|_{\mathcal I_T}]_{d-1}
 =j_{\mathcal I}^*\Gamma_k([X_k]_n),
 \label{eq:p171-revised2-incidence-restriction}
\end{equation}

where \(j:T\hookrightarrow P_k\) and \(j_{\mathcal I}\) is the induced
open immersion on the incidence scheme.  Here
\cite[Tag~02RA]{stacks-project} is used
only for the finite-type flat projection \(p_T\) and for open immersions,
never for the local-spectrum morphism
\(\operatorname{Spec}R\to P_k\).

Proper pushforward of coherent-sheaf cycles
\cite[Lemma~42.12.4, Tag~02R3]{stacks-project} and
\eqref{eq:p171-revised2-cokernel-sheaf} give

\begin{equation}
 [\mathcal C_T]_{d-1}
 =q_{T*}[L_T|_{\mathcal I_T}]_{d-1}
 =j^*q_{k*}\Gamma_k([X_k]_n).
 \label{eq:p171-revised2-open-pushforward}
\end{equation}

The last equality is also elementary: proper pushforward commutes with
restriction to an open Cartesian square, as is checked on each integral
cycle generator from
\cite[Definitions~42.12.1 and~42.14.1]{stacks-project}; the relevant
function-field degree is unchanged after restriction.  Finally, by the
definition of a coherent-sheaf cycle \cite[Tag~02QV]{stacks-project}, the
coefficient
of the height-one point \(\eta\) in the left side of
\eqref{eq:p171-revised2-open-pushforward} is
\(\operatorname{length}_{A_\mathfrak p}((\mathcal C_T)_\mathfrak p)\).
Combining this with \eqref{eq:p171-revised2-dvr-length} identifies the
coefficient at \(\eta\) of the two sides of
\eqref{eq:p171-revised2-divisor-cycle}.  This includes all residue-field
degrees through \cite[Tag~02R3]{stacks-project} and all generic component
lengths through \cite[Tag~02QV]{stacks-project}.

Since every codimension-one point was arbitrary and \(P_k\) is regular, a
section of a line bundle has divisor determined by these DVR orders.  This
proves \eqref{eq:p171-revised2-divisor-cycle}.  The cycle is precisely the
one in Chen--Moriwaki
\cite[Definition~1.6.3]{ChenMoriwaki21AdelicIntersection}; lower-dimensional
and embedded
components do not occur in \([X_k]_n\), while every top-dimensional generic
nilpotent thickness is the displayed local length.

It remains to determine the scalar and orientation.  Pass faithfully flatly
to a purely transcendental extension \(K/k\), hence to an infinite field.
Set \(Y_0=X_K\).  For \(0\leq j<n\), recompute the full finite set
\(\operatorname{Ass}(Y_j)\).  For each associated point, sections in the
complete basepoint-free system of \(C_j\) that vanish there form a proper
linear subspace.  A finite union of proper linear subspaces cannot exhaust
the section space over \(K\); choose \(t_j\) outside their union and put
\(Y_{j+1}=V(t_j)\subset Y_j\).  Multiplication by \(t_j\) is injective on
all of \(\mathcal O_{Y_j}\), so \(Y_{j+1}\) is an effective Cartier
divisor on the entire possibly nonpure and nonreduced scheme \(Y_j\).
Embedded and lower-dimensional components are included in the associated
point avoidance: they are either missed or cut, and none is discarded.
The choice is repeated on the new quotient rather than inferred from an
earlier associated-prime calculation.

After \(n\) steps, \(Y_n\) has dimension at most zero.  Indeed, on every
surviving positive-dimensional projective component the restriction of a
very ample line bundle cannot have a nowhere-vanishing section, so a regular
section cuts that component in dimension.  Choose \(t_n\) nonvanishing at
every point of the finite support of \(Y_n\).  The total intersection is
empty.

Let \(t=(t_0,\ldots,t_n)\).  Ducrot's successive positive restriction maps
and final norm \cite[Sections~6.1--6.2]{Ducrot97CubeIntersection} define
\(e_t\in I_{X_K/K}(C_0,\ldots,C_n)\setminus\{0\}\).  Write
\(r_K=e_tF_t\).  By construction \(F_t(t)=1\).  Let \(R_t\) be the
Chen--Moriwaki resultant normalized by the literal equality \(R_t(t)=1\),
as in \cite[Definition~1.6.7]{ChenMoriwaki21AdelicIntersection}.  The
divisor identity gives
\(\operatorname{div}(F_t)=\operatorname{div}(R_t)\).  Hence
\(F_t/R_t\) is a global unit on the integral projective coefficient product
\(P_K\).  Since \(\Gamma(P_K,\mathcal O_{P_K})=K\), that ratio is a scalar
in \(K^\times\); evaluation at \(t\) makes the scalar literally \(1\).
Thus \(F_t=R_t\), proving
\eqref{eq:p171-revised2-positive-normalization} over \(K\).  Equality
descends to \(k\) because scalar extension of a vector space is injective.

It remains to separate the rank-one, zero-dimensional endpoint.  If
\(n=0\), Ducrot \cite[Proposition~6.1.1(3)]{Ducrot97CubeIntersection}
identifies the intersection line with
the positive finite-flat norm, and
\eqref{eq:p171-revised2-dvr-length} is the determinant of multiplication on
the original finite algebra, as in Chen--Moriwaki
\cite[Example~1.6.4]{ChenMoriwaki21AdelicIntersection}.  If in
addition \(\dim_kV_{0,k}=1\), very ampleness makes the complete morphism a
closed immersion
\(X_k\hookrightarrow\mathbf P(V_{0,k}^\vee)=\mathbf P^0_k\), so it forces
\(X_k=\operatorname{Spec}k\).  Choose a nonzero basis section \(v\).  On
the affine coefficient line the universal section is \(a_0v\): for a
nonzero tuple \(a_0\ne0\), and the determinant norm polynomial is the
single linear coordinate \(a_0\).  On \(\mathbf P^0_k\), this coordinate
is nowhere vanishing as a section of \(\mathcal O(1)\), so the incidence
and projective divisor are empty.  Framing by the value at the unique
projective tuple makes the normalized positive norm the unit \(1\), not a
section of a dual norm line.
\end{proof}

\begin{proof}
We prove \cref{thm:p171-revised2-one-stage-bridge}.  The maps on
symmetric powers are ordinary contravariant substitution of linear forms,
so kernels of the \(a_i\) cause no algebraic problem.  We first prove the
comparison after passage to a field \(k\), without constructing a
projective map between coefficient spaces.

Let \(Z\) be an integral top-dimensional component of \(X_k\).  The
complete resultant polynomial, after affine substitution by the \(a_i\),
vanishes at \(e=(e_0,\ldots,e_n)\) exactly when the evaluated sections
\(a_i(e_i)\) have a common zero on \(Z\).  This is the degeneracy locus for
the closed immersions defined by the \(E_i\).  Chen--Moriwaki
\cite[Proposition~1.6.2 and Definition~1.6.3]{ChenMoriwaki21AdelicIntersection}
say that this locus is an irreducible
hypersurface.  The substituted polynomial is nonzero because \(n+1\)
general members of these basepoint-free ample systems have empty common
intersection.  It remains homogeneous in block \(i\) of the unchanged
positive degree

\[
 \delta_i(Z)=(C_0\cdots\widehat C_i\cdots C_n)_Z.
\]

Equality of zero sets gives a positive power of the irreducible equation;
comparison in any block gives exponent one.  Hence the substituted complete
polynomial is a scalar times the \(E_i\)-resultant of \(Z\).

For the fundamental cycle,
\cref{thm:p171-revised2-complete-resultant} and Chen--Moriwaki
\cite[Definition~1.6.3]{ChenMoriwaki21AdelicIntersection} factor the
complete polynomial as the product over \(Z\),
with exponent \(\ell(\mathcal O_{X_k,\eta_Z})\).  Affine substitution
distributes over this product, so the same generic lengths occur for the
\(E_i\)-resultants.  After a purely transcendental extension if needed,
choose a tuple from the \(E_i\) by successively avoiding all associated
points of every new quotient, as in the preceding proof.  The first \(n\)
members then form a regular sequence on all of \(X_k\), and the total
intersection is empty.  Evaluation of the substituted line-valued tensor
is the same positive norm element \(e_t\), so its scalar polynomial takes
the literal value \(1\).  Chen--Moriwaki's normalization
\cite[Definition~1.6.7]{ChenMoriwaki21AdelicIntersection} fixes the remaining
scalar.  Equality descends from the faithfully flat field extension.

Now let \(A_0\subset A\) be the one chosen stage and let
\(A_0\to A\to k\) be any factorization through a field.  Associativity of
tensor product says that the \(k\)-fibre of the substituted \(A_0\)-tensor
pulled to \(A\) is obtained as follows: first take the \(k\)-fibre of the
complete \(A_0\)-section, and then apply the linear substitutions
\((a_i)_k\).  Apply
\cref{thm:p171-revised2-complete-resultant} directly to the field
map \(\operatorname{Spec}k\to\operatorname{Spec}A_0\).  It identifies the
first object with the complete fundamental-cycle resultant over \(k\).
Apply the affine zero-set and block-degree argument just proved over \(k\)
to identify its substitution with the finite-system Chen--Moriwaki
resultant \cite[Definition~1.6.3]{ChenMoriwaki21AdelicIntersection}.  This
argument does not pull an \(A_0\)-regular open through
\(A\to k\), and it does not use a projective map associated to an
\(a_i\) having a kernel.  It proves fibre correctness without proving model
independence.  The intersection line, unlike Ducrot's construction
\cite[Section~6.3]{Ducrot97CubeIntersection}, is
the alternating finite difference of determinants of perfect complexes
\(Rf_*\bigotimes_{i\in I}C_i\); proper perfect base change
\cite[Tag~0B91]{stacks-project} gives its arbitrary base change.

For completeness, consider \(n=0\) and \(\dim_kV_{0,k}=1\).  The first
proof gives \(X_k=\operatorname{Spec}k\), and the surjective evaluation
makes \((a_0)_k:E_{0,k}\to V_{0,k}\) a nonzero linear map.  If \(z_0\)
is the affine coordinate on \(V_{0,k}\), the complete affine resultant is
\(z_0\), and its substitution is the nonzero linear form
\(z_0\circ(a_0)_k\).  Its projective zero divisor is exactly
\(\mathbf P(\ker((a_0)_k))\), the incidence divisor for the embedded point;
it is empty when \(E_{0,k}\) also has rank one.  Choosing
\(e\in E_{0,k}\) with \((a_0)_k(e)=v\) makes the normalized scalar value
literally \(1\), in agreement with the positive finite norm.  This is
separate from the complete rank-one projective coefficient space, whose
divisor is always empty.
\end{proof}

\begin{proof}
We prove \cref{cor:p171-revised2-veronese-resultants}.  Descend the
finite presentation, the line bundles, sections, evaluation maps, and the
original closed immersions to one finitely generated
\(\mathbf Z\)-subalgebra \(A_0\subset A\), enlarging it until the model is
flat.  The needed limit statements are
\cite[Tags~01ZM, 01ZR, 0B8W, and~05LY]{stacks-project}.  Since \(A_0\subset A\)
and \(A\) is integral, \(A_0\) is a Noetherian domain.  Proper flat finite
presentation and \cite[Tag~0D4R]{stacks-project} give constant fibre
dimension \(n\).

Relative Serre vanishing gives \(p_0\) such that \(B_{i,0}^p\gg0\) for all
\(i\) and \(p\geq p_0\).  The all-monomial system is the composite of the
original closed immersion with the degree-\(p\) Veronese closed immersion.
Apply \cref{thm:p171-revised2-one-stage-bridge}.  Ducrot
\cite[Lemma~6.3.1]{Ducrot97CubeIntersection} computes the block degree as

\[
 (B_{0,k}^p\cdots\widehat{B_{i,k}^p}\cdots B_{n,k}^p)
 =p^n(B_0\cdots\widehat B_i\cdots B_n),
\]

including the generic multiplicities of the fibre.  This proves the
corollary.
\end{proof}


\subsubsection{The power parameter and the exact fibre convention}

Let $A$ be an integral Banach ring.  For every $x\in\cM(A)$, set
\begin{equation}
 \rho(x):=\frac{\log\max\{1,\abs{2}_x\}}{\log 2}.
\label{eq:rho-base}
\end{equation}
This is a continuous function on $\cM(A)$.

\begin{lemma}
\label{lem:rho-split}
One has $\rho(x)=0$ exactly when the absolute value on the completed residue
field $\mathcal H(x)$ is non-Archimedean.  If $\rho(x)>0$, then
$0<\rho(x)\leq1$, and there is a character into $\C$, unique up to complex
conjugation, such that after completing the fraction field
\begin{equation}
 \abs{a}_x=\abs{\chi_x(a)}_\infty^{\rho(x)}.
\label{eq:character-power}
\end{equation}
Thus $\mathcal H(x)$ is $\R$ or $\C$, with the
$\rho(x)$-th power of the usual absolute value.
\end{lemma}

\begin{proof}
If $\abs{2}_x\leq1$, binary expansion first gives
$\abs{n}_x\leq1$ for all positive integers $n$; applying the binomial
formula to $(a+b)^m$ and taking $m$-th roots gives the ultrametric
inequality.  The converse follows from
$\abs{2}_x\leq\max\{\abs1_x,\abs1_x\}=1$.

If $\abs2_x>1$, the base-expansion proof of Ostrowski's theorem gives
$\abs n_x=n^{\rho(x)}$.  The ordinary triangle inequality gives
$\abs2_x\leq2$, hence $0<\rho(x)\leq1$.  The unpowered function
$v(a)=\abs a_x^{1/\rho(x)}$ is multiplicative, is usual on $\mathbf Q$,
and satisfies the triangle inequality: expanding $(a+b)^m$ yields
\[
 \abs{a+b}_x^m\leq(m+1)
 \bigl(v(a)+v(b)\bigr)^{m\rho(x)},
\]
and the power $1/(m\rho(x))$, followed by $m\to\infty$, proves the
claim.  The completion is a commutative real Banach division algebra, hence
$\R$ or $\C$ by Gelfand--Mazur.  The only isometric ambiguity in the
complex case is conjugation.  This is the character construction used in
\cref{thm:hybrid-weighted-philippon-gauss}.
\end{proof}

Fix henceforth an Archimedean point and abbreviate
$K_\rho=(K,\abs{\cdot}_\infty^\rho)$, where
$K\in\{\R,\C\}$ and $0<\rho\leq1$.  A metric weight is
additive: for a local frame $e$,
\begin{equation}
 \phi_\rho(e):=-\log\norm e_\rho,
 \qquad
 \phi_\rho(ue)=\phi_\rho(e)-\log\abs u_\rho.
\label{eq:metric-transition}
\end{equation}
Its unpowered metric is
\begin{equation}
 \overline\phi(e):=\rho^{-1}\phi_\rho(e),
 \qquad
 \norm e_{\overline\phi,\infty}:=\norm e_{\phi_\rho,\rho}^{1/\rho}.
\label{eq:metric-unpower}
\end{equation}
Division of~\eqref{eq:metric-transition} by $\rho$ gives the usual metric
transition law.

\begin{lemma}
\label{lem:PS-transport}
Under the canonical homeomorphism from a usual complex or real
analytification to the $K_\rho$-analytification, the restricted fibre metric
of \cite[Proposition~3.2]{PS23Hybrid} is transported to the usual complex
metric by \eqref{eq:metric-unpower}.  The fibrewise Monge--Amp\`ere measure
of \cite[Definition~4.8]{PS23Hybrid} is therefore the Boucksom--Eriksson
Archimedean measure of the unpowered metrics
\cite[formula~(8.1) and Proposition~8.3]{BE21}.  This is a consequence of
the stated fibre convention, not a new formal definition of a symbol
$\ddc_\rho$.
\end{lemma}

\begin{proof}
The exponent is explicit in \cite[Proposition~3.2]{PS23Hybrid}.  In the
hybrid-disc notation, put
\[
 \rho_t:=\frac{\log r}{\log\abs t}\quad(0<\abs t\leq r<1).
\]
For the canonical homeomorphism $\beta_t:X_t\to X_t^{\mathrm{hyb}}$, their
proof gives, for every local regular function $f$,
\begin{equation}
 \abs{f(\beta_t(z))}=\abs{f(z)}_\infty^{\rho_t},
 \qquad
 \norm{s(z)}_{\phi_t}
 =\norm{s(\beta_t(z))}_{\phi}^{1/\rho_t}.
\label{eq:PS32-transport}
\end{equation}
Equivalently,
$\log\norm{s}_{\phi_t}/\log_r\abs t$
extends as the hybrid logarithmic norm.  Thus the raw hybrid fibre weight is
$\rho_t$ times the ordinary complex weight, exactly as in
\eqref{eq:metric-unpower}.

For an arbitrary Archimedean $x\in\cM(A)$, the character
\eqref{eq:character-power} gives the same canonical map: on every affine
chart it sends a complex point $z$ to
$f\mapsto\abs{\chi_x(f)(z)}_\infty^{\rho(x)}$.  Formula
\eqref{eq:metric-unpower} is therefore the unique metric transport compatible
with local frames and this seminorm map; a different representing character
only conjugates the construction.

This transport also preserves the psh class used in \cite{PS23Hybrid}.  For
a tropical Fubini--Study metric,
\begin{equation}
 \rho^{-1}\frac1m\max_j
 \{\log\abs{s_j}_\rho+a_j\}
 =\frac1m\max_j
 \{\log\abs{s_j}_\infty+a_j/\rho\},
\label{eq:FS-transport}
\end{equation}
which is usual complex psh by \cite[Theorem~2.11]{PS23Hybrid}.
\cite[Proposition~2.21]{PS23Hybrid} and multiplication by the fixed scalar
$1/\rho$ transport uniform limits as well.

Finally, \cite[Definition~4.8]{PS23Hybrid} defines the family measure at $x$
by restricting the metrics to $X_{\mathcal H(x)}^{\mathrm{an}}$ and applying the
valued-field Monge--Amp\`ere operator of
\cite[Section~4.1]{PS23Hybrid}.  On an Archimedean field,
Boucksom--Eriksson first identify the field with $\R$ or $\C$ only
\emph{up to normalization of its absolute value}
\cite[Section~1.1]{BE21}; their Archimedean current is the classical
Bedford--Taylor current on the usual complex space.  Equations
\eqref{eq:PS32-transport}--\eqref{eq:FS-transport} identify the metric
supplied to that current as $\overline\phi$.

Two further source tests make this reading fail-closed.  First,
\cite[Theorem~4.11]{PS23Hybrid} writes the Archimedean members of a hybrid
family as $(\ddc\phi_t)^n$, where \cite[Proposition~3.2]{PS23Hybrid} has
already defined $\phi_t=\rho_t^{-1}\phi|_{X_t^{\mathrm{hyb}}}$.  Thus the
displayed complex operator acts on the unpowered metric.

Second, \cite[Theorem~8.14]{BE21} is stated over an arbitrary complete valued
field, and its proof uses
\begin{equation}
 \ddc\log\norm{s}_{\phi}\wedge\delta_X
 =\delta_Z-\ddc\phi\wedge\delta_X.
\label{eq:PL-source}
\end{equation}
For $K_\rho$, both
$\log\norm{s}_{\phi_\rho,\rho}$ and $\phi_\rho$ are $\rho$ times their
unpowered counterparts.  Applying the raw Euclidean operator to those
powered weights would turn the left side into
$\rho\delta_Z-\rho\ddc_\infty\overline\phi$, while the right side demanded
by~\eqref{eq:PL-source} would be
$\delta_Z-\rho\ddc_\infty\overline\phi$.  These differ for
$\rho\ne1$ and a nonzero divisor.  The source identity is restored exactly
by the transported operator
\begin{equation}
 \dc_\rho\phi_\rho:=\dc_\infty(\rho^{-1}\phi_\rho),
 \qquad
 \ddc_\rho\phi_\rho:=\ddc_\infty(\rho^{-1}\phi_\rho).
\label{eq:ddc-unpacking}
\end{equation}
Equation~\eqref{eq:ddc-unpacking} is our explicit unpacking of the source
convention, not a formula printed verbatim in either source.
\cite[Proposition~8.3(ii)]{BE21}, which fixes total mass to the
algebraic intersection number, supplies the independent mass check: the raw
Euclidean operator on $\rho\overline\phi$ would multiply an $n$-fold mass
by $\rho^n$.  Hence \cite[Definition~4.8]{PS23Hybrid} has the transported
reading asserted in the lemma.
\end{proof}

\subsubsection{Powered Archimedean normalization}

Let $X/K_\rho$ be projective of dimension $n$; nilpotents, generic
multiplicities, and several top-dimensional components are allowed.

\begin{theorem}
\label{thm:powered-normalization}
For continuous psh-regularizable metrics $\phi_{i,\rho}$, put
$\overline\phi_i=\rho^{-1}\phi_{i,\rho}$.  Under the canonical
identification of analytic spaces,
\begin{equation}
 \MA_\rho(\phi_{1,\rho},\ldots,\phi_{n,\rho})
 =\MA_\infty(\overline\phi_1,\ldots,\overline\phi_n).
\label{eq:PN1-r3}
\end{equation}
All generic lengths in the fundamental cycle occur on both sides.  If $e$
is a nonzero vector in the Deligne line $I_X(L_0,\ldots,L_n)$, then
\begin{equation}
 \lambda_\rho(e;\phi_{0,\rho},\ldots,\phi_{n,\rho})
 =\rho\lambda_\infty
 (e;\overline\phi_0,\ldots,\overline\phi_n).
\label{eq:PN2-r3}
\end{equation}
Equivalently, the recursively defined local height has exactly one factor
$\rho$ in every dimension.
\end{theorem}

\begin{proof}
We prove \eqref{eq:PN1-r3}.  \Cref{lem:PS-transport} identifies the smooth
complex fibre measure with
\[
 \ddc_\infty\overline\phi_1\wedge\cdots\wedge
 \ddc_\infty\overline\phi_n.
\]
For a general scheme, \cite[formula~(8.1)]{BE21} uses the fundamental cycle
\[
 [X]_n=\sum_Zm_Z[Z],
 \qquad
 m_Z=\operatorname{length}(\mathcal O_{X,\eta_Z}).
\]
The resulting measure is the sum of the component measures multiplied by
$m_Z$.
Applying the transport on each reduced top-dimensional component proves the
smooth statement with all generic lengths.

For continuous psh-regularizable metrics, uniformly approximate by smooth
semipositive nets.  Division by the fixed $\rho$ preserves uniform
convergence.  \cite[Theorem~8.4]{BE21}, as used in
\cite[Theorem~4.3]{PS23Hybrid}, gives the unique weakly continuous extension
and hence
\eqref{eq:PN1-r3}.  Over $\R$, both sources define the measure as the
push-forward from complexification.  Complexification commutes with the
transport and introduces no factor two.
\end{proof}

\begin{lemma}
\label{lem:zero-anchor-r3}
Let $X=\operatorname{Spec}A$ be finite and flat over $K_\rho$, and let
$\phi_{0,\rho}$ be a continuous dpsh metric on $L_0$.  Put
$\overline\phi_0=\rho^{-1}\phi_{0,\rho}$.  Under the algebraic
identification of the Deligne line with the norm line,
\begin{equation}
 \log\norm{\langle s\rangle}_{\langle\phi_{0,\rho}\rangle,\rho}
 =\rho\,
 \log\norm{\langle s\rangle}_{\langle\overline\phi_0\rangle,\infty}
\label{eq:zero-metric-r3}
\end{equation}
for every trivializing section $s$ of $L_0$.  Consequently the additive
Deligne weight and every dimension-zero local height are multiplied by
exactly $\rho$, with all residue degrees and nonreduced Artin lengths.
For $K=\R$, the same assertion uses the real-fibre push-forward convention
and has no factor two.
\end{lemma}

\begin{proof}
The algebraic input is \cite[Example~8.15]{BE21}: in dimension zero
$\langle L_0\rangle=N_{A/K}(L_0)$ and
$\langle us\rangle=N_{A/K}(u)\langle s\rangle$.
\cite[Proposition~6.1.1(3)]{Ducrot97CubeIntersection} fixes the positive
algebraic norm orientation, but
neither algebraic statement fixes the multiplicative constant of the
metric.

That constant is fixed by
\cite[Theorem~8.16, formula~(8.11)]{BE21}.  Reconstructed in the present
norm notation, the formula is
\begin{equation}
 \log\norm{\langle s\rangle}_{\langle\phi_0\rangle}
 =\int_{X^{\mathrm{an}}}\log\norm{s}_{\phi_0}\,\delta_X.
\label{eq:BE811-r3}
\end{equation}
Write $A=\prod_jA_j$, with $A_j$ local, residue field $K_j$, and Artin
length $m_j$.  \cite[Example~8.2]{BE21} gives, on the corresponding
analytic points $x_j$,
\begin{equation}
 \delta_X=\sum_j(\dim_KA_j)\delta_{x_j}
 =\sum_jm_j[K_j:K]\,\delta_{x_j}.
\label{eq:zero-cycle-r3}
\end{equation}
Under the canonical identification of the powered and unpowered analytic
spaces, this measure is unchanged, while metric unpowering gives
$\log\norm{s(x_j)}_{\phi_{0,\rho},\rho}
 =\rho\log\norm{s(x_j)}_{\overline\phi_0,\infty}$.
Applying~\eqref{eq:BE811-r3} twice and then~\eqref{eq:zero-cycle-r3}
therefore yields
\begin{align*}
 \log\norm{\langle s\rangle}_{\langle\phi_{0,\rho}\rangle,\rho}
 &=\sum_jm_j[K_j:K]\,
   \log\norm{s(x_j)}_{\phi_{0,\rho},\rho}\\
 &=\rho\sum_jm_j[K_j:K]\,
   \log\norm{s(x_j)}_{\overline\phi_0,\infty}\\
 &=\rho\,
   \log\norm{\langle s\rangle}_{\langle\overline\phi_0\rangle,\infty}.
\end{align*}
This proves~\eqref{eq:zero-metric-r3} and visibly preserves both the residue
degree $[K_j:K]$ and the Artin length $m_j$.  Frame compatibility follows
from \cite[Theorem~8.16, formula~(8.12)]{BE21}, namely
$\int\log\abs u\,\delta_X=\log\abs{N_{A/K}(u)}$; after powering, both
sides acquire the same factor $\rho$.  Finally,
\cite[Remark~8.1]{BE21} defines the real measure as the push-forward from
complexification.  Thus the
degree already present in~\eqref{eq:zero-cycle-r3} is retained and no
additional factor two appears.
\end{proof}

\medskip
\begin{proof}
We prove \eqref{eq:PN2-r3}.  Define the candidate powered Deligne weight by
\begin{equation}
 \widetilde\lambda_\rho
 (e;\phi_{0,\rho},\ldots,\phi_{n,\rho})
 :=\rho\lambda_\infty
 (e;\overline\phi_0,\ldots,\overline\phi_n).
\label{eq:PN2-candidate-r3}
\end{equation}
If $e'=ue$, then
$
\widetilde\lambda_\rho(e')=\widetilde\lambda_\rho(e)
-\rho\log\abs u_\infty
=\widetilde\lambda_\rho(e)-\log\abs u_\rho
$, so this is a metric for the powered absolute value.
At fixed $\rho$, unpowering is additive under tensor products; hence
symmetry and multilinearity of the usual Deligne metric give the same
properties for $\widetilde\lambda_\rho$.

Let $s\in H^0(X,L_0)$ be regular, with divisor $Z$.  The usual additive
restriction formula is
\begin{align*}
 \lambda_\infty(e;\overline\phi_0,\ldots,\overline\phi_n)
 ={}&\lambda_\infty(e|_Z;
       \overline\phi_1|_Z,\ldots,\overline\phi_n|_Z)\\
 &-\int_X\log\norm{s}_{\overline\phi_0,\infty}
       \,\MA_\infty(\overline\phi_1,\ldots,\overline\phi_n).
\end{align*}
Multiplication by $\rho$, the identity
$
 \log\norm{s}_{\phi_{0,\rho},\rho}
 =\rho\log\norm{s}_{\overline\phi_0,\infty}
$, and \eqref{eq:PN1-r3} turn this into the powered restriction formula.  The
dimension-zero lemma supplies the base case, and the uniqueness clause in
\cite[Theorem~8.16]{BE21} identifies the candidate with the target Deligne
metric.

The same calculation is the local-height induction: the restricted height
has one factor $\rho$, while the integral has one factor from its Green
function and none from its measure.  No new factor is accumulated at the
next dimension.
\end{proof}

\subsubsection{A fibre-uniform Chern--Levine--Nirenberg estimate}

For every $0\leq i\leq n$, let
$\phi_{i,\rho},\phi'_{i,\rho}$ be continuous psh-regularizable
semipositive metrics on $L_i$, and put
$f_\rho=\phi_{0,\rho}-\phi'_{0,\rho}$.  In particular, the slot-zero
test is a difference of two metrics in the semipositive domain of the CLN
inequality; it is not an arbitrary continuous function.  Define
\begin{equation}
 C:=
 \begin{cases}
 0,&n=0,\\[2mm]
 2\displaystyle\max_{1\leq i\leq n}
 (L_0L_1\cdots\widehat{L_i}\cdots L_n),&n\geq1.
 \end{cases}
\label{eq:CLN-C-r3}
\end{equation}

\begin{lemma}
\label{lem:CLN-r3}
Assume explicitly that both metrics in every pair
$\phi_{i,\rho},\phi'_{i,\rho}$, including $i=0$, are continuous
psh-regularizable and semipositive.  For $0<\rho\leq1$,
\begin{align}
 &\left|\int_Xf_\rho\,
 \MA_\rho(\phi_{1,\rho},\ldots,\phi_{n,\rho})
 -\int_Xf_\rho\,
 \MA_\rho(\phi'_{1,\rho},\ldots,\phi'_{n,\rho})\right|\notag\\
 &\hspace{30mm}\leq
 C\sum_{i=1}^n
 \norm{\phi_{i,\rho}-\phi'_{i,\rho}}_\infty.
\label{eq:CLN-r3}
\end{align}
The constant is independent of $\rho$.
\end{lemma}

\begin{proof}
If $n=0$, there are no metric slots to change, the two integrals coincide,
and both sides are zero by definition~\eqref{eq:CLN-C-r3}.

Assume $n\geq1$ and first suppose all metrics are smooth semipositive.
Put $\overline\phi_i=\rho^{-1}\phi_{i,\rho}$ and
$\overline\phi'_i=\rho^{-1}\phi'_{i,\rho}$ for every $0\leq i\leq n$,
and $\overline f=\overline\phi_0-\overline\phi'_0$.
The convention transport preserves semipositivity, and
\eqref{eq:PN1-r3} makes the left side equal to
\[
 \rho\left|\int_X\overline f\,
 \MA_\infty(\overline\phi_1,\ldots)
 -\int_X\overline f\,
 \MA_\infty(\overline\phi'_1,\ldots)\right|.
\]
\cite[Lemma~8.6]{BE21} applies because the slot-zero pair
$(\overline\phi_0,\overline\phi'_0)$ is also semipositive.  It gives one
inequality with constant~\eqref{eq:CLN-C-r3}.  Keep that test pair fixed and
interchange primed and unprimed metrics only in slots $1,\ldots,n$ to get
the opposite inequality.  Therefore
\begin{align*}
 \rho\left|\int\overline f\,\MA_\infty(\overline\phi_i)
 -\int\overline f\,\MA_\infty(\overline\phi'_i)\right|
 &\leq\rho C\sum_i
 \norm{\overline\phi_i-\overline\phi'_i}_\infty\\
 &=C\sum_i
 \norm{\phi_{i,\rho}-\phi'_{i,\rho}}_\infty.
\end{align*}
This is the two-sided cancellation: the test function contributes $\rho$,
and unpowering a metric difference contributes $1/\rho$.

For the stated continuous domain, uniformly approximate all
$2(n+1)$ metrics, including the slot-zero pair, by smooth semipositive
metrics.  \cite[Theorem~8.4]{BE21} gives weak convergence of the
Monge--Amp\`ere measures, while uniform convergence of the slot-zero
differences passes the integrals to the limit.  The same constant $C$
therefore survives.  This argument supplies no estimate for an arbitrary
continuous test function outside a difference of two such semipositive
metrics.
\end{proof}

\medskip
\begin{corollary}
\label{cor:p171-production-family-CLN}
Let $X/A$ be flat and projective, let the line bundles above come from $X$,
and assume that for every $0\leq i\leq n$ both metrics in the pair
$\phi_i,\phi'_i$ are continuous psh-regularizable and semipositive on the
fibres.  The same constant works on every fibre for the test
$\phi_0-\phi'_0$.  It is not asserted for an arbitrary continuous test
function.  At $\rho(x)>0$ this is \cref{lem:CLN-r3}; at
$\rho(x)=0$ it is \cite[Lemma~4.10]{PS23Hybrid} and its continuous
semipositive extension.  No division by $\rho$ occurs on the endpoint
locus.
\end{corollary}

\begin{proof}
For $n=0$, $C=0$.  For $n\geq1$, the constant is made only from mixed
algebraic intersection numbers, which are fibre-independent by flatness.
\cref{lem:CLN-r3} has no residual power parameter.  On the
non-Archimedean locus, \cite[Lemma~4.10]{PS23Hybrid} starts with FS metrics in
every slot, including the slot-zero pair.  Uniform FS approximation of all
these semipositive metrics and the same inequality pass the estimate to the
stated continuous domain.
\end{proof}


\subsubsection{Reduction to closed-immersion data}
\label{sec:hybrid-closed-immersion}
\begin{proposition}
\label{prop:p171-closed-immersion-finite-level-reduction}
Let \(S\) be a scheme, let \(X\) be a projective \(S\)-scheme, and fix a
closed immersion
\[
 \iota:X\hookrightarrow\mathbf P^r_S.
\]
Put \(H=\iota^*\mathcal O_{\mathbf P^r_S}(1)\), let
\(t_0,\ldots,t_r\) be the pulled-back coordinate sections, and, whenever
metrics are considered, put
\[
                         \eta=\max_{0\leq k\leq r}\log|t_k|.
\]
Let \(L\) be a line bundle on \(X\), let \(m\geq1\), and let
\(s_0,\ldots,s_N\in H^0(X,mL)\) have no common zero.  For real numbers
\(a_0,\ldots,a_N\), put
\[
                 \phi=\frac1m\max_{0\leq j\leq N}
                              \bigl(\log|s_j|+a_j\bigr).
\]
For every multi-index
\(\alpha=(\alpha_0,\ldots,\alpha_r)\in\mathbf N^{r+1}\) with
\(|\alpha|=m\), write \(t^\alpha=\prod_k t_k^{\alpha_k}\).  Then the
sections
\[
           s_jt^\alpha\in H^0\bigl(X,m(L+H)\bigr),
           \qquad 0\leq j\leq N,\quad |\alpha|=m,
\]
define a closed immersion into a projective space over \(S\), and their
Fubini--Study metric is exactly
\begin{equation}
\label{eq:p171-closed-immersion-metric-identity}
 \phi+\eta
 =\frac1m\max_{j,\,|\alpha|=m}
       \bigl(\log|s_jt^\alpha|+a_j\bigr).
\end{equation}
Thus every tropical Fubini--Study metric is the difference of two
tropical Fubini--Study metrics whose defining systems give closed
immersions:
\begin{equation}
\label{eq:p171-closed-immersion-difference}
                         \phi=(\phi+\eta)-\eta.
\end{equation}

Consequently, let \(A\) be an integral Banach ring, let \(X/A\) be an
\(n\)-dimensional projective scheme, and assume in addition that \(X\) is
flat over \(A\), as in
\cite[Definition~4.8 and Question~4.9]{PS23Hybrid}.  Let
\(L_1,\ldots,L_n\) carry tropical
Fubini--Study metrics \(\phi_1,\ldots,\phi_n\), and let \(\psi\) be a DFS
function.  Over every complete residue field of \(A\), the finite-level
quantity
\[
                       \int\psi\,
                       \operatorname{MA}(\phi_1,\ldots,\phi_n)
\]
is a finite signed sum of differences of Deligne-metric weights in which
every metric in every slot is induced by a system defining a closed
immersion.  The same finite expansion works simultaneously over the whole
Berkovich spectrum of \(A\); no extension of a resultant theorem from
closed immersions to arbitrary basepoint-free systems is needed for the
finite-level reduction.
\end{proposition}

\begin{remark}
\label{rem:p171-closed-immersion-reduction-scope}
This proposition is only a reduction.  By itself, it neither constructs a
relative resultant over an arbitrary Banach base nor identifies its fibre
with the fundamental-cycle resultant, and it does not prove that its
fieldwise Chow--Deligne norm assembles continuously.  These remaining
steps are supplied below.
\end{remark}

\begin{proof}
We prove \cref{prop:p171-closed-immersion-finite-level-reduction}.
The sections \(s_0,\ldots,s_N\) define an \(S\)-morphism
\[
                         f_s:X\longrightarrow\mathbf P^N_S.
\]
The degree-\(m\) monomials in the \(t_k\) define the composite of
\(\iota\) with the degree-\(m\) Veronese immersion
\[
       \nu_m\circ\iota:X\hookrightarrow\mathbf P^{R_m}_S,
       \qquad R_m=\binom{r+m}{m}-1.
\]
The product morphism
\[
        (f_s,\nu_m\circ\iota):X
        \longrightarrow\mathbf P^N_S\times_S\mathbf P^{R_m}_S
\]
is a closed immersion.  Indeed,
\((f_s,\operatorname{id}_X):X\to\mathbf P^N_S\times_S X\) is the graph
of \(f_s\), hence a closed immersion because \(\mathbf P^N_S\) is
separated over \(S\); the morphism
\[
 \mathbf P^N_S\times_S X
 \longrightarrow
 \mathbf P^N_S\times_S\mathbf P^{R_m}_S
\]
induced by \(\nu_m\circ\iota\) is a base change of a closed immersion.
Their composite is the displayed product morphism.  Composing it with the
Segre immersion gives a closed immersion whose homogeneous coordinate
sections are precisely the products \(s_jt^\alpha\).  This proves the
first assertion.

At every analytic point, one has
\[
 \max_{|\alpha|=m}\log|t^\alpha|
 =m\max_{0\leq k\leq r}\log|t_k|.
\]
The inequality from left to right follows because every \(t^\alpha\) is a
product of \(m\) factors; equality follows because the list of monomials
contains \(t_k^m\) for every \(k\).  Therefore
\begin{align*}
 \frac1m\max_{j,\,|\alpha|=m}
      \bigl(\log|s_jt^\alpha|+a_j\bigr)
 &=\frac1m\max_j\bigl(\log|s_j|+a_j\bigr)
   +\frac1m\max_{|\alpha|=m}\log|t^\alpha|\\
 &=\phi+\eta,
\end{align*}
which proves \eqref{eq:p171-closed-immersion-metric-identity}.  The
coordinate sections \(t_k\) themselves define the fixed closed immersion
\(\iota\), so both terms on the right of
\eqref{eq:p171-closed-immersion-difference} have closed-immersion defining
systems.  This proves the metric assertion, including arbitrary real
weights.

We prove the final reduction explicitly.  By
\cite[Definition~2.15]{PS23Hybrid}, write the DFS function as
\(\psi=\alpha-\beta\), where \(\alpha\) and \(\beta\) are tropical
Fubini--Study metrics on isomorphic semiample line bundles; fix the
isomorphism implicit in this DFS presentation.  Applying the first part to
both metrics and using the same \(H,\eta\) gives
\begin{equation}
\label{eq:p171-closed-immersion-dfs-expansion}
                    \psi=(\alpha+\eta)-(\beta+\eta),
\end{equation}
and the two systems on the right define closed immersions.  For each
\(1\leq i\leq n\), put
\[
                  \theta_i^+=\phi_i+\eta,
                  \qquad \theta_i^-=\eta.
\]
Then both defining systems are closed immersions and
\(\phi_i=\theta_i^+-\theta_i^-\).  Symmetry and multilinearity of the
fieldwise Monge--Amp\`ere pairing on continuous dpsh metrics over every
complete valued field
\cite[Proposition~8.5, Definition~8.7, and the paragraph following
Example~8.10]{BE21} give
\begin{equation}
\label{eq:p171-closed-immersion-ma-expansion}
 \operatorname{MA}(\phi_1,\ldots,\phi_n)
 =\sum_{\varepsilon\in\{+,-\}^n}
   (-1)^{\#\{i:\varepsilon_i=-\}}
   \operatorname{MA}(\theta_1^{\varepsilon_1},\ldots,
                     \theta_n^{\varepsilon_n}).
\end{equation}
For every summand, all displayed metrics are continuous dpsh metrics over
the complete residue field.  The fieldwise Deligne change-of-metric formula
\cite[Theorem~8.16 and formula~(8.9)]{BE21} therefore identifies the integral
of the difference in
\eqref{eq:p171-closed-immersion-dfs-expansion} against that summand with
the difference of the corresponding two Deligne-metric weights.  This is
the asserted finite signed expansion.

All sections, weights, tensor products, and the fixed embedding used in
the construction are defined over \(A\).  Restriction to any complete
residue field preserves the displayed formulas and the closed-immersion
property by base change.  Hence the index set and signs in
\eqref{eq:p171-closed-immersion-ma-expansion} are independent of the point
of the Berkovich spectrum.  This proves the simultaneous assertion.
\end{proof}


\subsubsection{Fieldwise and Banach-base comparison}
\label{sec:hybrid-fieldwise-comparison}
\begin{lemma}
\label{lem:p171-module-c-field-comparison}
Let \(K\) be a complete valued field and let \(Z\) be a projective
\(K\)-scheme of dimension \(n\).  Neither purity nor reducedness is
assumed.  Write
\[
 [Z]_n=\sum_T\ell(\mathcal O_{Z,\eta_T})[T].
\]
For \(0\leq i\leq n\), let \(E_i\) be a nonzero finite-dimensional
\(K\)-vector space, let \(E_i\otimes_K\mathcal O_Z\twoheadrightarrow L_i\)
define a closed immersion \(Z\hookrightarrow\mathbf P(E_i)\), and equip
\(E_i\), in the non-Archimedean case, with a diagonalizable ultrametric norm
(an orthogonal basis is fixed).  In the Archimedean case
\(K=K_\rho\), equip it with a powered norm
\(\lVert\cdot\rVert_{i,\rho}=\lVert\cdot\rVert_{i,\infty}^{\rho}\), where
\(\lVert\cdot\rVert_{i,\infty}\) is induced by an inner product.  Let
\(\vartheta_i\) be the induced orthogonal quotient metric on \(L_i\).  Let
\[
 \mathcal R\in I_{Z/K}(L_0,\ldots,L_n)\otimes_K
 \bigotimes_{i=0}^n\operatorname{Sym}^{\delta_i}(E_i^\vee)
\]
be the positive line-valued Chen--Moriwaki resultant
\cite[Definition~1.6.3]{ChenMoriwaki21AdelicIntersection} of \([Z]_n\).  Put
\[
 \delta_i=(L_0\cdots\widehat L_i\cdots L_n)[Z]_n.
\]
If \(e\ne0\) is a frame of \(I_{Z/K}(L_0,\ldots,L_n)\) and
\(\mathcal R=eF_e\), then
\begin{equation}
 \lambda_K(e;\vartheta_0,\ldots,\vartheta_n)=\Psi_K(F_e).
 \label{eq:p171-module-c-field-identity}
\end{equation}
Here \(\lambda_K=-\log\lVert\,\cdot\,\rVert\) is the additive
Boucksom--Eriksson Deligne weight \cite[Section~8.3]{BE21}, and the right
side is the following literal functional.
\begin{enumerate}
\item If the valuation is nontrivially non-Archimedean,
\[
 \Psi_K(F_e)=\log\lVert F_e\rVert_\varepsilon,
\]
where the norm is the epsilon tensor product of the epsilon symmetric
quotient norms on the dual spaces.
\item If the valuation is trivial, the same equality and the same epsilon
norm hold.
\item If \(K=K_\rho\) is \(\mathbf R\) or \(\mathbf C\) with
\(\lvert\cdot\rvert_\rho=\lvert\cdot\rvert_\infty^\rho\),
\(0<\rho\leq1\), unpower the norms and metrics by \(1/\rho\).  In
coordinates orthonormal for the unpowered norms,
\begin{equation}
 \Psi_{K_\rho}(F_e)
 =\rho\int_{\prod_iS(E_{i,\mathbf C})}
       \log\lvert F_e(z_0,\ldots,z_n)\rvert_\infty\,
       \mathrm d\eta(z)
 +\frac\rho2\sum_{i=0}^n\delta_i
       \sum_{\ell=1}^{\dim E_i-1}\frac1\ell.
 \label{eq:p171-module-c-powered-arch}
\end{equation}
For \(K=\mathbf R\), the spheres and polynomial are complexified and the
measure is pushed forward as in \cite[Remark~8.1]{BE21}; there is no factor
two.
\end{enumerate}
The assertions include \(n=0\).  In that case the Deligne line is the
positive finite norm line and all residue degrees and nonreduced Artin
lengths occur.  In every dimension the two sides of
\eqref{eq:p171-module-c-field-identity} are linear in the displayed
fundamental cycle; components of dimension below \(n\), whether embedded
or isolated, contribute no additional term.
\end{lemma}

\begin{theorem}
\label{thm:p171-module-c-synchronized-comparison}
Let \(A\) be an integral Banach ring and let
\(f:X\to\operatorname{Spec}A\) be flat and projective of finite
presentation, with nonempty fibres of dimension \(n\).  Fibre schemes may
be nonpure and nonreduced.  For \(0\leq i\leq n\), let \(M_i\) be an ample
line bundle, let \(m_i\geq1\), and let
\[
 s_{ij}\in H^0(X,M_i^{m_i}),\qquad 0\leq j\leq N_i,
\]
define a closed immersion on \(X\).  Fix real weights \(a_{ij}\), put
\(B_i=M_i^{m_i}\), and let
\[
 \theta_i=\frac1{m_i}\max_j\{\log\lvert s_{ij}\rvert+a_{ij}\}
\]
be the corresponding tropical Fubini--Study metric.  Put
\[
 d_i=(M_0\cdots\widehat M_i\cdots M_n)[X_x]_n,
\]
where the right side is independent of \(x\in\mathcal M(A)\); for \(n=0\)
it is the degree of \([X_x]_0\).  Set
\[
 \mathscr L=I_{X/A}(M_0,\ldots,M_n).
\]

There is \(p_0\) such that, for every \(p\geq p_0\),
\cref{cor:p171-revised2-veronese-resultants}, applied at one fixed
Noetherian stage, gives a raw tensor
\begin{equation}
 \mathcal R_p^{\mathrm{raw}}\in
 I_{X/A}(B_0^p,\ldots,B_n^p)\otimes_A
 \bigotimes_{i=0}^n\operatorname{Sym}^{\delta_{i,p}}(E_{i,p}^\vee),
 \label{eq:p171-module-c-relative-tensor}
\end{equation}
where
\begin{align}
 E_{i,p}&=\operatorname{Sym}^p(A^{N_i+1}),&
 K_{i,p}&=\operatorname{rank}E_{i,p}=\binom{N_i+p}{p},
 \label{eq:p171-module-c-veronese-module}\\
 Q_p&=\prod_{j=0}^n(m_jp),&
 \delta_{i,p}&=\frac{Q_p}{m_ip}d_i.
 \label{eq:p171-module-c-powers}
\end{align}
Let
\begin{equation}
 \iota_p:I_{X/A}(B_0^p,\ldots,B_n^p)
 \xrightarrow{\sim}\mathscr L^{\otimes Q_p}
 \label{eq:p171-module-c-line-power}
\end{equation}
be the ordered positive multiadditivity isomorphism, and put
\[
 \mathcal R_p=(\iota_p\otimes1)(\mathcal R_p^{\mathrm{raw}})
 \in\mathscr L^{\otimes Q_p}\otimes_A
 \bigotimes_i\operatorname{Sym}^{\delta_{i,p}}(E_{i,p}^\vee).
\]
For every field fibre, \(\mathcal R_p^{\mathrm{raw}}\) is the positive
fundamental-cycle resultant, and \(\mathcal R_p\) is its transport by
\(\iota_p\), with every generic length retained.

Let \(U\subset\mathcal M(A)\) be an analytic open on which
\(\mathscr L^{\mathrm{an}}\) has a frame \(e\), refined so that the finitely
many coefficients below lie in one local Banach analytic algebra.  With
\(e_{i,\beta}\), \(\lvert\beta\rvert=p\), denoting the monomial basis of
\(E_{i,p}\), and \(U_{i,\beta}\) its dual coordinate, write
\begin{align}
 \mathcal R_p&=e^{\otimes Q_p}F_{p,e},&
 F_{p,e}&=\sum_\alpha c_{p,\alpha}\mathbf U^\alpha,
 \label{eq:p171-module-c-local-polynomial}\\
 a_{i,\beta}&=\sum_j\beta_ja_{ij},&
 w_{p,\alpha}&=\sum_{i,\beta}\alpha_{i,\beta}a_{i,\beta}.
 \label{eq:p171-module-c-monomial-weights}
\end{align}
Every coefficient block has degree \(\delta_{i,p}\), and the coefficient
tuple is nowhere zero on \(U\).

For \(x\in U\), put
\begin{equation}
 \rho(x)=\frac{\log\max\{1,\lvert2\rvert_x\}}{\log2}.
 \label{eq:p171-module-c-rho}
\end{equation}
If \(\rho(x)=0\), define
\begin{equation}
 \Phi_{p,e}(x)=\max_\alpha
 \{\log\lvert c_{p,\alpha}(x)\rvert_x+w_{p,\alpha}\},
 \qquad \log0:=-\infty.
 \label{eq:p171-module-c-phi-na}
\end{equation}
If \(\rho(x)>0\), choose a character on the local Banach analytic algebra
containing the coefficients, with
\(\lvert b\rvert_x=\lvert\chi_x(b)\rvert_\infty^{\rho(x)}\), and define
\begin{align}
 G_{p,e,x}
 &=\sum_\alpha\chi_x(c_{p,\alpha})
   \exp(w_{p,\alpha}/\rho(x))\mathbf U^\alpha,
 \label{eq:p171-module-c-arch-polynomial}\\
 \Phi_{p,e}(x)
 &=\rho(x)\left(
 \int_{\prod_iS^{2K_{i,p}-1}}\log\lvert G_{p,e,x}(z)\rvert_
 \infty\,\mathrm d\eta(z)
 +\frac12\sum_i\delta_{i,p}
       \sum_{\ell=1}^{K_{i,p}-1}\frac1\ell\right).
 \label{eq:p171-module-c-phi-arch}
\end{align}
Then \(\Phi_{p,e}\) is choice-independent and continuous.  On an overlap
\(e'=ue\), one has the exact frame law
\begin{equation}
 F_{p,e'}=u^{-Q_p}F_{p,e},\qquad
 \Phi_{p,e'}=\Phi_{p,e}-Q_p\log\lvert u\rvert.
 \label{eq:p171-module-c-frame-law}
\end{equation}

For every \(x\), give \(E_{i,p,x}\) the diagonal ultrametric norm
\(\lVert e_{i,\beta}\rVert=\exp(-a_{i,\beta})\) when
\(\rho(x)=0\).  When \(\rho(x)>0\), give it the unpowered Hermitian norm
in which the monomial basis is orthogonal and
\(\lVert e_{i,\beta}\rVert_\infty
 =\exp(-a_{i,\beta}/\rho(x))\), form the unpowered quotient metric, and
then repower that metric by \(\rho(x)\).  In either case let
\(\widetilde\theta_{i,p,x}\) be the resulting quotient metric on
\(B_{i,x}^p=M_{i,x}^{m_ip}\), and put
\begin{equation}
 \theta_{i,p,x}=\frac1{m_ip}\widetilde\theta_{i,p,x}.
 \label{eq:p171-module-c-normalized-metric}
\end{equation}
If
\[
 \lambda_{p,e}(x)=
 \lambda_x(e(x);\theta_{0,p,x},\ldots,\theta_{n,p,x}),
\]
then the exact tensor/metric comparison is
\begin{equation}
 \boxed{\lambda_{p,e}(x)=Q_p^{-1}\Phi_{p,e}(x)}.
 \label{eq:p171-module-c-main-identity}
\end{equation}
It holds separately on trivially valued, nontrivially non-Archimedean,
powered real, and powered complex fibres.

Moreover, on every fibre
\begin{equation}
 0\leq\theta_{i,p,x}-\theta_{i,x}
 \leq\rho(x)\frac{\log K_{i,p}}{2m_ip}
 \leq\epsilon_{i,p}:=\frac{\log K_{i,p}}{2m_ip},
 \label{eq:p171-module-c-softmax}
\end{equation}
where the first difference is identically zero when \(\rho(x)=0\).  Put
\[
 \lambda_e(x)=
 \lambda_x(e(x);\theta_{0,x},\ldots,\theta_{n,x}).
\]
Then
\begin{equation}
 0\leq\lambda_{p,e}(x)-\lambda_e(x)
 \leq\sum_{i=0}^nd_i\epsilon_{i,p}\xrightarrow[p\to\infty]{}0
 \label{eq:p171-module-c-deligne-error}
\end{equation}
uniformly in \(x\).  Consequently the local functions \(\lambda_e\) are
continuous and obey
\(\lambda_{ue}=\lambda_e-\log\lvert u\rvert\); they define a continuous
additive metric on \(\mathscr L^{\mathrm{an}}\).
\end{theorem}

\begin{corollary}
\label{cor:p171-module-c-finite-level-continuity}
Let \(A\) be an integral Banach ring and let \(X/A\) be flat, projective,
and of finite type, with nonempty fibres of dimension \(n\).  Let
\(\phi_1,\ldots,\phi_n\) be tropical Fubini--Study metrics on line bundles
over \(X\), and let \(\psi\) be a DFS function.  Then
\begin{equation}
 x\longmapsto
 \int_{X_x^{\mathrm{an}}}\psi_x\,
 \operatorname{MA}_x(\phi_{1,x},\ldots,\phi_{n,x})
 \label{eq:p171-module-c-finite-level-function}
\end{equation}
is continuous on \(\mathcal M(A)\).  This includes \(n=0\) and the
fundamental-cycle multiplicities of every fibre.  The conclusion is only
the finite-level input; no density passage to arbitrary continuous psh
metrics or arbitrary continuous test functions is asserted here.
\end{corollary}

\begin{proof}
We first compare the two fieldwise constructions in
\cref{lem:p171-module-c-field-comparison}.  We use the same ordered
intersection line on both sides.  \cite[Section~8.3 and Appendix~A]{BE21}
defines the Deligne pairing by the determinant finite difference used in
\cite{Ducrot97CubeIntersection}; \cite[Theorem~A.22]{BE21} gives its
positive restriction isomorphism, and \cite[Example~8.15]{BE21} identifies
its dimension-zero member with the positive finite norm.  Thus the
evaluation element below is not silently replaced by an element of the dual
line.

We need the following comparison, which we prove rather than cite as a
change of terminology.  The recursive calculation is first made after the
complete isometric scalar extension of
\cite[Definition~3.6.1]{ChenMoriwaki21AdelicIntersection}, so that
\cite[Remark~3.7.3 and formula~(3.16)]{ChenMoriwaki21AdelicIntersection}
applies.  Descent to the original nontrivially valued field is justified
below by
\cite[Proposition~3.8.9 and Remark~3.8.10]{ChenMoriwaki21AdelicIntersection}
and by the corresponding Deligne base-change argument.  Suppose that
\(t=(t_0,\ldots,t_n)\), with \(t_i\in E_i\), is such that the first
\(n\) sections form a regular sequence on all of \(Z\) and the total
intersection is empty.  Let
\[
 e_t\in I_{Z/K}(L_0,\ldots,L_n)\setminus\{0\}
\]
be obtained by successive positive restriction and the final finite norm.
Then
\begin{equation}
 h^{t}_{\mathrm{CM}}(Z;\vartheta_0,\ldots,\vartheta_n)
 =\lambda_K(e_t;\vartheta_0,\ldots,\vartheta_n).
 \label{eq:p171-module-c-height-weight}
\end{equation}
For \(n=0\), over the complete algebraic closure,
\cite[formula~(3.15)]{ChenMoriwaki21AdelicIntersection} is
\(-\int\log\lVert t_0\rVert_{\vartheta_0}\,\delta_Z\).
\cite[Theorem~8.16, formula~(8.11)]{BE21} says that the logarithm of the norm
of \(e_t\) is the same integral with the positive sign.  Taking its negative
proves \eqref{eq:p171-module-c-height-weight} there.  The base-change
argument below descends the equality to \(K\), and
\cite[Example~8.2]{BE21} identifies the descended measure as
\[
 \delta_Z=\sum_z\operatorname{length}(A_z)[k_z:K]\,\delta_z
\]
for a finite scheme \(Z=\operatorname{Spec}\prod_zA_z\).  Thus the descended
base case retains every residue degree and every nonreduced Artin length.

Assume the equality in dimensions below \(n\).  The local intersection
number is symmetric by
\cite[Remark~3.6.2(1)]{ChenMoriwaki21AdelicIntersection}.  We may therefore
apply
\cite[Definition~3.7.1 and formula~(3.16)]{ChenMoriwaki21AdelicIntersection}
successively in the order
\(t_0,\ldots,t_{n-1}\).  At the \(j\)-th step, \(t_j\) is regular on the
entire current quotient, so its effective Cartier cycle is the top
fundamental cycle of that quotient, with the coefficients supplied by
Cartier intersection.  This recursion restricts the remaining
local height to that cycle and subtracts the integral of
\(\log\lVert t_j\rVert\) against the mixed measure in the other slots.
\cite[Proposition~8.5 and Theorem~8.16, formula~(8.8)]{BE21} applies the
same positive restriction to the Deligne evaluation element and subtracts
the same integral.  After \(n\) restrictions, \(t_n\) is nonvanishing on the
finite scheme and the already proved dimension-zero comparison applies.
Induction proves \eqref{eq:p171-module-c-height-weight}.  This also shows
directly that the comparison is linear in
\([Z]_n=\sum_T\ell(\mathcal O_{Z,\eta_T})[T]\).  The analytic measures in
\cite[formula~(8.1) and Proposition~8.5]{BE21} use that same cycle.
Thus a generic nilpotent thickness is counted by its length, while a
component of dimension below \(n\), whether embedded or isolated,
contributes no extra summand.

We also record the field-extension step, since normalization after an
unsupported extension would leave a gap.  If \(F/K\) is a complete
isometric valued extension, the algebraic Deligne line commutes with base
change.  The pullback of its metric satisfies the defining restriction
formula of \cite[Theorem~8.16]{BE21}:
\cite[Proposition~8.3(iv)]{BE21} pushes the mixed measure on \(Z_F\) to the
mixed measure on \(Z\), and algebraic restriction commutes with base change.
\cite[Theorem~8.16, formula~(8.11)]{BE21} gives the same assertion in
dimension zero.  Uniqueness in \cite[Theorem~8.16]{BE21} therefore gives
\begin{equation}
 \lambda_F(e\otimes1;\vartheta_{0,F},\ldots,\vartheta_{n,F})
 =\lambda_K(e;\vartheta_0,\ldots,\vartheta_n).
 \label{eq:p171-module-c-deligne-base-change}
\end{equation}
On the Chen--Moriwaki side,
\cite[Proposition~3.8.9]{ChenMoriwaki21AdelicIntersection} proves invariance
for non-Archimedean complete valued extensions, and
\cite[Remark~3.8.10]{ChenMoriwaki21AdelicIntersection} gives the Archimedean
case.  \cite[Definition~3.8.2 and Remark~3.8.3]{ChenMoriwaki21AdelicIntersection}
gives the corresponding scalar-extension identity for the epsilon polynomial
norm.  These are the required transport results if a normalizing tuple is
selected only after a complete isometric valued extension;
\cite[Remark~1.6.8]{ChenMoriwaki21AdelicIntersection} alone would not
suffice.
Consequently the recursively proved equality
\eqref{eq:p171-module-c-height-weight} descends from the complete algebraic
closure to every nontrivially valued complete field (and from \(\mathbf C\)
to \(\mathbf R\)).

Assume first that the valuation is nontrivial.  Then \(K\) is infinite.
Choose the sections successively.  Before choosing \(t_j\), form, for every
subset \(I\subset\{0,\ldots,j-1\}\), the scheme-theoretic intersection
\(Z_I=\bigcap_{i\in I}V(t_i)\), with \(Z_\varnothing=Z\), and avoid the
vanishing subspaces attached to every point in the finite union
\(\bigcup_I\operatorname{Ass}(Z_I)\).  Each forbidden set is a proper
linear subspace because the system is basepoint-free, and a finite union of
proper subspaces does not exhaust a vector space over an infinite field.
Thus every chosen section is a nonzerodivisor on every nonempty preceding
subset intersection.  In particular, each \(t_j\) is Cartier on \(Z\), the
family of divisors intersects properly, and
\(t_0,\ldots,t_{n-1}\) is a regular sequence on the full current
quotients.  After \(n\) cuts the full intersection has dimension at most
zero.  Choose \(t_n\) subject to the same associated-point conditions and
also nonvanishing at every point of that finite support.  The total
intersection is empty.  The same construction may be made after a complete
isometric valued extension; in that form descent uses
\cite[Proposition~3.8.9 and Remark~3.8.10]{ChenMoriwaki21AdelicIntersection}.

Put \(c_t=F_e(t)\).  Positive normalization of the line-valued resultant
is the pair of literal equalities
\begin{equation}
 e_t=e\,c_t,\qquad R_Z^{t}=F_e/c_t,
 \qquad R_Z^{t}(t)=1.
 \label{eq:p171-module-c-normalization-cancellation}
\end{equation}
The scalar \(c_t\) is nonzero.  For a nontrivially non-Archimedean
valuation, \cite[Theorem~3.8.7]{ChenMoriwaki21AdelicIntersection} and
\eqref{eq:p171-module-c-height-weight} give
\[
 \lambda_K(e_t)=\log\lVert R_Z^{t}\rVert_\varepsilon
 =\log\lVert F_e\rVert_\varepsilon-\log\lvert c_t\rvert.
\]
On the other hand, the frame law for the Deligne metric and the first
identity in \eqref{eq:p171-module-c-normalization-cancellation} give
\(\lambda_K(e_t)=\lambda_K(e)-\log\lvert c_t\rvert\).  The two scalar
terms cancel, proving \eqref{eq:p171-module-c-field-identity} with the
claimed sign.

For the usual complex absolute value,
\cite[Theorem~3.8.7]{ChenMoriwaki21AdelicIntersection} gives
\(h^t_{\mathrm{CM}}=\operatorname{ph}(R_Z^t)\), where
\(\operatorname{ph}\) is exactly the sphere integral plus the displayed
harmonic correction.  Its logarithmic homogeneity and
\eqref{eq:p171-module-c-normalization-cancellation} give the same
cancellation.  The real assertion follows from the scalar extension to
\(\mathbf C\) in \cite[Remark~3.8.10]{ChenMoriwaki21AdelicIntersection},
\cite[Remark~8.1]{BE21}, and
\eqref{eq:p171-module-c-deligne-base-change}; conjugation changes neither
the integral nor the resultant.  In particular, conjugate points are
pushed forward and not averaged, so no factor two occurs.

For the powered Archimedean field \(K_\rho\), unpower every norm and metric.
By \cref{thm:powered-normalization}, \eqref{eq:PN1-r3} identifies the mixed
measure with the usual unpowered measure, while \eqref{eq:PN2-r3}
identifies the additive Deligne weight with \(\rho\) times the usual
weight.  The dimension-zero anchor is \cref{lem:zero-anchor-r3}.
Multiplying the whole usual formula by \(\rho\), including the harmonic
correction, proves \eqref{eq:p171-module-c-powered-arch}.

Finally suppose that \(K\) is trivially valued.  If a normalizing tuple is
not defined over \(K\), pass to the complete isometric nontrivially valued
extension \(F=K((T))\), whose \(T\)-adic absolute value is trivial on
\(K^\times\), and choose the tuple over the infinite field \(F\).  We do
not apply \cite[Proposition~3.8.9]{ChenMoriwaki21AdelicIntersection} across
this change, since the source scope of
\cite[Section~3.8]{ChenMoriwaki21AdelicIntersection} starts with a
nontrivial valuation.  Instead,
\eqref{eq:p171-module-c-deligne-base-change} transports the left side
itself.  Base change of the line-valued tensor, together with
\cite[Remark~1.6.8]{ChenMoriwaki21AdelicIntersection}, identifies
\(F_e\otimes 1\) with the unnormalized resultant polynomial in the
base-changed frame, while
\cite[Remark~3.8.3]{ChenMoriwaki21AdelicIntersection} preserves its epsilon
norm.  The already proved
nontrivially non-Archimedean identity over \(F\) therefore descends as an
unnormalized identity over \(K\).
\cite[Proposition~3.9.3, formula~(3.27)]{ChenMoriwaki21AdelicIntersection}
independently identifies this same epsilon norm with the trivially valued
local intersection number; its proof uses precisely a nontrivial Gauss
extension.  This proves the trivial branch without asserting an unsupported
transport of a chosen normalized local height.

When \(n=0\), all preceding inductions stop at
\cite[Theorem~8.16, formula~(8.11)]{BE21}, and the
line is the positive norm line.  If also \(\dim E_0=1\), the closed
immersion \(Z\hookrightarrow\mathbf P^0_K\) forces
\(Z=\operatorname{Spec}K\); the affine resultant is one linear coordinate,
the projective divisor is empty, and normalization makes its value the
literal unit \(1\).  Thus this endpoint has neither a hidden dualization
nor a missing degree.  This proves
\cref{lem:p171-module-c-field-comparison}.
\end{proof}

\begin{proof}
We prove \cref{thm:p171-module-c-synchronized-comparison}.  The relative
intersection numbers \(d_i\) are constant because they are mixed leading
coefficients of the fibre Hilbert polynomial of a flat projective family
over the connected scheme \(\operatorname{Spec}A\).  Apply
\cref{cor:p171-revised2-veronese-resultants}.  It chooses one Noetherian
integral stage and produces \(\mathcal R_p^{\mathrm{raw}}\); transport by
\(\iota_p\) gives \(\mathcal R_p\), for every sufficiently large \(p\).
On each field fibre it first applies the complete-system theorem
at that stage and only then performs affine monomial substitution.  Thus a
kernel in a monomial evaluation map is harmless and no projective map is
asserted.  The fibre is the positive resultant of \([X_x]_n\).

Ordered multiadditivity gives
\[
 I_{X/A}(M_0^{m_0p},\ldots,M_n^{m_np})
 \simeq I_{X/A}(M_0,\ldots,M_n)^{\otimes\prod_i(m_ip)},
\]
which is \eqref{eq:p171-module-c-line-power}.  The degree of the resultant
in block \(i\) is
\begin{align*}
 (B_0^p\cdots\widehat{B_i^p}\cdots B_n^p)[X_x]_n
 &=\left(\prod_{j\ne i}m_jp\right)d_i\\
 &=\frac{Q_p}{m_ip}d_i,
\end{align*}
which proves \eqref{eq:p171-module-c-powers}, including
\(Q_p=m_0p\) and \(\delta_{0,p}=d_0\) when \(n=0\).

We spell out the analytic localization used next.  The invertible
\(A\)-module underlying \(\mathscr L\) is trivial after inverting one
element \(b\notin\ker(x)\) at any prescribed
\(x\in\mathcal M(A)\).  Since \(\lvert b\rvert_x>0\), there is a rational
Banach analytic neighbourhood of \(x\) on which \(b\) is invertible.  The
chosen algebraic frame and the finitely many coefficients of
\(\mathcal R_p\) then belong to the Banach algebra of that neighbourhood.
These neighbourhoods cover \(\mathcal M(A)\); no finite rational cover and
no global analytic frame is required.

On a local analytic frame, the coefficients in
\eqref{eq:p171-module-c-local-polynomial} are analytic.  If they all
vanished at a point \(x\), the fibre tensor would be the zero polynomial,
contrary to the fibre assertion in
\cref{cor:p171-revised2-veronese-resultants}.  Hence the coefficient tuple
is nowhere zero.  Apply
\cref{thm:hybrid-weighted-philippon-gauss}, with block sizes
\(K_{i,p}\), degrees \(\delta_{i,p}\), and weights
\(w_{p,\alpha}\).  Its definitions are exactly
\eqref{eq:p171-module-c-rho}--\eqref{eq:p171-module-c-phi-arch}.  It proves
choice independence and continuity of \(\Phi_{p,e}\).  Since
\(e'^{\otimes Q_p}=u^{Q_p}e^{\otimes Q_p}\), equality of the two
expressions for \(\mathcal R_p\) gives
\(F_{p,e'}=u^{-Q_p}F_{p,e}\).  Logarithmic homogeneity gives the second
formula in \eqref{eq:p171-module-c-frame-law}, with the displayed minus
sign.

Fix \(x\in U\), put \(K=\mathcal H(x)\), and put
\(\widetilde e_{p,x}=\iota_{p,x}^{-1}(e(x)^{\otimes Q_p})\), a frame of
the raw intersection line
\(I_{X_x/K}(B_{0,x}^p,\ldots,B_{n,x}^p)\).  The definition by transport
gives the literal scalarization
\[
 \mathcal R_{p,x}^{\mathrm{raw}}
 =\widetilde e_{p,x}F_{p,e,x}.
\]
If \(\rho(x)=0\), the dual basis
vector \(U_{i,\beta}\) has norm \(\exp(a_{i,\beta})\).  The epsilon
symmetric quotient norm is exactly
\begin{equation}
 \log\lVert F_{p,e,x}\rVert_\varepsilon
 =\max_\alpha\{\log\lvert c_{p,\alpha}(x)\rvert_x+w_{p,\alpha}\}.
 \label{eq:p171-module-c-exact-gauss}
\end{equation}
For completeness, choose in each tensor word the sorted representative of a
commutative monomial.  This gives a lift with the displayed maximum norm.
Conversely, in any lift the coefficient of a commutative monomial is the
sum of the coefficients of its tensor-word representatives; the
ultrametric inequality bounds its weighted modulus by the norm of the
lift.  Taking the quotient infimum proves equality.  No multinomial
coefficient and no injectivity of the monomial evaluation map is involved.

If \(\rho(x)>0\), the vectors
\(\widehat e_{i,\beta}=\exp(a_{i,\beta}/\rho(x))e_{i,\beta}\)
are orthonormal for the unpowered Hermitian norm.  Their dual coordinates
satisfy
\(U_{i,\beta}=\exp(a_{i,\beta}/\rho(x))\widehat U_{i,\beta}\).
Consequently the coefficient of multiexponent \(\alpha\) in orthonormal
coordinates is multiplied by
\(\exp(w_{p,\alpha}/\rho(x))\), with the positive sign appearing in
\eqref{eq:p171-module-c-arch-polynomial}.  Thus
\cref{lem:p171-module-c-field-comparison}, including its powered real and
complex branches, gives in all four valuation cases
\begin{equation}
 \lambda_x\!\left(\widetilde e_{p,x};
       \widetilde\theta_{0,p,x},\ldots,
       \widetilde\theta_{n,p,x}\right)=\Phi_{p,e}(x).
 \label{eq:p171-module-c-raw-identity}
\end{equation}
The quotient metrics in this formula are continuous semipositive metrics,
so the cited Deligne metric is in the domain of
\cite[Theorem~8.16]{BE21}.

By definition,
\(\widetilde\theta_{i,p,x}=m_ip\,\theta_{i,p,x}\).  Metric
multiadditivity makes \(\iota_{p,x}\) an isometry and gives
\[
 \lambda_x\!\left(\widetilde e_{p,x};
       m_0p\,\theta_{0,p,x},\ldots,m_np\,\theta_{n,p,x}\right)
 =Q_p\lambda_x(e(x);\theta_{0,p,x},\ldots,\theta_{n,p,x}).
\]
Combining this equality with \eqref{eq:p171-module-c-raw-identity} proves
\eqref{eq:p171-module-c-main-identity}.  Notice that the exponent is
\(Q_p\), not a block degree and not its reciprocal.

It remains to prove the uniform approximation, independently of the
coefficient continuity.  Work at a point \(y\in X_x^{\mathrm{an}}\) in a
local frame \(\tau_i\) of \(M_i\), and write
\(s_{ij}=f_{ij}\tau_i^{m_i}\).  At a non-Archimedean point, including a
trivially valued point, the reciprocal quotient norm is
\[
 \max_{\lvert\beta\rvert=p}
 \lvert f_i^\beta(y)\rvert\exp(a_{i,\beta})
 =\left(\max_j\lvert f_{ij}(y)\rvert\exp(a_{ij})\right)^p.
\]
The equality uses the pure powers and multiplicativity.  After taking the
logarithm and dividing by \(m_ip\), it says literally
\(\theta_{i,p,x}=\theta_{i,x}\).

Now let \(\rho=\rho(x)>0\).  Under the canonical unpowered real/complex
analytic identification, choose a complex representative \(y_\infty\) of
\(y\) (the two representatives in the real case are conjugate), and put
\[
 b_j=\log\lvert f_{ij}(y_\infty)\rvert_\infty+a_{ij}/\rho,
 \qquad b_\beta=\sum_j\beta_jb_j.
\]
The Hilbert quotient formula for the declared diagonal monomial basis gives
\begin{equation}
 -\log\lVert\tau_i\rVert_{\theta_{i,p,x}}
 =\frac\rho{2m_ip}\log\sum_{\lvert\beta\rvert=p}\exp(2b_\beta).
 \label{eq:p171-module-c-softmax-formula}
\end{equation}
There are no multinomial coefficients: the norm was declared on the
monomial basis itself, rather than taken as the symmetric-power quotient of
the degree-one Hermitian norm.  Since there are \(K_{i,p}\) terms and the
pure powers give
\(\max_{\lvert\beta\rvert=p}b_\beta=p\max_jb_j\), one has
\[
 \max_\beta b_\beta
 \leq\frac12\log\sum_\beta\exp(2b_\beta)
 \leq\max_\beta b_\beta+\frac12\log K_{i,p}.
\]
The target frame weight is
\[
 -\log\lVert\tau_i\rVert_{\theta_{i,x}}
 =\rho m_i^{-1}\max_jb_j.
\]
Substitution in
\eqref{eq:p171-module-c-softmax-formula} proves
\eqref{eq:p171-module-c-softmax}, including its sign and factor \(1/2\).
The inequality \(\rho\leq1\) follows from the ordinary triangle inequality
\(\lvert2\rvert_x\leq2\).

Change the \(n+1\) normalized Deligne metrics one slot at a time from
\((\theta_i)_i\) to \((\theta_{i,p})_i\).
\cite[Theorem~8.16, formula~(8.9)]{BE21} expresses the change in slot \(i\)
as the integral of \(\theta_{i,p}-\theta_i\) against the positive mixed
measure in the other slots.  At a powered Archimedean fibre this is the
transported formula: \eqref{eq:PN2-r3} contributes one factor \(\rho\),
\eqref{eq:PN1-r3} leaves the unpowered mixed measure, and unpowering
\(\theta_{i,p}-\theta_i\) contributes \(1/\rho\), so no residual power
remains.  Its total mass is \(d_i\), independently of which
already changed semipositive metrics occur there.  Hence
\[
 0\leq\lambda_{p,e}(x)-\lambda_e(x)
 \leq\sum_i d_i\lVert\theta_{i,p,x}-\theta_{i,x}\rVert_\infty
 \leq\sum_i d_i\epsilon_{i,p},
\]
which is \eqref{eq:p171-module-c-deligne-error}.  Equivalently, on the raw
line the exact ledger is
\(\delta_{i,p}(m_ip)=Q_pd_i\), and division by \(Q_p\) gives the same
bound.  This check prevents loss of a Veronese exponent.  When \(n=0\),
the measure in the empty set of other slots has mass
\(d_0=\sum_z\operatorname{length}(A_z)[k_z:K]\), so the same calculation
is literal.  On a real fibre the push-forward convention retains this
mass and adds no factor two.

For fixed \(N_i\),
\(\log K_{i,p}=O(\log p)\), so the right side of
\eqref{eq:p171-module-c-deligne-error} tends to zero independently of
\(x\).  Each \(\lambda_{p,e}=Q_p^{-1}\Phi_{p,e}\) is continuous by the
weighted coefficient theorem.  Therefore \(\lambda_e\) is its uniform
limit and is continuous.  Dividing \eqref{eq:p171-module-c-frame-law} by
\(Q_p\) gives
\(\lambda_{p,ue}=\lambda_{p,e}-\log\lvert u\rvert\); passing to the
uniform limit proves the same law for \(\lambda_e\).  This completes the
proof of \cref{thm:p171-module-c-synchronized-comparison}.
\end{proof}

\begin{proof}
We prove \cref{cor:p171-module-c-finite-level-continuity}.  The source
hypotheses imply finite presentation.  Indeed,
\cite[Section~2]{PS23Hybrid} assumes that \(A\) is integral and \(X\) is
projective of finite type over \(A\).  Flatness and local finite type over an
integral base imply local finite presentation by
\cite[Tag~\texttt{05IN}]{stacks-project}.  Finite type is
quasi-compact, while projective is separated and hence quasi-separated.
The definition in \cite[Tag~\texttt{01TP}]{stacks-project} therefore makes
\(X\to\operatorname{Spec}A\) of finite presentation.

Apply \cref{prop:p171-closed-immersion-finite-level-reduction}.  Its exact
identities \eqref{eq:p171-closed-immersion-dfs-expansion} and
\eqref{eq:p171-closed-immersion-ma-expansion}, followed by the Deligne
change-of-metric formula, express
\eqref{eq:p171-module-c-finite-level-function} as a finite signed sum of
differences of Deligne weights.  The same finite index set and signs work
for every \(x\), and every defining section system in every slot gives a
closed immersion.  The underlying line bundles are therefore ample.

Apply \cref{thm:p171-module-c-synchronized-comparison} separately to each
of those finitely many \((n+1)\)-tuples.  Each Deligne weight in the signed
sum is continuous on \(\mathcal M(A)\); hence so is their finite signed
sum.  The argument also applies for \(n=0\), where it is the difference of
two positive finite-norm weights and includes the Artin lengths and residue
degrees recorded above.  No continuity of the desired integral was used in
constructing the coefficient functions: continuity first came from the
independently verified weighted polynomial theorem, the fieldwise
identification came next, and only then was the uniform limit taken.  Thus
the proof is not circular and establishes only the stated finite-level
conclusion.
\end{proof}


\subsubsection{Passage to continuous data}
\label{sec:hybrid-continuous-data}
\begin{theorem}
\label{thm:p171-full-assembly}
Let \(A\) be an integral Banach ring and let \(X\) be an
\(n\)-dimensional projective \(A\)-scheme of finite type, flat over
\(A\), in the sense of
\cite[Definition~4.8 and Question~4.9]{PS23Hybrid}.  Let
\(L_1,\ldots,L_n\) be semiample line bundles on
\(X\), let
\[
 \phi_i\in\operatorname{CPSH}(L_i),\qquad 1\leq i\leq n,
\]
and let \(\psi\in C^0(X^{\mathrm{an}})\).  Then the function
\begin{equation}
 \mathcal I_{\psi,\boldsymbol\phi}(x)
 =\int_{X_{\mathcal H(x)}^{\mathrm{an}}}
       \psi_x\,
       \operatorname{MA}_x
       \bigl((\phi_1)_x,\ldots,(\phi_n)_x\bigr)
 \label{eq:p171-full-target}
\end{equation}
is continuous on \(\mathcal M(A)\).

No purity or reducedness is assumed.  On each fibre all intersection
numbers, Deligne weights, and Monge--Amp\`ere measures are evaluated on
the top fundamental cycle
\begin{equation}
 [X_x]_n=\sum_T
 \operatorname{length}(\mathcal O_{X_x,\eta_T})[T].
 \label{eq:p171-full-fundamental-cycle}
\end{equation}
Thus every generic nilpotent multiplicity is retained; embedded and
isolated components of dimension below \(n\) do not create extra
top-cycle terms.  The assertion includes \(n=0\), with the empty mixed
Monge--Amp\`ere product equal to the degree measure of
\([X_x]_0\).
\end{theorem}

\begin{proof}
We prove \cref{thm:p171-full-assembly}.  We first check that the source
hypotheses place the family in the exact
finite-presentation scope of the preceding algebraic results.
Pille--Schneider assume that \(A\) is integral and that \(X\) is projective
of finite type over \(A\), while their question adds flatness
\cite[Section~2 and Question~4.9]{PS23Hybrid}.  Flatness and local finite
type over an integral base imply local finite presentation by
\cite[Tag~05IN]{stacks-project}.  Finite type makes the
morphism quasi-compact, while projectivity makes it separated and hence
quasi-separated.  The definition in \cite[Tag~01TP]{stacks-project}
therefore makes
\(X\to\operatorname{Spec}A\) a morphism of finite presentation.  No
Noetherian or reducedness hypothesis has been added.

We separate the finite-level assertion from the density passage.  Fix a
projective embedding
\[
 \iota:X\hookrightarrow\mathbf P_A^r,
 \qquad H=\iota^*\mathcal O(1),
 \qquad \eta=\max_{0\leq k\leq r}\log|t_k|,
\]
where the \(t_k\) are the pulled-back coordinate sections.  Let
\(\phi_i^\circ\) be a tropical Fubini--Study metric on \(L_i\), and let
\(\psi^\circ\) be a DFS function.  By
\cite[Definition~2.15]{PS23Hybrid}, after fixing the isomorphism in the
DFS presentation,
write
\[
 \psi^\circ=\alpha-\beta
\]
with \(\alpha\) and \(\beta\) tropical Fubini--Study metrics on the same
semiample line bundle \(H_0\).
\Cref{prop:p171-closed-immersion-finite-level-reduction} gives
\begin{align}
 \zeta^+&=\alpha+\eta,&
 \zeta^-&=\beta+\eta,&
 \psi^\circ&=\zeta^+-\zeta^-,
 \label{eq:p171-full-dfs-expansion}\\
 \theta_i^+&=\phi_i^\circ+\eta,&
 \theta_i^-&=\eta,&
 \phi_i^\circ&=\theta_i^+-\theta_i^-.
 \label{eq:p171-full-fs-expansion}
\end{align}
Write \(G_i^+=L_i+H\) and \(G_i^-=H\) for the underlying line
bundles of \(\theta_i^+\) and \(\theta_i^-\), respectively.
Every displayed metric is defined by a finite system whose associated
morphism is a closed immersion: the construction is the Segre product of
the original basepoint-free morphism with a Veronese power of \(\iota\).
Consequently each underlying line bundle in the displayed positive terms
is ample.

For \(\varepsilon=(\varepsilon_1,\ldots,\varepsilon_n)\in
\{+,-\}^n\), set
\[
 s(\varepsilon)=(-1)^{\#\{i:\varepsilon_i=-\}},
 \qquad
 \boldsymbol\theta^\varepsilon
 =(\theta_1^{\varepsilon_1},\ldots,
   \theta_n^{\varepsilon_n}).
\]
Symmetry and multilinearity of the fieldwise mixed Monge--Amp\`ere
pairing give
\begin{equation}
 \operatorname{MA}(\phi_1^\circ,\ldots,\phi_n^\circ)
 =\sum_{\varepsilon\in\{+,-\}^n}s(\varepsilon)
   \operatorname{MA}(\boldsymbol\theta^\varepsilon).
 \label{eq:p171-full-ma-expansion}
\end{equation}
For a fixed \(\varepsilon\), the two metrics \(\zeta^+\) and
\(\zeta^-\) live on the same line bundle.  Let
\[
 \mathscr D_\varepsilon
 =I_{X/A}(H_0+H,G_1^{\varepsilon_1},\ldots,
   G_n^{\varepsilon_n})
\]
be the corresponding algebraic intersection line and choose a local frame
\(e_\varepsilon\).  The Deligne change-of-metric formula
\cite[Theorem~8.16, formula~(8.9)]{BE21} identifies their integral difference with
the difference of the two additive Deligne weights.  Thus, on any local
frame of the relevant intersection line,
\begin{align}
 &\int_{X_x^{\mathrm{an}}}\psi_x^\circ\,
     \operatorname{MA}_x
       ((\phi_1^\circ)_x,\ldots,(\phi_n^\circ)_x)\notag\\
 &\quad=
 \sum_{\varepsilon\in\{+,-\}^n}s(\varepsilon)
 \left(
 \lambda_x(e_\varepsilon(x);
           \zeta_x^+,\boldsymbol\theta_x^\varepsilon)
 -\lambda_x(e_\varepsilon(x);
           \zeta_x^-,\boldsymbol\theta_x^\varepsilon)
 \right).
 \label{eq:p171-full-signed-weights}
\end{align}
The right side consists literally of \(2\cdot2^n=2^{n+1}\) signed
Deligne weights.  If \(n=0\), the set \(\{+,-\}^0\) has one element,
so \eqref{eq:p171-full-signed-weights} has exactly two weights.  For each
fixed \(\varepsilon\) the two frame-change terms cancel, and hence their
difference is a globally defined real function.

We now justify continuity of every weight occurring in
\eqref{eq:p171-full-signed-weights}.  Consider one of its
closed-immersion tuples and write it in the notation of
\cref{thm:p171-module-c-synchronized-comparison}:
\[
 \vartheta_i=\frac1{m_i}\max_j
 \{\log|s_{ij}|+a_{ij}\},
 \qquad
 s_{ij}\in H^0(X,M_i^{m_i}),
 \qquad 0\leq i\leq n.
\]
Each \(M_i\) is ample and each displayed section system defines a closed
immersion.  Put
\begin{equation}
 d_i=(M_0\cdots\widehat M_i\cdots M_n)[X_x]_n,
 \qquad
 Q_p=\prod_{j=0}^n(m_jp),
 \qquad
 \delta_{i,p}=\frac{Q_p}{m_ip}d_i.
 \label{eq:p171-full-degree-ledger}
\end{equation}
Flatness makes the \(d_i\), hence the block degrees, independent of \(x\).
For \(n=0\), the same ledger reads
\(Q_p=m_0p\) and \(\delta_{0,p}=d_0\), where \(d_0\) is the degree of
\([X_x]_0\).

Put
\[
 \mathscr L=I_{X/A}(M_0,\ldots,M_n).
\]
By \cref{cor:p171-revised2-veronese-resultants}, for every sufficiently
large \(p\) there is a raw tensor
\[
 \mathcal R_p^{\mathrm{raw}}\in
 I_{X/A}(M_0^{m_0p},\ldots,M_n^{m_np})\otimes_A
 \bigotimes_{i=0}^n
 \operatorname{Sym}^{\delta_{i,p}}(E_{i,p}^\vee),
 \qquad
 E_{i,p}=\operatorname{Sym}^p(A^{N_i+1}).
\]
Ordered positive multiadditivity gives the isomorphism
\[
 \iota_p:
 I_{X/A}(M_0^{m_0p},\ldots,M_n^{m_np})
 \xrightarrow{\sim}\mathscr L^{\otimes Q_p}.
\]
Set
\[
 \mathcal R_p=(\iota_p\otimes1)
  (\mathcal R_p^{\mathrm{raw}})
 \in\mathscr L^{\otimes Q_p}\otimes_A
 \bigotimes_{i=0}^n
 \operatorname{Sym}^{\delta_{i,p}}(E_{i,p}^\vee).
\]
The construction uses one fixed Noetherian stage.  On every residue-field
fibre,
\(\mathcal R_p^{\mathrm{raw}}\) is the positive Chen--Moriwaki resultant
\cite[Section~1.6]{ChenMoriwaki21AdelicIntersection} of \([X_x]_n\),
including every top generic length, and \(\mathcal R_p\)
is its transport under \(\iota_p\).  Affine monomial
substitution is performed only after the complete-fibre theorem; kernels
in the substitution map therefore require neither an inherited norm open
nor a projective coefficient map.

Let \(x_0\in\mathcal M(A)\).  Because only finitely many intersection
lines occur in \eqref{eq:p171-full-signed-weights}, there is one common
analytic neighbourhood \(U\) of \(x_0\), independent of \(p\), on which
all of them have frames.  For each fixed \(p\) and each point of \(U\),
refine locally to a Banach analytic domain whose algebra contains the
finitely many tensor coefficients.  These domains cover \(U\), so a
coefficient weight proved continuous on each of them is continuous on
\(U\).  For a frame \(e\) of \(\mathscr L^{\mathrm{an}}\), write
\[
 \mathcal R_p=e^{\otimes Q_p}F_{p,e},
 \qquad
 F_{p,e}=\sum_\gamma c_{p,\gamma}\mathbf U^\gamma.
\]
The fibre resultant is nonzero at every point, so the coefficient tuple is
nowhere zero on \(U\).
\Cref{thm:hybrid-weighted-philippon-gauss} therefore makes the following
coefficient functional continuous.  With
\[
 \rho(x)=\frac{\log\max\{1,|2|_x\}}{\log2},
 \qquad
 w_{p,\gamma}=\sum_{i,\beta}
       \gamma_{i,\beta}\sum_j\beta_ja_{ij},
\]
it is the weighted Gauss maximum
\[
 \Phi_{p,e}(x)=
 \max_\gamma\{\log|c_{p,\gamma}(x)|_x+w_{p,\gamma}\},
 \qquad \log0:=-\infty,
\]
when \(\rho(x)=0\).  When \(\rho(x)>0\), it is
\begin{align*}
 \Phi_{p,e}(x)=\rho(x)\Bigg(&
 \int_{\prod_iS^{2K_{i,p}-1}}
 \log\left|
  \sum_\gamma\chi_x(c_{p,\gamma})
  \exp(w_{p,\gamma}/\rho(x))\mathbf U^\gamma
 \right|_\infty\,\mathrm d\eta\\
 &+\frac12\sum_i\delta_{i,p}
    \sum_{\ell=1}^{K_{i,p}-1}\frac1\ell\Bigg),
 \qquad K_{i,p}=\binom{N_i+p}{p}.
\end{align*}
The factor \(\rho(x)\) multiplies both the sphere integral and the
harmonic correction.  The representing character is chosen only on the
local Banach algebra and is unique up to conjugation.  The weighted theorem
proves choice independence and continuity through \(\rho=0\).  On an
overlap (e'=ue), its exact frame law is
\begin{equation}
 F_{p,e'}=u^{-Q_p}F_{p,e},
 \qquad
 \Phi_{p,e'}=\Phi_{p,e}-Q_p\log|u|.
 \label{eq:p171-full-frame-law}
\end{equation}

\Cref{lem:p171-module-c-field-comparison,thm:p171-module-c-synchronized-comparison}
prove on the
trivially valued, nontrivially non-Archimedean, powered real, and powered
complex fibres that
\begin{equation}
 \lambda_{p,e}(x)=Q_p^{-1}\Phi_{p,e}(x).
 \label{eq:p171-full-module-c-identity}
\end{equation}
Here the quotient metrics use the diagonal monomial basis and all powers in
\eqref{eq:p171-full-degree-ledger}.  The real measure is pushed forward
from complexification with no factor two.  In dimension zero the identity
is the positive finite norm, including residue degrees and Artin lengths.

The same theorem proves, without using continuity of the target weight,
that its normalized quotient metrics satisfy
\begin{equation}
 0\leq\vartheta_{i,p,x}-\vartheta_{i,x}
 \leq
 \frac{\rho(x)\log K_{i,p}}{2m_ip}
 \leq\frac{\log K_{i,p}}{2m_ip},
 \label{eq:p171-full-softmax}
\end{equation}
with equality in the first comparison at every non-Archimedean point.
The positively oriented change-of-metric formula
\cite[formula~(8.9)]{BE21} and the
degree identity
\(\delta_{i,p}(m_ip)=Q_pd_i\) then give
\begin{equation}
 0\leq\lambda_{p,e}(x)-\lambda_e(x)
 \leq E_p:=\sum_{i=0}^n
 d_i\frac{\log K_{i,p}}{2m_ip}
 \xrightarrow[p\to\infty]{}0.
 \label{eq:p171-full-global-error}
\end{equation}
The number \(E_p\) is independent of \(x\) and of the chosen local frame.
Thus the approximation is global even though its coefficients and frames
are read on local Banach neighbourhoods.  Since
\(\log K_{i,p}=O(\log p)\), it converges to zero.  Equations
\eqref{eq:p171-full-module-c-identity} and
\eqref{eq:p171-full-global-error} show that every target Deligne weight is
a local uniform limit of continuous coefficient weights.  It is therefore
continuous, and division of \eqref{eq:p171-full-frame-law} by \(Q_p\)
gives the correct gluing law.  This order avoids circularity: algebraic
tensor first, independently continuous coefficient functional second,
fieldwise metric identity third, and only then the uniform limit.

\Cref{cor:p171-module-c-finite-level-continuity} applies to
\(\psi^\circ,\phi_1^\circ,\ldots,\phi_n^\circ\).  Equivalently, combine
the \(2^{n+1}\) continuous weights in
\eqref{eq:p171-full-signed-weights}.  Hence
\begin{equation}
 x\longmapsto
 \int_{X_x^{\mathrm{an}}}\psi_x^\circ\,
 \operatorname{MA}_x
 ((\phi_1^\circ)_x,\ldots,(\phi_n^\circ)_x)
 \label{eq:p171-full-finite-level}
\end{equation}
is continuous for every DFS test and every tuple of tropical
Fubini--Study metrics.  This proves the finite-level continuity needed
below.

We now pass to arbitrary data by a cofinal product net.
Let
\[
 M=(L_1\cdots L_n)[X_x]_n\geq0
\]
be the common fibre mass; for \(n=0\), this is the degree of
\([X_x]_0\).  By \cite[Theorem~2.16]{PS23Hybrid}, there is a net of DFS functions
\(\psi_q\to\psi\) uniformly.  Write
\(\psi_q=\alpha_q-\beta_q\) on one semiample line bundle \(H_q\), using
the isomorphism in its DFS presentation.  By
\cite[Proposition~2.21]{PS23Hybrid}, independently for each \(i\), there is a net of tropical
Fubini--Study metrics
\(\phi_{i,a_i}\to\phi_i\) uniformly.

\Cref{cor:p171-production-family-CLN} applies to the slot-zero pair
\((\alpha_q,\beta_q)\); the arbitrary continuous function \(\psi\) is
never inserted into that CLN slot.  Put
\begin{equation}
 C_q=
 \begin{cases}
  0,&n=0,\\[1mm]
  2\displaystyle\max_{1\leq i\leq n}
   (H_qL_1\cdots\widehat L_i\cdots L_n)[X_x]_n,&n\geq1.
 \end{cases}
 \label{eq:p171-full-cln-constant}
\end{equation}
Flatness makes \(M\) and every \(C_q\) independent of \(x\).  At powered
Archimedean points, \eqref{eq:PN1-r3} leaves the unpowered mixed measure
unchanged, while \eqref{eq:PN2-r3} contributes one factor \(\rho\), and the
reciprocal
\(1/\rho\) from unpowering the varying metric difference cancels it.
At \(\rho=0\), \cite[Lemma~4.10]{PS23Hybrid} supplies the endpoint without
division by \(\rho\).  Hence the same \(C_q\) works on every fibre.  No
bound uniform in \(q\) is asserted or needed.

Let \(Q\) and \(A_i\) be the directed sets indexing these nets and put
\begin{equation}
 P=(0,1]^{\mathrm{op}}\times Q\times\prod_{i=1}^nA_i
 \label{eq:p171-full-product-set}
\end{equation}
with coordinatewise order.  Let \(D\subset P\) consist of
\(d=(\epsilon,q,\mathbf a)\) satisfying
\begin{align}
 (1+M)\|\psi_q-\psi\|_\infty&\leq\epsilon/2,
 \label{eq:p171-full-D-test}\\
 (1+C_q)\sum_{i=1}^n
 \|\phi_{i,a_i}-\phi_i\|_\infty&\leq\epsilon/2.
 \label{eq:p171-full-D-metric}
\end{align}
The set \(D\) is directed.  Given two of its elements, first choose a
smaller \(\epsilon\), next a common later \(q\) satisfying
\eqref{eq:p171-full-D-test}, and only after that \(q\) is fixed choose
common later \(a_i\)'s satisfying \eqref{eq:p171-full-D-metric}.

The cofinality is equally literal.  Given
\(p_0=(\epsilon_0,q_0,\mathbf a_0)\in P\), first choose
\(0<\epsilon\leq\epsilon_0\).  Uniform convergence of the DFS net gives
\(q\geq q_0\) satisfying \eqref{eq:p171-full-D-test}; after this
particular \(q\) and \(C_q\) are fixed, uniform convergence of the finitely
many metric nets gives \(a_i\geq a_{0,i}\) satisfying
\eqref{eq:p171-full-D-metric}.  Hence
\(d=(\epsilon,q,\mathbf a)\in D\) and \(d\geq p_0\).

For \(d\in D\), put
\begin{align*}
 F_d(x)&=\int_{X_x^{\mathrm{an}}}\psi_{q,x}\,
   \operatorname{MA}_x
   ((\phi_{1,a_1})_x,\ldots,(\phi_{n,a_n})_x),\\
 F(x)&=\mathcal I_{\psi,\boldsymbol\phi}(x).
\end{align*}
Each \(F_d\) is continuous by \eqref{eq:p171-full-finite-level}.  Positivity
and common mass control the test-function error, and
\cref{cor:p171-production-family-CLN}, applied only to
\(\alpha_q-\beta_q\), controls the metric error.  Uniformly for every
\(x\in\mathcal M(A)\),
\begin{align}
 |F_d(x)-F(x)|
 &\leq M\|\psi_q-\psi\|_\infty
 +C_q\sum_{i=1}^n
   \|\phi_{i,a_i}-\phi_i\|_\infty\notag\\
 &\leq\epsilon.
 \label{eq:p171-full-product-error}
\end{align}
Thus \(F_d\to F\) uniformly along the cofinal directed set \(D\).  A
uniform limit of a net of continuous real-valued functions is continuous
on an arbitrary topological space.  This proves
\cref{thm:p171-full-assembly}.
\end{proof}


\setcounter{equation}{\savedPseventeenEquation}
\endgroup

\section{What is actually converging}
\label{sec:hybrid-what-converges}

Four measures appear in degeneration problems and should be kept in
separate columns:
\[
\begin{array}{c|c}
\text{input and normalization}&\text{central support and mechanism}\\
\hline
\begin{array}{c}
\text{continuous hybrid psh metric}\\[-1mm]V^{-1}\MA
\end{array}
&
\begin{array}{c}
X_K^{\mathrm{an}}\\[-1mm]\text{hybrid MA continuity}
\end{array}\\[3mm]
\begin{array}{c}
\text{model metric}\\[-1mm]\text{mass }V
\end{array}
&
\begin{array}{c}
\text{divisorial vertices}\\[-1mm]\text{intersection formula}
\end{array}\\[3mm]
\begin{array}{c}
\text{pluricanonical volume form}\\[-1mm]
|t|^{-2\kappa_{\min}}(2\pi\log|t|^{-1})^{-d}
\end{array}
&
\begin{array}{c}
\text{top minimum faces}\\[-1mm]\text{residue calculation}
\end{array}\\[3mm]
\begin{array}{c}
\text{K\"ahler--Einstein metric}\\[-1mm]V^{-1}\MA
\end{array}
&
\begin{array}{c}
\text{canonical-model vertices}\\[-1mm]
\text{hybrid regularity plus MA}
\end{array}.
\end{array}
\tag{66.27}\label{eq:hybrid-four-measure-ledger}
\]
The first row is a continuity theorem for a metric already defined on the
whole hybrid space.  The second is its finite model core.  The third begins
with a volume form and uses a different, dimension-dependent renormalizing
factor.  The fourth requires a geometric regularity theorem before it enters
the first row.  Equality of their supports in a special example does not
identify their coefficients.

\subsection{A model calculation on the nodal chart}
\label{sec:hybrid-node-measures}

Suppose a semistable curve has two components $E_x,E_y$ meeting at the
node \eqref{eq:hybrid-node-chart}, and a relatively nef model line bundle
has degrees $d_x,d_y$.  Its unnormalized central Monge--Amp\`ere measure is
\[
                     d_x\delta_{v_x}+d_y\delta_{v_y},
                 \qquad d_x+d_y=\deg L.
\tag{66.28}\label{eq:hybrid-node-model-measure}
\]
The points $v_x,v_y$ are the endpoints $\alpha=0,1$ of
\eqref{eq:hybrid-node-monomial-valuation}.  By contrast, a normalized
pluricanonical volume form whose weight is constant along the edge can
converge to a multiple of affine length measure on the whole interval.
Both limits live on the same skeleton, but one is atomic and the other is
one-dimensional.  This is the smallest example in which the two entries in
\eqref{eq:hybrid-four-measure-ledger} prevent a false
identification.

The mass test remains decisive.  In the model case
\[
 \int1\,\rd\MA^{\NA}(\phi^{\NA})=d_x+d_y=\deg L,
\tag{66.29}\label{eq:hybrid-node-mass-test}
\]
whereas the pluricanonical residue measure acquires its mass only after the
normalization displayed in the third row of
\eqref{eq:hybrid-four-measure-ledger}.  There is no reason for the two
unnormalized masses to agree.

\paragraph{Source boundary for measure convergence.}
The first row of \eqref{eq:hybrid-four-measure-ledger} is exactly
\cite[Theorem~4.11]{PS23Hybrid}; the third is
\cite{BJ17Hybrid} under the log-smooth and subklt hypotheses stated in
\Cref{thm:hybrid-volume-forms}; and the fourth uses
\cite{PS22HybridKE,PS23Hybrid} under the hypotheses of
\Cref{thm:hybrid-ke}.  The table is a separation ledger, not a new theorem
identifying these measures.

\section{Polynomial escape potentials and hybrid limits}
\label{sec:hybrid-polynomial-escape}
\index{escape potential!hybrid metric}
\index{polynomial automorphism!hybrid limit}
\index{hybrid convergence!common compact support}

Polynomial dynamics gives a useful application of the metric-to-measure
passage. The application also shows why an affine potential, a relative
metric weight and an invariant measure must be distinguished. The
discussion below is an attributed explanation with explicit analytic
inputs. It does not assert that every step in the recent proof of
\cite{Irokawa26RegularV1} has been verified.

Let $f_t:\mathbf C^N\to\mathbf C^N$, $N\geq2$, be a meromorphic
family of nonlinear regular polynomial automorphisms, of fixed forward
and inverse degrees $d_+,d_-\geq2$. Its coefficients are holomorphic
on the punctured disc and have finite pole orders at zero. Assume that
the induced map $f_K$ over $K=\mathbf C((t))$ is regular with these
same degrees. Here regularity means that the indeterminacy loci
$I_+=I(f)$ and $I_-=I(f^{-1})$ of the projective extensions are
disjoint. We use the classical dimension and degree relations
\[
 l_+=\dim I_++1,\qquad l_-=\dim I_-+1,
 \qquad l_++l_-=N,\qquad d_+^{l_-}=d_-^{l_+}.
\tag{66.29a}\label{eq:hybrid-regular-degree-balance}
\]
These are geometric inputs of the regular-automorphism setting,
as recalled in \cite{Kawaguchi13RegularPublished,Lee12RegularPairV2};
they are not statements about affine linear automorphisms.

\paragraph{Balanced approximants and the reference metric.}
Put $F_t=f_t^{l_-}$, $B_t=f_t^{-l_+}$ and
$D=d_+^{l_-}=d_-^{l_+}$. Fix the max norm on affine coordinates and
write $q_n=D^n$. The balanced finite escape potential is
\[
 g_{n,t}(z)=\frac{1}{q_n}
  \log\max\{1,\|F_t^n(z)\|,\|B_t^n(z)\|\}.
\tag{66.29b}\label{eq:hybrid-balanced-escape-approximant}
\]
Use the same formula over $K$, evaluating the polynomial coordinates
at analytic points. When the individual normalized escape sequences
converge locally uniformly, their balanced subsequences give
\[
 G_t=\max\{G_{+,t},G_{-,t}\},\qquad
 G_{\pm,t}=\lim_{m\to\infty}
        d_\pm^{-m}\log^+\|f_t^{\pm m}(z)\|,
 \qquad g_{n,t}\longrightarrow G_t.
\tag{66.29c}\label{eq:hybrid-balanced-escape-limit}
\]
The analogous $K$-limits, their continuity and their precise psh class
are also inputs, rather than consequences of a change of notation.

Homogenize the two coordinate tuples in
\eqref{eq:hybrid-balanced-escape-approximant} to degree $q_n$ and
adjoin $S^{q_n}$, where $S$ is the last projective coordinate.
Check that this combined tuple is nowhere zero on each projective
fibre, including the generic $K$-fibre. It then defines a morphism
to projective space and, after division by $q_n$, a semipositive
metric on $\mathcal O(1)$. Individual forward maps need not extend
as projective morphisms. It is the combined tuple that matters.
The regular-pair projectivization is the geometric construction used
in \cite{Lee12RegularPairV2,Irokawa23HenonV2,Irokawa26RegularV1}.

For an exact formula choose the continuous max reference metric,
whose local potential in $S=1$ is
\[
 \rho_t(z)=\log\max\{1,\|z\|\},\qquad
 \rho_K(z)=\log\max\{1,\|z\|_K\}.
\tag{66.29d}\label{eq:hybrid-escape-reference}
\]
Both are ordinary-norm weights, in the convention of this chapter.
The globally defined relative weights are
\[
 \varphi_{n,t}=g_{n,t}-\rho_t,\qquad
 \varphi_{n,K}=g_{n,K}-\rho_K.
\tag{66.29e}\label{eq:hybrid-escape-relative-weight}
\]
Equivalently, in homogeneous coordinates $X$, take
$q_n^{-1}\log\max_j|s_{j,n}(X)|$ and subtract the reference
$\log\max_i|X_i|$.
The frame factors cancel; on the affine chart the reference term
does not disappear. A smooth Fubini--Study reference instead gives
$\varphi_{n,t}=g_{n,t}-\rho_{\mathrm{FS},t}$, with
$\rho_{\mathrm{FS},t}=\tfrac12\log(1+\sum_i|z_i|^2)$.
Changing reference changes the relative weight and its background
curvature together, while preserving the absolute metric.

In particular, if $\Theta_t=dd^c\rho_t$ is the reference curvature,
then on the affine chart
\[
 \Theta_t+dd^c\varphi_{n,t}=dd^c g_{n,t},\qquad
 \Theta_t+dd^c\varphi_t=dd^c G_t,
 \quad \varphi_t=G_t-\rho_t.
\tag{66.29f}\label{eq:hybrid-escape-curvature-reference}
\]
The same relation is used centrally in the applicable continuous
metric and current formalism, with $\varphi_K=G_K-\rho_K$.
Neither $dd^c\varphi_t$ alone nor
a raw degree-$q_n$ model function is the normalized escape curvature.
Special-fibre multiplicities in a model computation remain those of
\eqref{eq:hybrid-model-ma}. Semipositivity is supplied by the
globally generated tuple or a nef model, not by the mere existence
of a uniform limit of arbitrary model metrics.

\paragraph{The uniform estimate needed before taking measures.}
The finite tuples give continuous hybrid metrics after the usual
model modifications. Relative to the hybrid max reference, their
weights on a nonzero fibre are
$\Phi_n=\lambda(t)\varphi_{n,t}$, with central value
$\varphi_{n,K}$ and the positive $\lambda(t)$ of
\eqref{eq:hybrid-lambda}. A sufficient approximation input is
\[
 \|\varphi_{n,t}-\varphi_t\|_{\mathbf P^N}
       \leq e_n(1+\log|t|^{-1}),\qquad
 \|\varphi_{n,K}-\varphi_K\|_{\mathbf P^{N,\mathrm{an}}_K}
       \leq e_n,\qquad e_n\longrightarrow0.
\tag{66.29g}\label{eq:hybrid-escape-uniformity-input}
\]
The estimates hold for all sufficiently small nonzero $t$ in one
fixed disc. Since $\lambda(t)\log|t|^{-1}=-\log r$, they give
uniform hybrid convergence of $\Phi_n$ to $\Phi$, hence a continuous
semipositive hybrid metric. This is the required order of quantifiers:
choose the approximation index before taking the fibre limit.

The continuous projective metric theorem
\Cref{thm:hybrid-ma-continuity} now gives probability measures
\[
 \mu_t^{\mathrm{proj}}=\MA(\Phi|_{X_t}),\qquad
 \mu_0^{\mathrm{proj}}=\MA^{\NA}(\varphi_K),\qquad
 \mu_t^{\mathrm{proj}}\longrightarrow\mu_0^{\mathrm{proj}}.
\tag{66.29h}\label{eq:hybrid-escape-projective-measures}
\]
Here $(\mathcal O(1)^N)=1$; on complex fibres the powered absolute
value is undone before taking ordinary complex curvature, as in
\Cref{thm:hybrid-ma-continuity}. Thus the affine restriction is
$(dd^cG_t)^N$, without an extra $\lambda(t)^N$ in its probability
mass. The model atoms retain their degrees and multiplicities before
the same volume normalization.

The uniformity input must actually be proved for the family under
consideration. The constants in the published local escape estimates
of \cite{Kawaguchi13RegularPublished} control Nullstellensatz
certificates as well as map coefficients; complex estimates also
count monomials. The product-pair method in
\cite{Lee12RegularPairV2} retains forward invariance of the selected
domain when it telescopes. Membership of the diagonal at the first
step does not establish membership along its product-map orbit.
These requirements remain explicit here. The covering constants
used in \cite[Proposition~5.2]{Irokawa26RegularV1} fail the source
check recorded in \Cref{sec:provenance-hybrid-polynomial}; this
discussion assumes \eqref{eq:hybrid-escape-uniformity-input} and
does not certify that proof or a complete replacement.

\paragraph{Returning to the affine space.}
Suppose all the limiting fibre measures in
\eqref{eq:hybrid-escape-projective-measures} are supported in one
closed hybrid polydisc $W$ contained in the affine chart, over the
fixed closed base disc. This is a separate support input. For example,
complex coordinate bounds $|z_i|\leq C|t|^{-A}$ give
\[
 |z_i|_{\mathrm{hyb}}
       \leq C^{\lambda(t)}r^{-A}
       \leq\max\{1,C\}r^{-A}.
\tag{66.29i}\label{eq:hybrid-escape-support-bound}
\]
The central support needs the corresponding $K$-bound. The resulting
closed polydisc is compact in the projective hybrid space and avoids
infinity. The affine hybrid projection itself is not proper.

For a continuous affine hybrid test $h$, choose a continuous cutoff
$\chi$ supported in the affine chart and equal to one on $W$.
The extension by zero of $\chi h$ is a continuous projective test.
Applying \eqref{eq:hybrid-escape-projective-measures} to this test
gives the affine weak limit. This argument works with nets and does
not require first countability. If common support is unavailable,
an appropriate tightness argument must replace it; a bare restriction
of projective weak convergence is insufficient.

\paragraph{Identifying the invariant measure is another input.}
The functional equations of the escape potentials give
$f^*T_+=d_+T_+$ and $(f^{-1})^*T_-=d_-T_-$, where
$T_\pm=dd^cG_\pm$ has the positive orientation. When the products
are defined, \eqref{eq:hybrid-regular-degree-balance} explains the
invariance factor $d_+^{l_-}/d_-^{l_+}=1$ for
$T_+^{l_-}\wedge T_-^{l_+}$.

To identify this product with the maximum-potential measure requires
the escape-rank and maximum/product interfaces. The complex argument
in \cite[Lemma~2.12]{GuedjSibony02Author} uses nonnegative continuous
psh $u,v$ and the ranks
\[
\begin{gathered}
 (dd^cu)^{l_-}=0\ \text{on }\{u>0\},\qquad
 (dd^cv)^{l_+}=0\ \text{on }\{v>0\},\\
 (dd^c\max\{u,v\})^N=(dd^cu)^{l_-}\wedge(dd^cv)^{l_+}.
\end{gathered}
\tag{66.29j}\label{eq:hybrid-escape-rank-identification}
\]
Its truncation, locality and local uniform limiting operations explain
the desired identification. In the non-Archimedean setting the
continuous locally approachable product formalism of \cite{CLD25}
supplies the selected operations in its specified class. One must
still verify that the dynamical potentials belong to that class and
that the required ranks and identification hold. A nonzero tuple
of escape coordinates does not make one fixed coordinate invertible
on an entire compact set. Local choices and the eventual inactivity
of the constant in a truncated maximum cannot be omitted.

The older \cite[Remark~2.6]{Irokawa23HenonV2} leaves the central
maximum/product equality open; \cite[Theorem~4.4]{Irokawa26RegularV1}
claims to establish it in higher dimensions. The present application
retains its non-Archimedean identification as an explicit input.
Neither historical wording nor a theorem assertion changes the
book's canonical problem statuses without verification.

\paragraph{A scalar exponent check and the cocycle boundary.}
Weak convergence of invariant measures is not by itself convergence
of Lyapunov exponents. For the H\'enon family
$H_t(x,y)=(p_t(x)-a(t)y,x)$, the determinant is the scalar $a(t)$.
If $a(t)=t^m b(t)$ with $b(0)\neq0$, the sum of the complex
Lyapunov exponents, when they are defined with their usual
integrability hypotheses, is the direct determinant identity
\[
 \lambda_{1,t}+\lambda_{2,t}
   =\log|a(t)|
   =\frac{\log|t|^{-1}}{-\log r}\log|a|_K+O(1).
\tag{66.29k}\label{eq:hybrid-henon-total-exponent}
\]
For $a(t)=t$ the sum is $\log|t|<0$. This fixes the sign without
an appeal to a measure limit. The first exponent concerns the
derivative cocycle $DH_t^n$, a product evaluated along the orbit;
one-step derivative eigenvalues cannot replace it. Interchanging
the parameter limit, the long-orbit limit and integration needs
additional estimates. No first-exponent asymptotic is imported here.

\section{Sequential rational dynamics and relative projective lines}
\label{sec:hybrid-sequential-rational}
\index{rational map!sequential hybridization}
\index{Robinson field!dynamical degeneration}
\index{Lyapunov exponent!singular integrability}

A meromorphic one-parameter family has a specified central field.
An arbitrary degenerating sequence of rational maps requires a different
base. The construction of \cite{FavreGong26SequentialV4} records the
sequence in a Banach ring and places its ultrafilter limits over the
Stone--\v{C}ech compactification of the integers. This construction
does not identify all such limits with one map over $\mathbf C((t))$.

Fix $d\geq2$. For a degree-$d$ complex rational map let $R(f)$ be
the absolute resultant of a homogeneous lift whose largest coefficient
has absolute value one. Optimize this value over conjugacy and write
$R_n=\sup_{M\in\mathrm{PGL}_2(\mathbf C)}R(Mf_nM^{-1})$.
Suppose the classes $[f_n]$ leave every compact subset of the degree-$d$
moduli space. In the source construction $R_n\to0$, and a suitable
constant $C_d$ gives the positive scales
\[
 \varepsilon_n=\frac{1}{\log(C_d/R_n)}\in(0,1],\qquad
 \varepsilon_n\longrightarrow0,\qquad
 R_n^{\varepsilon_n}=e^{-1}C_d^{\varepsilon_n}.
\tag{66.29l}\label{eq:hybrid-sequential-resultant-scale}
\]
The last equality is an immediate logarithmic identity. In particular,
the optimized resultant has a nonzero boundary value after rescaling.
Properness of the optimized resultant on the degree-$d$ moduli
space is a source input; the fixed-field resultant-minimum theorem of
\cite{Rumely13ResultantV2} concerns a separate finite-tree minimization.
It does not make arbitrary bounded subsets of a strongly metrized
Berkovich line compact.

Define the complete Banach ring
\[
 A_{\varepsilon}=\left\{(a_n)\in\mathbf C^{\mathbf N}:
                  \sup_n|a_n|_\infty^{\varepsilon_n}<\infty\right\},
 \qquad \|(a_n)\|=\sup_n|a_n|_\infty^{\varepsilon_n}.
\tag{66.29m}\label{eq:hybrid-sequential-Banach-ring}
\]
Its idempotents include the characteristic functions of subsets of
$\mathbf N$, so it has zero divisors. The integral-base hypothesis of
\Cref{thm:p171-full-assembly} therefore does not cover this ring.
The source identifies $\mathcal M(A_{\varepsilon})$ with
$\beta\mathbf N$ and constructs a compact relative analytic line
\[
 \pi:\mathbf P^{1,\mathrm{an}}_{A_{\varepsilon}}
                         \longrightarrow\beta\mathbf N.
\tag{66.29n}\label{eq:hybrid-sequential-relative-line}
\]
At a principal point $n$ the field is
$\mathbf C^{\varepsilon_n}=(\mathbf C,|\cdot|_\infty^{\varepsilon_n})$.
At a nonprincipal ultrafilter $\omega$ it is a Robinson field
$\mathcal H(\omega)$, with the quotient norm given by the ultrafilter
limit of the powered norms. The source-attested boundary fields are
algebraically closed and spherically complete, with characteristic and
residue characteristic zero and value group $\mathbf R_{>0}$.
These are properties of this construction, rather than properties
of arbitrary Banach-ring residue fields.

Choose optimized normalized lifts $[P_n:Q_n]$. The coefficient sequences
belong to $A_{\varepsilon}$ and the resultant sequence is a unit, by
\eqref{eq:hybrid-sequential-resultant-scale}. They consequently define
a degree-$d$ relative map $F$, with fibre maps $f_\omega$.
The optimized boundary maps have no potential good reduction in
\cite[Theorem~3.17]{FavreGong26SequentialV4}. The word ``potential''
matters: a Julia set can be a singleton at a non-Gauss point after
a change of coordinates. Different nonprincipal ultrafilters can
give different fields and maps; an ultrafilter conclusion is not
an assertion that the original sequence has a unique ordinary limit.

\paragraph{Uniform potentials before weak convergence.}
The following estimate explains the measure-continuity mechanism,
while keeping its relative analytic inputs visible. Let $N_\omega$
be the homogeneous reference norm: the powered Euclidean norm on
a complex fibre and the max norm on a non-Archimedean fibre. For a
nowhere-zero homogeneous lift $\widetilde F_\omega$ set
\[
 g_{0,\omega}([z])=
       \frac1d\log N_\omega(\widetilde F_\omega(z))
                          -\log N_\omega(z),\qquad
 G_{k,\omega}=\sum_{j=0}^{k-1}d^{-j}g_{0,\omega}\circ f_\omega^j.
\tag{66.29o}\label{eq:hybrid-sequential-Green-sum}
\]
Homogeneity makes the relative weight $g_0$ independent of the
representative of $[z]$. Assume the reference metric and this weight
are continuous on the compact relative line, with $|g_0|\leq M$.
The sum converges uniformly, with
\[
 \|G_{\infty,\omega}-G_{k,\omega}\|_\infty
                 \leq\frac{M}{1-d^{-1}}d^{-k}
                  \quad\text{for every }\omega.
\tag{66.29p}\label{eq:hybrid-sequential-Green-tail}
\]

Write $\Delta_\omega$ for the positively oriented fibre Laplacian,
so that logarithmic zeros have positive mass. On the ordinary complex
line $\Delta_\infty=2dd^c$ in the book's convention; when its input is
an intrinsic $\varepsilon_n$-scaled weight, the operator divides by
$\varepsilon_n$. Thus the fibre probability curvature is unchanged
by powering the absolute value. If $\mu_{\mathrm{can},\omega}$
is the reference curvature probability, the normalized pullbacks are
\[
 \mu_{k,\omega}=d^{-k}(f_\omega^k)^*\mu_{\mathrm{can},\omega}
      =\mu_{\mathrm{can},\omega}+\Delta_\omega G_{k,\omega},
 \qquad
 \mu_{f_\omega}=\mu_{\mathrm{can},\omega}
                                   +\Delta_\omega G_{\infty,\omega}.
\tag{66.29q}\label{eq:hybrid-sequential-pullback-probabilities}
\]
For the relative passage one needs continuity of the reference measures
and finite weighted traces, together with a dense family of continuous
tests $h$ satisfying $\|\Delta_\omega h\|_{\mathrm{TV}}\leq C_h$
uniformly in $\omega$. Integration by parts then gives
\[
 \left|\int h\,\rd\mu_{f_\omega}
                  -\int h\,\rd\mu_{k,\omega}\right|
       \leq C_h\frac{M}{1-d^{-1}}d^{-k}.
\tag{66.29r}\label{eq:hybrid-sequential-test-tail}
\]
Choose $k$ before taking a fibre limit. Continuity for finite pullbacks,
followed by this estimate and uniform approximation of continuous
tests, gives weak continuity for the equilibrium measures. Basepoint-free
model tests, the reference curvature and their uniform Laplacian bounds
are actual inputs. Bounded families of continuous functions need not
be compact, and density for arbitrary Banach rings does not follow
from the model-function assertion used here.

\cite[Theorem~5.8]{FavreGong26SequentialV4} asserts measure continuity
over $A_\varepsilon$ and additionally continuity of certain singular
quasi-model integrals. The exact-v4 proof has the singular-passage
obligation recorded in \Cref{sec:provenance-hybrid-sequential}; this
exposition gives the uniform-potential mechanism and does not certify
a complete repair of that additional assertion.

\paragraph{The distinct good-base relative-line theorem.}
For a good Banach ring $A$, an analytic base $Y$ and a finite flat
endomorphism $\varphi$ of $Y\times_A\mathbf P^{1,\mathrm{an}}_A$
of degree $d\geq2$, polarized by $\mathcal O(1)$, the equilibrium
measure family is weakly continuous. This is
\cite[Theorem~A and Theorem~6.26]{Poi25DynamicsZV5}.
Its definition of a good ring includes local arc connectedness of
the spectrum and the specified non-Archimedean part. It does not
automatically include the product ring $A_\varepsilon$, whose spectrum
$\beta\mathbf N$ is totally disconnected and has nondiscrete
nonprincipal points.

Poineau's proof uses a continuous relative potential, locally uniform
metric approximation and dense macologue tests. Its test approximation
has the form $2\eta+\|\rho_k\|_\infty M_\eta$: fix the continuous-test
error $\eta$, then make the potential tail $\rho_k$ small. The test
class permits its specified rational constants and is distinct from
the model class above. The good-base theorem, the sequential
construction and the integral-base theorem in this chapter have
separate domains; none extends the others merely because all three
use Banach rings. No higher-dimensional or energy-pairing conclusion
is imported from the relative-line theorem.

\paragraph{Lyapunov exponents require the singular integral.}
For $f_\omega=[P:Q]$ put
$J=P_{z_0}Q_{z_1}-P_{z_1}Q_{z_0}$. The intrinsic chordal derivative is
\[
 \|df_\omega\|([z])=
       |J(z)/d|_\omega\,
       \frac{N_\omega(z)^2}{N_\omega(P(z),Q(z))^2},
 \qquad
 \chi(f_\omega)=\int\log\|df_\omega\|\,\rd\mu_{f_\omega}.
\tag{66.29s}\label{eq:hybrid-sequential-valued-derivative}
\]
The scalar $1/d$ is inside the valued norm. On
$\mathbf C^{\varepsilon_n}$ this gives the usual spherical derivative
raised to $\varepsilon_n$, and hence
$\chi(f_n\text{ over }\mathbf C^{\varepsilon_n})
=\varepsilon_n\chi_\infty(f_n)$.
On the boundary Robinson field $|d|=1$, whereas characteristic zero
alone would not imply this for an arbitrary non-Archimedean field.
For $f(z)=z^d$, the principal values are $\varepsilon_n\log d$
and the Gauss-fibre value is zero.

The logarithm in \eqref{eq:hybrid-sequential-valued-derivative}
has critical-point singularities. A sufficient additional interface
is continuity of the truncated functions
$\ell_A=\max\{\log\|df\|,-A\}$ on the relative space, together with
\[
 \lim_{A\to\infty}\sup_\omega
       \int\bigl(-\log\|df_\omega\|-A\bigr)_+
                                         \,\rd\mu_{f_\omega}=0.
\tag{66.29t}\label{eq:hybrid-sequential-critical-uniform-integrability}
\]
The derivative is bounded above on the compact relative space.
Truncation, weak continuity and this tail estimate then justify
continuity of the integral. Nonzero fibrewise critical sections and
the continuous reference corrections are retained in this step.
Weak convergence alone does not justify it.

\cite[Corollary~5.9]{FavreGong26SequentialV4} asserts the intrinsic
continuity, which would give
$\lim_\omega\varepsilon_n\chi_\infty(f_n)=\chi(f_\omega)$.
The statement is source-attested; the exact proof is subject to the
singular-integral boundary just described. The corrected scalar
formula is not a counterexample to intrinsic continuity, nor does
absence of potential good reduction imply a strictly positive
non-Archimedean Lyapunov exponent here.

The tree comparison of \cite{Luo21LimitsV4,Luo21TreesV2} is also
compatible with the sequential viewpoint. Its asymptotic-cone quotient
identifies $(x_n)$ with $(y_n)$ when their distance divided by $r_n$
tends to zero along $\omega$. Unscaled distance is not the quotient
criterion. Selected correspondence
and branched-covering results are source-attested; a complete audit
of both scale-comparison inequalities is not imported.

\section{Hyperbolic K3 eigen-currents and the hybrid interface}
\label{sec:hybrid-K3-eigen-currents}
\index{K3 surface!hybrid invariant measure}
\index{eigen-current!symmetric normalization}

K3 dynamics uses two reciprocal eigenclasses and a mixed product,
rather than a single ample equilibrium metric. This distinction
determines both the normalization and the scope of a hybrid application.
Let $X$ be a complex projective K3 surface and $f$ an automorphism
of positive entropy, with spectral radius $\Lambda>1$ on
$N^1(X)_{\mathbf R}$. There are nef classes $\alpha_+,\alpha_-$ with
\[
 f^*\alpha_+=\Lambda\alpha_+,\qquad
 f^*\alpha_-=\Lambda^{-1}\alpha_-,\qquad
 \alpha_+^2=\alpha_-^2=0,\qquad
 q=\alpha_+\cdot\alpha_->0.
\tag{66.29u}\label{eq:hybrid-K3-eigenclasses}
\]
The classes lie in the real Picard space. This space need not equal
all of $H^{1,1}(X,\mathbf R)$.

\paragraph{The fixed-fibre normalization.}
Write an ample class as
\[
 H=a\alpha_++b\alpha_-+w,\qquad
 a=\frac{H\cdot\alpha_-}{q}>0,\quad
 b=\frac{H\cdot\alpha_+}{q}>0,\quad
 w\perp\alpha_+,\alpha_-.
\tag{66.29v}\label{eq:hybrid-K3-ample-decomposition}
\]
The intersection form is negative definite on this orthogonal
complement, and the powers of $f^*$ are bounded there. For a smooth
positive curvature form $\omega_H$ representing $H$, the correct
approximants in both directions are
\[
 T_{+,n}=\Lambda^{-n}(f^n)^*\omega_H,\qquad
 T_{-,n}=\Lambda^{-n}(f^{-n})^*\omega_H.
\tag{66.29w}\label{eq:hybrid-K3-symmetric-approximants}
\]
They converge to $aT_+$ and $bT_-$, where $[T_\pm]=\alpha_\pm$.
The existence and uniqueness of the complex eigen-currents use
\cite[Theorem~2.4]{Cantat01K3Published}; their continuous local
potentials use its Theorem~3.1. Locally uniform potential convergence
gives the corresponding Bedford--Taylor mixed-product limit. Curvature
forms here use squared-norm weights with the book's $dd^c$; equivalently,
ordinary-norm weights use $2dd^c$. The conventions are not mixed.

The intersection computation gives
\[
 M_n:=\int_X T_{+,n}\wedge T_{-,n}
       =abq(1+\Lambda^{-4n})+O(\Lambda^{-2n})
       \longrightarrow abq>0.
\tag{66.29x}\label{eq:hybrid-K3-approximant-mass}
\]
Without the factor $\Lambda^{-2n}$ the raw mixed-pullback mass is
$abq(\Lambda^{2n}+\Lambda^{-2n})+O(1)$ and diverges. A tensor-power
factor $m^{-2}$ removes the choice of $H^{\otimes m}$; it cannot
replace dynamical normalization. The probability measures are
\[
 \mu_n=M_n^{-1}T_{+,n}\wedge T_{-,n},\qquad
 \mu_f=q^{-1}T_+\wedge T_-,\qquad
 \mu_n\longrightarrow\mu_f.
\tag{66.29y}\label{eq:hybrid-K3-invariant-probability}
\]
Their limiting invariance is explained by the product of eigenvalues
$\Lambda\Lambda^{-1}=1$. Rescaling either eigenclass changes its
current and $q$ together, so the probability does not depend on
that choice.

For a complete non-Archimedean field of characteristic and residue
characteristic zero, \cite[Theorems~5.6.4, 5.6.5 and~5.6.7]{Filip19K3V1}
constructs the corresponding homogeneous torsor potentials and positive
eigen-currents in its projective K3 setting.
Its Corollary~5.6.8 gives the nonzero invariant mixed measure and
explicitly retains $\Lambda^{-n}$ in each direction.
The case $K=\mathbf C((t))$ satisfies the characteristic assumptions.
These are known fixed-field results in their specified potential and
product class. Existence of arbitrary real psh functions alone would
not establish that a mixed Chambert--Loir--Ducros product is available.

\paragraph{What a family-uniform hybrid argument still needs.}
Consider a meromorphic projective K3 family with a relative
automorphism and a fixed hyperbolic action on a marked relative Picard
space. After any required finite base change, suppose there is a
common relative Picard torsor, fixed $\Lambda$ and compatible reference
potentials. The normalized finite iterates have changing real
homogeneities; the limiting eigenpotentials have the limiting characters.
Compare their local weights only after tracking these homogeneities
and their respective reference terms on compact model or torsor domains.

A sufficient family approximation input, for these relative weights,
is
\[
 \|\psi_{\pm,n,t}-\psi_{\pm,t}\|
         \leq e_n(1+\log|t|^{-1}),\qquad
 \|\psi_{\pm,n,K}-\psi_{\pm,K}\|\leq e_n,
 \qquad e_n\longrightarrow0,
\tag{66.29z}\label{eq:hybrid-K3-uniformity-input}
\]
uniformly on the chosen finite compact cover and for all sufficiently
small nonzero $t$ in one disc. The reference potentials must themselves
have compatible hybrid limits. Multiplication of the ordinary complex
weights by the positive factor $\lambda(t)=\log r/\log|t|$ then makes
the potential tails uniformly small in the hybrid space. A global
bound on the entire noncompact torsor is not needed and generally
does not follow from the fixed-field contraction theorem.

There is a second input: the limiting real nef metrics must lie in
a compatible continuous semipositive class with a mixed-curvature
continuity theorem and uniform bounds for its test integrals.
\Cref{thm:hybrid-ma-continuity} directly treats a relatively ample
line bundle. Applying it to these real nef eigenclasses requires a
justified polarization or ample-perturbation passage, with its
uniformity preserved. A formal expansion of powers is not a substitute
for that passage. In local estimates keep the reference curvature,
all cutoff derivatives and special-fibre multiplicities; these are
the quantities used by the mixed-product integration-by-parts bounds.

Under those two inputs and compatible positive mass normalization,
mixed-curvature continuity gives the conditional hybrid limit
\[
 \mu_{f_t}=q_t^{-1}T_{+,t}\wedge T_{-,t}
                 \longrightarrow
 \mu_{f_K}=q_K^{-1}T_{+,K}\wedge T_{-,K}.
\tag{66.29aa}\label{eq:hybrid-K3-conditional-limit}
\]
Here the classes and products have the normalization of
\eqref{eq:hybrid-K3-eigenclasses}, and $q_t\to q_K>0$ is part of
the compatible family intersection data. The complex measures use
ordinary complex curvature after undoing the powered norm.
Choose the approximation index first, take the fibre limit second,
and only then remove the uniform tail.

\cite[Theorem~1.1]{Irokawa24K3V1} asserts this hybrid measure limit.
Its exact-v1 backward and mixed approximant formulas fail the
fixed-fibre test \eqref{eq:hybrid-K3-approximant-mass}; its family
uniformity and product interfaces have the separate obligations
recorded in \Cref{sec:provenance-hybrid-K3}. The correctly defined main
limit is not thereby refuted. The discussion above incorporates
the known symmetric normalization and states an explicit conditional
hybrid route, without certifying a complete repair of that proof.

\subsection{Dense existence in hyper-K\"ahler families}
\label{sec:hybrid-HK-dense-entropy}
\index{hyper-K\"ahler manifold!positive entropy in families}
\index{K3 surface!dense existence of positive entropy}

The existence of a hyperbolic map on many fibres raises a separate
question from the convergence of measures for a specified relative map.
For a projective irreducible holomorphic symplectic manifold $Y$
and a birational self-map $g$, the recent source
\cite{DenisiEtAl26EntropyV1} uses the H2-level algebraic entropy
\[
 h_2(g):=\log\lambda_1(g),\qquad
 \lambda_1(g)=
 \rho\bigl(g^*|_{H^{1,1}(Y,\mathbf C)}\bigr).
\tag{66.29ab}\label{eq:hybrid-HK-entropy-convention}
\]
Here $\rho$ denotes spectral radius. The H2 action is functorial
under powers because a birational hyper-K\"ahler self-map is an
isomorphism in codimension one. The convention in
\eqref{eq:hybrid-HK-entropy-convention} concerns the first dynamical
degree; it is not a definition of full topological entropy in every
dimension. For a K3 automorphism it is the positive-entropy condition
used in \eqref{eq:hybrid-K3-eigenclasses}.

Theorem~1.1 of the exact eight-page arXiv:2610.05173v1 reports
the following dense existence statement. Let
$\pi:\mathcal Y\to\Delta$ be a non-isotrivial polarized smooth
family of irreducible holomorphic symplectic manifolds over a small
complex disc, with $b_2(Y_t)\geq6$. Then there is a Euclidean-dense
subset $D\subset\Delta$ such that
\[
 \text{for every }t\in D\text{ there is }
 g_t\in\operatorname{Bir}(Y_t)
 \text{ with }\lambda_1(g_t)>1.
\tag{66.29ac}\label{eq:hybrid-HK-dense-existence}
\]
This is a source-attested statement with the proof qualifications
in \Cref{sec:provenance-HK-group-entropy}. Its Corollary~1.3
reports the analogous Euclidean density over an irreducible complex
analytic base when the fibres are projective. Projectivity and
non-isotriviality remain hypotheses; the statement does not extend
to arbitrary strict Calabi--Yau families.

Two features of the argument explain the scope of this information.
The common subgroup in
\cite[Theorem~1.1]{DenisiEtAl26GroupsV1} is realized on fibres
outside a finite exceptional set, and its finite index is asserted
in the very general group. A positive-entropy element in that group
can be replaced by a power lying in the common subgroup, since
$\lambda_1(g^m)=\lambda_1(g)^m$. This produces compatible lattice
actions on the stated fibres; it does not produce a morphism of
the whole relative space. The other input is strong intersection
density of suitable period divisors, as formulated in
\cite[Definition~2.8 and Theorem~2.10]{AmerikSoldatenkovVerbitsky26V1}.
Ambient density alone would not ensure dense intersection with a
given nonconstant period disc. The fixed deformation type and
finite exceptional strata used by the lattice argument are retained.

On K3 surfaces the birational maps are regular, so
\eqref{eq:hybrid-HK-dense-existence} supplies many fibres with
the fixed-fibre eigen-currents discussed above. The map $g_t$ may
depend on $t$, and neither its dynamical degree nor its eigenclasses
are specified by that assertion. In particular, it does not supply
one relative automorphism, a common Picard torsor, the uniform
potential estimate \eqref{eq:hybrid-K3-uniformity-input}, or a
mixed-curvature continuity theorem. Those are the compatibility
data required for \eqref{eq:hybrid-K3-conditional-limit}.
Dense existence of maps therefore adds geometric context to the
hybrid problem while leaving its family approximation inputs open.


\section{Skeleta and canonical measures}
\label{sec:hybrid-skeleta}

Assume now that $\mathcal X$ is an SNC model. Write
\[
                         \mathcal X_0=\sum_{i\in I}b_iE_i.
\tag{66.30}\label{eq:hybrid-snc-fibre}
\]
The dual complex $\Delta_{\mathcal X}$ embeds in
$X_K^{\mathrm{an}}$. Each connected component $Y$ of a nonempty intersection
$E_J:=\bigcap_{j\in J}E_j$ gives a face parametrizing monomial valuations with barycentric coordinates $(w_j)_{j\in J}$ satisfying
\[
                         w_j\geq0,
             \qquad \sum_{j\in J}b_jw_j=1.
\tag{66.31}\label{eq:hybrid-simplex}
\]

\begin{definition}[Skeleton and retraction]
\label{def:hybrid-skeleton}
The skeleton of the SNC model is
\[
                     \Sk(\mathcal X):=\Delta_{\mathcal X}
                     \subset X_K^{\mathrm{an}}.
\tag{66.32}\label{eq:hybrid-skeleton}
\]
The monomial retraction
\[
               \rho_{\mathcal X}:X_K^{\mathrm{an}}
                    \longrightarrow\Sk(\mathcal X)
\tag{66.33}\label{eq:hybrid-retraction}
\]
records the values of the vertical components $E_i$.
\end{definition}

A model Monge--Amp\`ere measure is atomic at vertices. Canonical
degenerating volume forms can instead limit to Lebesgue measure on
higher-dimensional faces. These are different constructions.

\begin{definition}[Weight-minimum skeleton]
\label{def:hybrid-essential-skeleton}
For a nonzero pluricanonical form $\Omega$, its model-independent weight function has the Kontsevich--Soibelman minimum locus
\[
 \Sk(X_K,\Omega)=\{x:\mathrm{wt}_\Omega(x)=\min\mathrm{wt}_\Omega\}.
\tag{66.34}\label{eq:hybrid-essential-skeleton}
\]
The essential skeleton is the union of these loci over all nonzero pluricanonical forms, as in \Cref{def:skeleton-essential}.
\end{definition}

The weight is rational piecewise affine on an SNC skeleton; horizontal
zeroes can produce bends. On a model adapted to its divisor it is affine
on each relevant face. A minimum locus of one form must not be identified
with the union without an additional geometric hypothesis.

\begin{theorem}\label{thm:hybrid-volume-forms}
Let $(X,\mathcal B)\to\Delta^*$ be a projective meromorphic log-smooth pair: $X\to\Delta^*$ is smooth, the support of $\mathcal B$ has relative simple normal crossings, and every coefficient of $\mathcal B$ is less than one.  Let $\psi$ be a continuous metric on the relative $\mathbf Q$-line bundle $K_{(X,\mathcal B)/\Delta^*}$, with analytic singularities at the origin.  Choose a log-smooth model $(\mathcal X,\overline{\mathcal B})\to\Delta$ and a $\mathbf Q$-line bundle $\mathcal L$ extending $K_{(X,\mathcal B)/\Delta^*}$ such that $\psi$ extends continuously to $\mathcal L$, and write
\[
 \mathcal X_0=\sum_i b_iE_i,\qquad
 K^{\log}_{(\mathcal X,\overline{\mathcal B})/\Delta}
       =\mathcal L+\sum_i a_iE_i,
\]
and let $\Delta(\mathcal L)$ be the subcomplex of $\Sk(\mathcal X)$ spanned by the vertices corresponding to the indices where $a_i/b_i$ is minimal.  Put
\[
 \kappa_i=\frac{a_i}{b_i},\qquad
 \kappa_{\min}=\min_i\kappa_i,\qquad
 d=\dim\Delta(\mathcal L),
\]
and let
\[
 \nu_t=\rme^{2(\psi_t-\phi_{\mathcal B_t})}
\]
be the fibrewise volume measure, where $\phi_{\mathcal B_t}$ is the divisor metric of $\mathcal B_t$.  When $\mathcal B=0$ and $\psi$ is induced by a nowhere-vanishing $m$-pluricanonical form $\Omega$, this is $\nu_t=|\Omega_t|^{2/m}$, and $\Delta(\mathcal L)$ is the weight-minimum skeleton $\Sk(X_K,\Omega)$.  Then
\[
 \mu_t:=
 \frac{\nu_t}
 {|t|^{2\kappa_{\min}}
  (2\pi\log|t|^{-1})^d}
 \ \longrightarrow\
 \mu_0:=
 \sum_{\dim\sigma=d}
 \left(\int_{Y_\sigma}\operatorname{Res}_{Y_\sigma}(\psi)\right)
 b_\sigma^{-1}\lambda_\sigma
\]
weakly on the hybrid space. The residual coefficient is defined by adjunction: for $Y=Y_\sigma$ cut out by the components indexed by $J$,
\[
\begin{aligned}
 B_Y^{\mathcal L}&=\overline{\mathcal B}|_Y+
 \sum_{i\notin J}\bigl(1-(a_i-\kappa_{\min}b_i)\bigr)E_i|_Y,\\
 \operatorname{Res}_Y(\psi)
 &=\rme^{2(\psi|_Y-\phi_{B_Y^{\mathcal L}})}.
\end{aligned}
\tag{66.34a}\label{eq:hybrid-residual-boundary}
\]
On a top minimum face this is a finite subklt volume measure. Here $\sigma$ runs over the top-dimensional faces of $\Delta(\mathcal L)$, $Y_\sigma$ is the corresponding stratum, $b_\sigma=\gcd\{b_i:E_i\supset Y_\sigma\}$, and $\lambda_\sigma$ is integral affine Lebesgue measure.
\end{theorem}

\begin{proof}
On the adapted simplex, the affine weight is
$\kappa(w)=\sum_i a_iw_i=\sum_i\kappa_i(b_iw_i)$,
where $\sum_i b_iw_i=1$.  Its minimum locus is therefore exactly the
subcomplex spanned by the vertices with
$\kappa_i=\kappa_{\min}$, namely $\Delta(\mathcal L)$.

We now compute on one coordinate chart meeting the stratum
$E_0\cap\cdots\cap E_p$.  After absorbing a unit into one coordinate,
\[
                 t=z_0^{b_0}\cdots z_p^{b_p}.
\]
After increasing the coordinate system if necessary, the horizontal
components of $\overline{\mathcal B}$ through the chart are
$B_j=\{z_j=0\}$, $p<j\leq m$, with coefficients $c_j<1$.
Order the vertical components so that
\[
 \kappa_0=\cdots=\kappa_q
       =\min_{0\leq i\leq p}\kappa_i,
 \qquad
 \kappa_i>\kappa_0\quad(q<i\leq p).
\]
The logarithmic
coordinates
\[
              w_i=\frac{\log|z_i|}{\log|t|}
\]
lie in $\sum b_iw_i=1$.  The relative logarithmic canonical form is
obtained by contracting
\[
 \frac{\rd z_0}{z_0}\wedge\cdots\wedge
 \frac{\rd z_p}{z_p}\wedge
 \rd z_{p+1}\wedge\cdots\wedge\rd z_n
\]
with $\rd t/t$.  The boundary contributes the additional density
$\prod_{j=p+1}^{m}|z_j|^{-2c_j}$.  Integrating the $p$ independent
angular variables contributes $(2\pi)^p$, while the radial logarithmic
coordinates contribute the corresponding power of $\log|t|^{-1}$.
With Haar probability on the possibly disconnected kernel of
$(\mathbf R/\mathbf Z)^{p+1}\to\mathbf R/\mathbf Z$,
$(\theta_i)\mapsto\sum b_i\theta_i$, the logarithmic-coordinate Jacobian
relative to integral affine measure is $b_\sigma^{-1}$.  Thus the
arithmetic factor occurs exactly once.

Write $\varepsilon(t)=(\log|t|^{-1})^{-1}$, not to be confused with the
hybrid parameter $\lambda(t)$ of \eqref{eq:hybrid-lambda}, from which it
differs by the factor $-\log r$.  After division by the
normalizing factor in the statement, the remaining radial integral has
mass
\[
 O\!\left(
   |t|^{2(\kappa_0-\kappa_{\min})}
   \varepsilon(t)^{d-q}\right).
\]
Thus the chart contributes zero if its least weight is larger than
$\kappa_{\min}$, or if its minimum face has dimension $q<d$.
When $q=d$, split $b=(b',b'')$ into its first $q+1$ and remaining
coordinates. The constraint-preserving substitution is
\[
\begin{gathered}
 Q_t(w',x'')=
 \bigl((1-\varepsilon(t)b''\!\cdot x'')w',\varepsilon(t)x''\bigr),\\
 b'\!\cdot w'=1,\qquad b''\!\cdot x''\leq\varepsilon(t)^{-1}.
\end{gathered}
\tag{66.34b}\label{eq:hybrid-residue-change-variables}
\]
It keeps $b\cdot Q_t=1$. In these coordinates the variable part of the
Jacobian is $(1-\varepsilon(t)b''\cdot x'')^q\varepsilon(t)^{p-q}$;
the integral affine lattice factors are those of
\cite[Lemma~3.8]{BJ17Hybrid}. This separates normal variables from the
minimum face without leaving the simplex.  Dominated
convergence gives, for every compactly supported test function $\chi$,
\[
 (\operatorname{Log}_{\mathcal X})_*(\chi\mu_t)
 \longrightarrow
 \left(\int_{Y_\sigma}
       \chi\,\operatorname{Res}_{Y_\sigma}(\psi)\right)
 b_\sigma^{-1}\lambda_\sigma .
\]
The inequalities $c_j<1$ are exactly what makes the horizontal normal
integrals and the residual measure finite.  The logarithmic
polar-coordinate estimate
and this dominated-convergence step are the foundational local input
\cite[Lemma~3.5]{BJ17Hybrid}; its subklt extension is carried out in
\cite[Theorem~8.4]{BJ17Hybrid}.

Choose a partition of unity by such charts.  Faces of dimension less than
$d$ have zero $\lambda_\sigma$-measure, so overlaps introduce no
extra mass.  Summing the local limits yields the formula for $\mu_0$.
On each fixed hybrid SNC model, logarithmic pushforward convergence
gives weak convergence by \cite[Theorem~3.4 and Lemma~3.6]{BJ17Hybrid}.
Apply the calculation on all dominating SNC models. Continuity of the
transition maps, as in \cite[Proposition~4.8]{BJ17Hybrid}, and uniqueness of
each weak limit give compatible pushforwards of the residual measures.
The full hybrid space is their compact
inverse limit. Pullbacks of continuous functions from the models form a
separating algebra, hence are uniformly dense by Stone--Weierstrass.
The local mass limits give a uniform total-mass bound; uniform test
approximation therefore passes convergence from that algebra to every
continuous hybrid test. This is \cite[Corollary~4.11]{BJ17Hybrid}, and its
log-pair extension \cite[Theorem~8.4]{BJ17Hybrid}. It supplies the final
passage to the full hybrid space and the asserted support.
\end{proof}

\begin{remark}[Hypotheses are part of the result]
\label{rem:hybrid-volume-hypotheses}
\Cref{thm:hybrid-volume-forms} is not a theorem about every family of smooth volume forms. The logarithmic pluricanonical origin, the log-smooth model, the singularity condition, and the stated normalization are all used to identify the weight-minimizing faces and their residual coefficients.
\end{remark}

\begin{example}[The logarithmic volume on a nodal annulus]
\label{ex:hybrid-nodal-curve}
On the annular part $|t|<|x|<1$ of $xy=t$, take $\Omega=\rd x/x$. With $s=\log|t|^{-1}$, write $x=\rme^{-s\alpha+i\theta}$, where $0<\alpha<1$. The density $|\Omega|^2$ is $s\,\rd\alpha\,\rd\theta$. Thus its mass is $2\pi s$, and for every continuous test $h$ of the retraction coordinate,
\[
 \frac1{2\pi s}\int h(\alpha)|\Omega|^2
       =\int_0^1 h(\alpha)\,\rd\alpha.
\tag{66.34c}\label{eq:hybrid-node-length-calculation}
\]
This is the local residue contribution with $\kappa_{\min}=0$, $d=1$ and $b_\sigma=1$. It has length measure on the whole edge. A relatively nef curve model with component degrees $1$ and $2$ instead has raw measure $\delta_{v_x}+2\delta_{v_y}$ and probability measure $\tfrac13\delta_{v_x}+\tfrac23\delta_{v_y}$. The two calculations have different supports and different raw masses.
\end{example}

There is a related Hodge-theoretic comparison under different hypotheses.
For a torically embedded semistable family whose extended tropicalization
has local tropical Poincar\'e duality and whose open central strata are
wundersch\"on, the tropical cohomology computes the weight-graded limiting
cohomology and nearby Hodge numbers
\cite[Theorem~7.1]{AAPS25Cohomologically}.
\Cref{sec:appG-cohomologically-tropical} states the all-strata and
embedding conditions. Applying it to a particular family requires checking
those conditions in addition to the residue-measure hypotheses above.

\section{Intermediate Calabi--Yau degenerations}
\label{sec:hybrid-intermediate-cy}
\index{Calabi--Yau degeneration!intermediate limit}

The residue measure also guides a more detailed metric description in a
specific class of Calabi--Yau degenerations. Let $M$ be a smooth complex
Fano manifold of dimension $n+1$, let $L$ be ample, and choose positive
integers $d_0,\ldots,d_m$ such that
$-K_M\simeq L^{\otimes(d_0+\cdots+d_m)}$, with $1\leq m\leq n-1$.
Choose sufficiently generic sections
$F_i\in H^0(M,L^{\otimes d_i})$ and $F\in H^0(M,-K_M)$ so that their
divisors are smooth and their intersections are transverse. The family is
\[
 X_t=\{F_0\cdots F_m+tF=0\}\subset M.
\tag{66.34d}\label{eq:hybrid-intermediate-family}
\]
Its smooth nearby fibres are polarized by $L|_{X_t}$, and its essential
skeleton is the simplex
$\Delta_m=\{x_i\geq0:\sum_{i=0}^m x_i=1\}$.
The deepest intersection $Y=\{F_0=\cdots=F_m=0\}$ is a connected
Calabi--Yau manifold of dimension $n-m$.

Li describes the limiting metric by a convex optimal-transport potential
$u$ on $\Delta_m$ \cite[Sections~3--5]{Li25Intermediate}. The target
density in the source's dual coordinates $p$ has weight
\[
 W(p)=\left(1+\sum_{i=0}^m d_i p_i\right)^{n-m}.
\tag{66.34e}\label{eq:hybrid-intermediate-transport}
\]
The normalizing constant makes the source and target transport measures
have the same mass. This weight records the positive-dimensional
Calabi--Yau factor $Y$. The open transport chambers $\Delta_k$ cover
$\Delta_m$ up to a set of Lebesgue measure zero; $u$ is smooth on each
chamber \cite[Corollary~4.6 and Lemma~4.7]{Li25Intermediate}.

Let $T=|\log|t||$. The convex potential determines a continuous
semipositive ansatz in $c_1(L|_{X_t})$, smooth and K\"ahler on the generic
chamber regions \cite[Section~5]{Li25Intermediate}. Denote its curvature
there by $\omega_{\mathrm{ans},t}$ and the Ricci-flat metric in the same
class by $\omega_{\mathrm{CY},t}$. The construction uses an ambient
K\"ahler form whose restriction to $Y$ is the Ricci-flat representative of
$c_1(L|_Y)$; it therefore specifies the ansatz as well as its limiting
potential. After suitable scalar normalization, the Ricci-flat line-bundle
norms raised to $1/T$ converge globally in the source's hybrid convention
\cite[Theorem~1.1 and Proposition~6.13]{Li25Intermediate}. In the book's
convention $|t|_K=r$, the corresponding exponent is
$\lambda(t)=(-\log r)/T$ as in \eqref{eq:hybrid-lambda}; the limiting
logarithmic weights are multiplied by the same fixed factor $-\log r$.

Li--Tosatti's generic regularity result gives metric convergence in this
same family \cite[Theorem~1.1]{LT26Intermediate}. To state the region and
normalization explicitly, put
\[
 V=\int_{X_t}c_1(L)^n,\qquad
 \mu_t=V^{-1}\omega_{\mathrm{CY},t}^{n}.
\]
Thus $\mu_t$ is a probability measure, while the raw Monge--Amp\`ere
mass is $V$. For every fixed $0<\delta<1$, choose compact subsets
$K_k\Subset\Delta_k$ whose complement has sufficiently small Lebesgue
measure and fixed neighborhoods inside the chamber regions. Their
logarithmic lifts, with the local neighborhoods used in the source,
give a generic region $U_{t,\delta}\subset X_t$. For all sufficiently
small $t$, depending on $\delta$, one has
\[
 \mu_t(U_{t,\delta})>1-\delta,\qquad
 \lim_{t\to0}
 \left\|\omega_{\mathrm{CY},t}-\omega_{\mathrm{ans},t}\right\|_
 {C^0(U_{t,\delta},\omega_{\mathrm{ans},t})}=0.
\tag{66.34f}\label{eq:hybrid-intermediate-generic-metric}
\]
Here the tensor norm is measured relative to the ansatz metric. The
estimate therefore compares the two metrics at their collapsing scales,
and it is unchanged if both are multiplied by $T^{-1}$.

Higher regularity is expressed in coordinates adapted to those scales
\cite[Section~2.4, Corollary~5.1 and the proof of Theorem~1.1]{LT26Intermediate}.
After locally unwrapping the torus directions, write
$\zeta_j=-T^{-1}\log z_j$ for $1\leq j\leq m$. The stretched coordinates
are $T^{1/2}\zeta_1,\ldots,T^{1/2}\zeta_m,z_{m+1},\ldots,z_n$,
where the last coordinates lie along $Y$. Derivative estimates use these
coordinates and the associated local reference metrics; they cannot be
read as estimates in fixed coordinates on the entire degenerating fibre.

This example connects three levels of the chapter: the skeletal residue
measure, the globally converging hybrid potential, and local metric
regularity on regions of arbitrarily large probability mass. Its family
and chamber conditions are part of the imported results. The estimates
leave the metric geometry near the omitted chamber walls and simplex
boundary to a separate analysis, which is needed for a global metric-limit
or fibration statement.

\section{Canonically polarized degenerations}
\label{sec:hybrid-ke}

\begin{theorem}\label{thm:hybrid-ke}
Let $X\to\Delta^*$ be a meromorphic degeneration of $n$-dimensional canonically polarized manifolds. Assume that the family admits semistable reduction. Make the required finite base change once, and use $t$ for the new parameter. The canonical bundle metrics induced by the negatively curved K\"ahler--Einstein volume forms define a continuous semipositive hybrid metric on the relative canonical bundle. Their normalized K\"ahler--Einstein measures converge weakly to
\[
              V^{-1}\operatorname{MA}^{\NA}
                     (\phi_{\mathrm{can}}^{\NA}),
\tag{66.35}\label{eq:hybrid-ke-limit}
\]
where $\phi_{\mathrm{can}}^{\NA}$ is the model metric determined by the relative canonical model. The limit is a finite sum of Dirac masses at the divisorial points corresponding to components of that canonical model.
\end{theorem}

\begin{proof}
We separate the birational, analytic, and pluripotential inputs. The
canonical metric is normalized by the K\"ahler--Einstein volume form;
curvature alone would leave an arbitrary base-dependent additive constant
and would not determine the hybrid extension.

\smallskip
\noindent\emph{Step 1: the canonical model.}
After the finite base change implicit in semistable reduction, finite
generation of the relative canonical ring produces the unique relative
canonical model
\[
 \mathscr X_{\mathrm{can}}
   =\operatorname{Proj}_{\Delta}
       \bigoplus_{m\geq0}
       \pi_*\mathcal O_{\mathscr X}
                    (mK_{\mathscr X/\Delta}),
\tag{66.35a}\label{eq:hybrid-relative-canonical-model}
\]
and $K_{\mathscr X_{\mathrm{can}}/\Delta}$ is relatively ample.  The
finite-generation and singularity assertions are deep Minimal Model
Program inputs, recalled in \cite[Theorem~5.3]{PS23Hybrid}; they are not
reproved here.  The relatively ample canonical bundle defines the model
metric $\phi_{\mathrm{can}}^{\NA}$ on the central Berkovich fibre.

\smallskip
\noindent\emph{Step 2: the radial-scale estimate.}
Choose a hybrid Fubini--Study reference metric $\Phi_{\mathrm{FS}}$ from
a relative pluricanonical embedding of the canonical model.  Denote its
ordinary norm weight on a nonzero fibre by $\phi_{\mathrm{FS},t}$, and
write
\[
             \phi_{\mathrm{KE},t}
             =\phi_{\mathrm{FS},t}+\psi_t,
 \qquad
             \varphi_{\mathrm{KE},t}=2\phi_{\mathrm{KE},t}.
\tag{66.35b}\label{eq:hybrid-ke-reference-potential}
\]
Thus $\varphi_{\mathrm{KE},t}$ is the squared-norm potential used by the
complex Monge--Amp\`ere convention in \eqref{eq:hybrid-complex-ma}.
Schumacher's theorem gives psh variation and logarithmic growth.  Fix the strict transform $\widetilde D_{i_0}$ of one component of the
canonical central fibre on the semistable resolution.  Write $\varrho_i$ for
divisorial weights along the other strict transforms and $\eta_j$ for
divisorial weights along the exceptional components meeting
$\widetilde D_{i_0}$.  The additional estimates of Song, in the form
recorded in \cite[Proposition~5.7]{PS23Hybrid}, give a constant $C$ with
$\sup_{X_t}\psi_t\leq C$, and, for every $\varepsilon>0$, a lower
bound
\[
 \psi_t\geq
 \varepsilon\!\left(
    \sum_{i\neq i_0}\varrho_i+
    \sum_{j\in J_{i_0}}\eta_j
 \right)-C_\varepsilon.
\tag{66.35c}\label{eq:hybrid-ke-song-bound}
\]
Choose fixed smooth metrics for the divisor weights. Each is bounded
above on the compact model neighbourhood. The sum of all weights of the
reduced semistable fibre is $\log|t|+O(1)$. Since every omitted weight is
bounded above, the partial sum in \eqref{eq:hybrid-ke-song-bound} is
bounded below by $\log|t|-O(1)$ globally. This uses neither boundedness
below of an omitted weight nor an identification of the partial product
with $t$.

Divide the upper bound and \eqref{eq:hybrid-ke-song-bound} by
$|\log|t||$.  The upper error tends to zero, while the lower error has
liminf at least $-C'\varepsilon$.  Since $\varepsilon$ is arbitrary
we obtain the uniform
radial estimate
\[
       \sup_{X_t}
       \frac{|\psi_t|}{|\log|t||}
             \longrightarrow0.
\tag{66.35d}\label{eq:hybrid-ke-radial-error}
\]
The Song--Sturm--Wang estimates underlying
\eqref{eq:hybrid-ke-song-bound} are the foundational analytic input; the
deduction of \eqref{eq:hybrid-ke-radial-error} is the normalization
argument just given.

By the chart criterion for the hybrid topology,
\eqref{eq:hybrid-ke-radial-error} says precisely that the metric obtained
by placing $\phi_{\mathrm{can}}^{\NA}$ on the central fibre and
$\phi_{\mathrm{KE},t}$ on the nonzero fibres is continuous.
\Cref{thm:hybrid-psh-extension} and the continuity criterion in
\cite[Proposition~3.16]{PS23Hybrid} then make it a continuous semipositive
hybrid metric.  This is the regularity theorem
\cite[Theorem~5.1]{PS23Hybrid}; no measure limit has yet been taken.

\smallskip
\noindent\emph{Step 3: pass to Monge--Amp\`ere measures.}
The hypotheses of \Cref{thm:hybrid-ma-continuity} now hold, so for every
continuous hybrid test function $F$,
\[
 \lim_{t\to0}
 \frac1V\int_{X_t}F\,(\ddc\varphi_{\mathrm{KE},t})^n
 =
 \frac1V\int_{X_K^{\mathrm{an}}}F\,
       \MA^{\NA}(\phi_{\mathrm{can}}^{\NA}).
\tag{66.35e}\label{eq:hybrid-ke-limit-test}
\]
This proves \eqref{eq:hybrid-ke-limit}; the limit exchange is exactly the
one justified in the proof of \Cref{thm:hybrid-ma-continuity} and is not
repeated with a fibre-dependent approximation.

Finally write
$\mathscr X_{\mathrm{can},0}=\sum_E b_EE$.  Since the central metric is a
model metric, \eqref{eq:hybrid-model-ma} gives
\[
 \frac1V\MA^{\NA}(\phi_{\mathrm{can}}^{\NA})
  =\frac1V\sum_E
       b_E\bigl(c_1(K_{\mathscr X_{\mathrm{can}}/\Delta}|_E)^n\bigr)
       \delta_{v_E}.
\tag{66.35f}\label{eq:hybrid-ke-atomic-formula}
\]
Thus the limiting measure is the asserted finite atomic probability
measure.  This last step is intersection theory; the hard PDE regularity
has already been isolated in Step~2.
\end{proof}

\begin{example}[Weights for stable curves]
\label{ex:hybrid-ke-curves}
Let the stable central fibre have components $C_i$, and let $g_i$ be the genus of the normalization of $C_i$. Adjunction gives the raw mass
\[
 2g_i-2+\operatorname{val}(C_i),
\tag{66.36}\label{eq:hybrid-stable-curve-weight}
\]
where valence counts branches of nodes: a self-node counts twice. Stability makes each coefficient positive. If the connected dual graph has $v$ vertices and $e$ edges, its arithmetic genus is $g=\sum_i g_i+e-v+1$. Summing the displayed masses gives $2\sum_i g_i-2v+2e=2g-2$, as required. For two elliptic components meeting at one node, $g=2$ and both masses are $1$; the probability weights are $1/2,1/2$. For an irreducible rational curve with two self-nodes, $g_1=0$ and the valence is $4$, so $g=2$ and the single raw mass is $2$. Its probability measure is one Dirac mass. Using arithmetic genus $2$ in place of normalization genus $0$ would incorrectly give mass $6$ in the latter example.
\end{example}

\begin{remark}[Two kinds of limiting measure]
\label{rem:hybrid-two-limits}
The skeletal measure of \Cref{thm:hybrid-volume-forms} and the atomic measure of \Cref{thm:hybrid-ke} answer different normalization problems. One should not identify them merely because both live on $X_K^{\mathrm{an}}$.
\end{remark}

\section{Energy and entropy along the hybrid parameter}
\label{sec:hybrid-energy-entropy}

Let $\phi_{\mathrm{ref},t}$ be a reference metric induced by a
projective model, and let $\varphi_{\mathrm{ref},t}$ be its complex
squared-norm potential. The relative complex Monge--Amp\`ere energy is
\[
 E_t(\varphi_t)
 =\frac{1}{(n+1)V}
   \sum_{j=0}^{n}
   \int_{X_t}
     (\varphi_t-\varphi_{\mathrm{ref},t})
     (\ddc\varphi_t)^j
     (\ddc\varphi_{\mathrm{ref},t})^{n-j}.
\tag{66.37}\label{eq:hybrid-energy}
\]
The non-Archimedean energy below is relative to
$\phi_{\mathrm{ref}}^{\NA}$ and normalized by $V$, with reference energy
zero, as in \Cref{sec:cma-finite-energy}.

\begin{proposition}\label{prop:hybrid-energy-continuity}
Under the hypotheses of \Cref{thm:hybrid-ma-continuity}, suppose that the reference metric is also continuous and semipositive on the hybrid line bundle. Then
\[
 \lim_{t\to0}
 \frac{E_t(\varphi_t)}{2\log_r|t|}
 =E^{\NA}(\phi^{\NA}).
\tag{66.38}\label{eq:hybrid-energy-limit}
\]
The same conclusion holds for mixed energies obtained by polarization.
\end{proposition}

\begin{proof}
Put $u=\Phi-\Phi_{\mathrm{ref}}$, a continuous function on the compact
hybrid space. On a complex fibre its value is
$u=\lambda(t)(\varphi_t-\varphi_{\mathrm{ref},t})/2$.
Polarization of \Cref{thm:hybrid-ma-continuity} gives convergence of all
mixed probability measures $V^{-1}(\ddc\varphi_t)^j
(\ddc\varphi_{\mathrm{ref},t})^{n-j}$ to their central mixed measures.
Since $1/\log_r|t|=\lambda(t)$, each summand of the scaled energy is the
integral of the fixed continuous test $u$ against such a measure.
Their $n+1$ limiting integrals, averaged, are precisely the relative
non-Archimedean energy. This proves \eqref{eq:hybrid-energy-limit} and
its mixed version, including dimension zero.

The approximation estimate is also explicit. For squared-norm potentials
with the same reference, monotonicity and the cocycle formula give
\[
 |E_t(\varphi_t)-E_t(\vartheta_t)|
       \leq\|\varphi_t-\vartheta_t\|_\infty.
\tag{66.39}\label{eq:hybrid-energy-lipschitz}
\]
A uniform hybrid error $\varepsilon$ therefore gives a complex error
$2\log_r|t|\varepsilon$, hence a scaled energy error $\varepsilon$.
Continuous hybrid semipositivity already supplies uniform Fubini--Study
approximation; no separate all-degree limiting assertion is needed.
\end{proof}

\begin{proposition}\label{prop:hybrid-energy-slope}
Let $(\mathcal X,\mathcal L)$ be a normal relatively ample test configuration for a connected polarized projective manifold, and let $(\varphi_s)_{s\geq0}$ be its filtration weak geodesic ray from a smooth positive reference in the book's squared-norm convention, with $s=-\log|t|$. Write $\widehat\phi_{\mathcal L}^{\NA}$ for its unit-normalized test-configuration metric on $X^{\mathrm{an,triv}}$. It is obtained from the $K$-metric by Gauss restriction and division of weights by $-\log r$, as in \Cref{chap:master-dictionary}. Then
\[
       \lim_{s\to+\infty}\frac{E(\varphi_s)}{2s}
             =E^{\NA}(\widehat\phi_{\mathcal L}^{\NA}),
       \qquad
       \lim_{s\to+\infty}\frac{J(\varphi_s)}{2s}
             =J^{\NA}(\widehat\phi_{\mathcal L}^{\NA}).
\tag{66.40}\label{eq:hybrid-ej-slopes}
\]
\end{proposition}

\begin{proof}
Let $\mathcal F$ be the associated finitely generated filtration,
$\Gamma_\tau$ its closed maximal test curve, and $c$ its cutoff.
By \Cref{thm:rf-phong-sturm}, the relative squared-norm potential is
$\ell_{2s}^{\mathcal F}$. By \Cref{thm:rf-energy-slope}, its normalized
energy slope is the first moment of the probability DH measure, which is
$E^{\NA}(\widehat\phi_{\mathcal L}^{\NA})$.

For the linear term in $J$, use the actual Legendre formula
$\ell_{2s}(x)=\sup_{\tau\leq c}(\Gamma_\tau(x)+2s\tau)$.
For each $\tau<c$, $\Gamma_\tau$ is not identically $-\infty$.
Choose a countable sequence of these levels increasing to $c$.
Outside their pluripolar exceptional sets, division by $2s$ gives
$\liminf\ell_{2s}(x)/(2s)\geq\tau$ for every chosen level.
The bounds $C_-\leq\ell_{2s}/(2s)\leq c$ from
\Cref{thm:rf-filtration-test-curve} imply convergence to $c$ almost
everywhere for the smooth reference probability measure $V^{-1}\theta^n$.
Dominated convergence yields
$\lim_{s\to\infty}V^{-1}\int\ell_{2s}\theta^n/(2s)=c$.
By definition $J(\ell)=V^{-1}\int\ell\theta^n-E(\ell)$, so its slope is
$c-E^{\NA}$. The cutoff equals the metric value at the trivial valuation
by \Cref{cor:rf-normalization-same-ray}; thus this difference is
$J^{\NA}(\widehat\phi_{\mathcal L}^{\NA})$. In dimension zero both $J$
values vanish directly. This argument proves the weak-ray statement from
its filtration identity; the smooth-ray intersection statement in
\Cref{thm:rf-bhj-unaffected-slopes} is not used as a regularity theorem
for the weak ray.
\end{proof}

\subsection{Compatibility with the ray parameter}
\label{sec:hybrid-ray-parameter}

For a test configuration, the punctured-disc coordinate gives
$s=-\log|t|$.  If $\psi_s$ is the ordinary norm weight used in the
source slope formula and $\varphi_s=2\psi_s$ is the book's squared-norm
potential, then
\[
 \frac{E(\varphi_s)}{2s}=\frac{E_{\mathrm{wt}}(\psi_s)}s.
\tag{66.41}\label{eq:hybrid-ray-factor-two-test}
\]
Here $E_{\mathrm{wt}}$ uses ordinary norm weights and the source operator
$i\partial\bar\partial/\pi$, with corresponding reference curvature and
probability normalization. The energies obey $E(2\psi)=2E_{\mathrm{wt}}(\psi)$
because their curvature measures agree. This changes metric conventions
and reference weights together.

Adding weight $c$ to the unit-normalized test-configuration metric adds
$2cs$ to the squared-norm ray; the raw central hybrid weight changes by
$(-\log r)c$. Since normalized energy translates by constants,
\[
 \frac{E(\varphi_s+2cs)-E(\varphi_s)}{2s}=c,
\tag{66.42}\label{eq:hybrid-energy-translation-test}
\]
matching $E^{\NA}(\widehat\phi^{\NA}+c)=E^{\NA}(\widehat\phi^{\NA})+c$.
Equation \eqref{eq:hybrid-energy-translation-test} is the fastest check that
the factor two and the additive normalization have both been transported.

Entropy is more delicate. Fix a Radon probability measure
$\nu_t$ on $X_t$ and put
\[
  \operatorname{Ent}_{\nu_t}(\mu_t)
      :=
      \begin{cases}
       \displaystyle\int_{X_t}
          \log\!\left(\frac{\rd\mu_t}{\rd\nu_t}\right)\,\rd\mu_t,
          &\mu_t\ll\nu_t,\\[6pt]
       +\infty,&\text{otherwise}.
      \end{cases}
\tag{66.43}\label{eq:hybrid-entropy}
\]
Weak convergence of $\mu_t$ alone does not imply convergence of
\eqref{eq:hybrid-entropy}.

\begin{proposition}\label{prop:hybrid-entropy-liminf}
Suppose that $\mu_t$ and $\nu_t$ are Radon probability measures on the compact hybrid space, supported on their respective fibres, and that $\mu_t\to\mu_0$ and $\nu_t\to\nu_0$ weakly. Then, with extended values,
\[
 \operatorname{Ent}_{\nu_0}(\mu_0)
       \leq\liminf_{t\to0}\operatorname{Ent}_{\nu_t}(\mu_t).
\tag{66.44}\label{eq:hybrid-entropy-liminf}
\]
\end{proposition}

\begin{proof}
For Radon probability measures on a compact Hausdorff space $H$, the
relative-entropy variational formula is
\[
 \operatorname{Ent}_{\nu}(\mu)
 =\sup_{f\in C(H)}
   \left\{\int f\,\rd\mu-\log\int\rme^f\,\rd\nu\right\}.
\tag{66.45}\label{eq:hybrid-entropy-dual}
\]
Its upper bound follows from Jensen's inequality applied to the tilted
probability measure $\rme^f\nu/\int\rme^f\rd\nu$. The reverse bound uses
bounded truncations of $\log(\rd\mu/\rd\nu)$, followed by regularity of
$\mu+\nu$ and continuous approximation. For fixed truncation, approximate
in $L^1(\mu+\nu)$ by uniformly bounded continuous functions; the exponential
is Lipschitz on that bounded interval. Then let the truncation grow.
If $\mu$ has a singular part, compact subsets of a $\nu$-null set and
Urysohn functions with shrinking $\nu$-mass give arbitrarily large values
of the supremum. This proves the formula also for infinite entropy.
Use $H=\mathcal X^{\mathrm{hyb}}$. For each fixed $f$, both integrals
converge weakly, and the exponential integral is positive. Take the lower
limit and then the supremum to obtain the assertion. Entropy computed on
a fibre equals ambient entropy since both measures are supported there.
No density bound is needed for this lower semicontinuity statement.
\end{proof}

\begin{example}[Total variation convergence with divergent entropy]
\label{ex:hybrid-entropy-concentration}
On $[0,1]$, let $\nu$ be Lebesgue probability measure and, for $j\geq2$, put $a_j=1/j$, $\delta_j=\rme^{-j^2}$ and
\[
 \frac{\rd\mu_j}{\rd\nu}
   =1-a_j+\frac{a_j}{\delta_j}\mathbf1_{[0,\delta_j]}.
\tag{66.45a}\label{eq:hybrid-entropy-spike}
\]
The density has integral one. Its total variation norm difference from $\nu$ is $2a_j(1-\delta_j)\to0$. On the short interval the density is at least $a_j/\delta_j>1$ and its mass is at least $a_j$. On the complement its contribution to entropy is at least $(1-a_j)\log(1-a_j)$. Consequently
\[
\begin{aligned}
 \operatorname{Ent}_\nu(\mu_j)
 &\geq a_j\log(a_j/\delta_j)+(1-a_j)\log(1-a_j)\\
 &=j-\frac{\log j}{j}+(1-1/j)\log(1-1/j)\longrightarrow+\infty.
\end{aligned}
\tag{66.45b}\label{eq:hybrid-entropy-spike-bound}
\]
The limiting entropy is $\operatorname{Ent}_\nu(\nu)=0$. This elementary measure example shows why a weak measure theorem, even strengthened to total variation convergence, cannot supply entropy equality.
\end{example}

\begin{remark}[No entropy equality from weak measure convergence]
\label{fire:hybrid-entropy}
Neither \Cref{thm:hybrid-ma-continuity} nor \Cref{prop:hybrid-entropy-liminf} proves entropy convergence. Equality requires recovery sequences or uniform integrability of the logarithmic densities. We will not use the disputed entropy-regularization step attached to \cite[Theorem~5.6]{BHJ19}; its published erratum is discussed in \Cref{chap:ytd-bridge}.
\end{remark}

\section{A convergence checklist}
\label{sec:hybrid-checklist}

\begin{proposition}\label{prop:hybrid-four-layers}
For a degenerating family of metrics, the following implications are valid only with the indicated additional inputs:
\begin{enumerate}
\item regularity upgrades a hybrid psh extension to a continuous
hybrid metric;
\item the projective Monge--Amp\`ere theorem then gives weak measure
convergence;
\item continuous hybrid metrics give energy convergence by mixed measure
convergence; and
\item entropy convergence needs its own density or recovery input.
\end{enumerate}
There is no downward implication from weak Monge--Amp\`ere convergence to entropy convergence.
\end{proposition}

\begin{proof}
The extension is \Cref{thm:hybrid-psh-extension}; its continuity needs a
separate input, such as \Cref{rem:hybrid-extension-not-continuity} or
\Cref{thm:hybrid-ke}. The measure step is
\Cref{thm:hybrid-ma-continuity}, and the energy step is
\Cref{prop:hybrid-energy-continuity}. Entropy convergence
is an additional scalar assertion about a family whose weak convergence
has already been obtained in the second layer; convergence of entropy
values alone would not determine a weak limit. The failure of the implication from weak convergence to entropy
convergence is the explicit calculation in
\Cref{ex:hybrid-entropy-concentration}.
\end{proof}

\begin{example}[The constant-weight product family]
\label{ex:hybrid-safe-procedure}
Take $\mathbf P^1\times\Delta$ polarized by $\mathcal O(1)$, its smooth Fubini--Study norm weight $\psi_{\mathrm{FS}}$, and $a\in\mathbf R$. The metric $\psi_t=\psi_{\mathrm{FS}}+a\log_r|t|$ is psh on the punctured family since the added term is harmonic there. Its hybrid extension is the continuous Fubini--Study hybrid metric plus the constant $a$; hence its central metric is $\phi_{\mathrm{FS}}^{\NA}+a$. Here $V=1$, and the complex probability measures are the unchanged Fubini--Study measures. \Cref{thm:hybrid-ma-continuity} gives their hybrid limit $\delta_{v_G}$ at the product model's Gauss point. Relative to the unshifted reference, $E_t=2a\log_r|t|$, so the scaled energy and central energy both equal $a$. Taking $\nu_t=\mu_t$, including $\nu_0=\mu_0=\delta_{v_G}$, gives entropy zero on every fibre. Thus growth, continuity, semipositivity, the measure limit, the energy value and the separate entropy computation can all be checked in one family. Changing the metric by a constant does not alter its measure.
\end{example}

\section*{Source notes}
\label{sec:hybrid-source-notes}

The hybrid topology and the skeletal limits of logarithmic volume
forms are taken from \cite{BJ17Hybrid}. The K\"ahler--Einstein
measure theorem is \cite[Theorem~A]{PS22HybridKE}; the continuous hybrid
psh formalism, projective Monge--Amp\`ere continuity, and continuity of
the canonically polarized hybrid metric are respectively
\cite[Theorems~3.13, 4.11, and~5.1]{PS23Hybrid}; the last measure
convergence is also \cite[Corollary~5.2]{PS23Hybrid}.
The weak-ray energy slope and the accompanying $J$ slope use the
filtration identities of \Cref{chap:rays-filtrations} and the explicit
linear-term argument above; their intersection values lie in the
unaffected range of \cite{BHJ17,BHJ19}. The general
Banach-family Monge--Amp\`ere problem was posed in
\cite[Question~4.9]{PS23Hybrid}; it is answered by
\cref{thm:p171-full-assembly}.  Entropy equality remains outside the
theorem statements.  The literal
Favre--Pille-Schneider hybrid-continuity question is disproved by
\Cref{thm:hybrid-endpoint-continuity-counterexample}, even after imposing
zero-supremum normalization relative to a fixed continuous model metric on
every complex fibre and on the non-Archimedean fibre.
Energy continuity uses the polarized hybrid measure theorem; weak-ray
slopes use the filtration ray, probability DH measure and normalized
intersection value. No disputed entropy-regularization step is used.

The intermediate Calabi--Yau discussion uses the published edition of
\cite{Li25Intermediate} and version~2 of \cite{LT26Intermediate}.
The former supplies the transport potential and global hybrid norm limit;
the latter supplies generic metric regularity for
\eqref{eq:hybrid-intermediate-family}, with the region and probability
normalization in \eqref{eq:hybrid-intermediate-generic-metric}.

The preceding geometric and variational source imports have five distinct boundaries.
\cite{BJ17Hybrid} supplies the hybrid topology and the residue-theoretic
volume-form limit; \cite{PS23Hybrid} supplies the canonical psh extension,
projective Monge--Amp\`ere continuity, and the canonically polarized
regularity statement at the theorem numbers listed above;
\cite{PS22HybridKE} supplies the K\"ahler--Einstein measure limit;
\cite{Li25Intermediate,LT26Intermediate} supply the specialized
intermediate Calabi--Yau description; and the energy/J slope argument uses
the filtration identities and intersection formulas stated above. The smooth-ray range of
\cite{BHJ19} remains subject to the boundary in \cite{BHJ19err}.  None of these imports by
itself supplies entropy equality or the integral Banach-family theorem;
the latter is the verified synthesis in
\cref{sec:hybrid-banach-family-ma}.  Continuity still requires a separate regularity input: endpoint
continuity alone is disproved by
\Cref{thm:hybrid-endpoint-continuity-counterexample}.

The dynamical additions retain their own scopes. The polynomial route
uses explicit uniformity, rank and closed-support inputs. The sequential
relative line of \cite{FavreGong26SequentialV4} has a product-ring base
and a separate singular-integral obligation; the good-base theorem of
\cite{Poi25DynamicsZV5} has its stated local topological hypotheses.
The K3 normalization follows \cite{Cantat01K3Published,Filip19K3V1},
while the family limit retains uniformity and real nef mixed-product
inputs. The exact-source boundaries are recorded in
\Cref{sec:provenance-hybrid-sequential,sec:provenance-hybrid-K3}.
